% De Finibus Bonorum et Malorum (de-finibus) --- English translation
% Translated by Claude (Anthropic) from the Latin (Perseus).
% Style guide: STYLE.md.

\ciceroBook{1}

\ciceroSection{1} I was not unaware, Brutus, that when I committed to Latin letters what the philosophers had treated in Greek, working through those subjects with the greatest minds and with painstaking learning, this labor of mine would run into criticism of various kinds. For there are those — and among them some who are not altogether unlearned — to whom the whole enterprise of philosophizing in this way is distasteful. Others do not censure it so much if it is pursued with some restraint, but think that so much zeal and so great an expenditure of effort ought not to be given to it. There will also be those, learned in Greek letters, who will look down on Latin and say they would rather spend their time reading Greek. Finally, I suspect there will be some who will call me to other kinds of writing and deny that this genre, though it may be elegant enough, is suited to my person and standing.

\ciceroSection{2} Against all these I think I should speak briefly. Philosophy's detractors have already received a sufficient answer in the book in which I defended and praised philosophy when it was accused and attacked by Hortensius. That book, since it seemed approved by you and by those whom I judged capable of judging, I went on to add more works — afraid that I might appear to stir men's enthusiasm while being unable to sustain it. As for those who, granting that this pursuit is welcome, yet want it conducted with more restraint, they demand a kind of impossible moderation in a thing that, once admitted, cannot be checked or suppressed; so that we would do nearly better to take the side of those who would draw us away from philosophy altogether than of those who set limits to an unlimited subject and demand mediocrity in a pursuit that, the greater it is, the better.

\ciceroSection{3} For if wisdom can be reached, we must not only acquire it but enjoy it as well; if it is difficult to reach, there is still no means of tracking down the truth except by finding it, and to grow weary of seeking is shameful when what is sought is supremely beautiful. And besides: if we take delight in writing, who is so jealous as to pull us away from it? If we labor at it, who is to set a limit on another man's industry? There is, to be sure, the Chremes of Terence — not an unkind neighbor, yet one who does not want the newcomer next door "to dig or plow, or bring \textit{anything} in the end" — for he is not deterring the man from industry but from degrading toil. In the same way these busybodies are put out by our labor, which is anything but disagreeable to us.

\ciceroSection{4} It is harder, then, to satisfy those who profess contempt for Latin writing. And here the first thing that surprises me is why serious subjects should give them no pleasure in their native tongue, when the very same people read Latin comedies — translated word for word from the Greek — without reluctance. Who is so hostile to the very name of Rome as to spurn or reject Ennius's \textit{Medea} or Pacuvius's \textit{Antiope}, on the grounds that he takes delight in those same plays of Euripides but hates Latin letters? "Am I," he says, "to read Caecilius's \textit{Synepheboi} or Terence's \textit{Andria} rather than both originals by Menander?" I disagree so strongly with such a view that, even though Sophocles wrote the \textit{Electra} supremely well, I think I must read Atilius's poor translation — of whom Lucilius says:

\ciceroSection{5} "an iron writer" — truly a writer, I suppose, nonetheless, to the extent that he must be read. To be entirely ignorant of our own poets is the mark of the laziest indolence or the most fastidious disdain. In my own judgment, no one is adequately educated who is ignorant of what is ours. Do we not read "O that in the forest…" — that passage — no less than the same line in Greek? And shall the discussions of Plato on living well and happily not be welcome when set out in Latin?

\ciceroSection{6} What of the fact that we are not performing the work of translators but are expounding the views of those we approve, while adding our own judgment and the arrangement of our own writing — what reason have they to prefer Greek to works that are both splendidly expressed and not merely translated from Greek? If they will say that these subjects have already been treated by the Greeks, then there is no reason even to read as many Greek authors as should be read. For what has Chrysippus left untreated among the Stoics? Yet we read Diogenes, Antipater, Mnesarchus, Panaetius, and many others — above all our friend Posidonius. Again: does not Theophrastus give us moderate delight when he works over ground already worked by Aristotle? Do the Epicureans cease to write on the same subjects on which Epicurus himself and the ancients wrote, in whatever way seems good to them? And if Greeks are read by Greeks on the same subjects arranged in different form, why should our own writers not be read by our own?

\ciceroSection{7} Though indeed, if I were translating Plato or Aristotle as our poets translated the tragedies — word for word — I should, I suppose, be doing a poor service to my countrymen if I were to bring those divine intelligences within their reach. But that I have not done as yet, nor do I think I am forbidden to do it. I shall translate certain passages when occasion arises and it can be done aptly — above all from those I have just named — just as Ennius does from Homer, or Afranius from Menander. Nor, unlike our Lucilius, shall I refuse to let everyone read what I write. He wished Persius would read him — and far more wished it of Scipio and Rutilius — and it was out of dread of their judgment that he said he wrote for the people of Tarentum and Consentia and Sicily. A witty remark, like his others; but the readers in his day were not so learned as to require him to labor for their verdict, and his writings are on the lighter side — the wit is supreme, the learning only moderate.

\ciceroSection{8} As for me, what reader should I fear, when I dare write for you — a man who yields to no Greek in philosophy? Though I am, to be sure, moved to write by you yourself, by that most welcome book you sent me on virtue. But I believe some people come to be repelled by Latin writing because they have happened upon certain rough and unpolished pieces — bad Greek material turned into Latin worse still. With such people I agree, provided they think the same Greek authors not worth reading either. Genuinely good matter expressed in well-chosen words, set forth with weight and distinction — who would not read it? Unless it were someone who wished to be taken as a thoroughgoing Greek, like the Albucius whom Scaevola, when praetor at Athens, greeted accordingly.

\ciceroSection{9} That scene Lucilius handles with great grace and full seasoning of wit, giving Scaevola these fine words: "You chose, Albucius, to be called a Greek rather than a Roman and a Sabine — a fellow citizen of Pontius, of Tritanus, of the centurions, men of the first rank and standard-bearers. And so, as praetor at Athens, I greet you in Greek, as you preferred, when you come to call on me: 'Greetings,' I say, 'Titus!' My lictors, my whole retinue, my staff — 'Greetings, Titus!' Hence Albucius is my enemy, hence my foe." Mucius, in short, had justice on his side.

\ciceroSection{10} As for myself, I cannot wonder enough at the source of this so unusual contempt for what is one's own. This is not the place for a lecture on the matter; but my settled view — one I have often argued — is that the Latin language is not merely not impoverished, as people commonly suppose, but is actually richer than Greek. When have we ever lacked, or our good orators and poets lacked — especially once there was a model to imitate — any ornament of speech, whether abundant or elegant? I, for my part, since I cannot feel that I abandoned my post in the work of the Forum, its labors and its dangers — the post at which the Roman people stationed me — surely owe it, to the full extent of my power, to labor in this field as well, so that my fellow citizens may be better educated through my work, my zeal, my toil. Nor need I fight so hard against those who prefer to read Greek, provided they really do read it and do not merely pretend to; and I am glad to serve those who either wish to use both literatures or, if they have ours, do not miss the other too greatly.

\ciceroSection{11} Those who prefer that I write other things should be fair: much has already been written — more than any other of our countrymen has produced — and more will perhaps be written if life holds out. And yet anyone who makes a habit of reading carefully what I am committing to writing on philosophy will judge that there is nothing better worth reading than this. For what in life is worth seeking so earnestly as, among all the things that philosophy offers, the very thing these books seek — what is the end, what the final point, what the ultimate goal, to which all the counsels of living well and acting rightly are to be referred; what nature points to as the highest of things to be desired, what she bids us flee as the extreme of evils? Since on this question there is the greatest disagreement among the most learned, who would think it unworthy of whatever standing people may allow me to investigate what is best and truest in the entire conduct of life?

\ciceroSection{12} What — the question whether the offspring of a slave woman is to count as \textit{fructus}, as produce — that is debated between the leading men of the state, Publius Scaevola and Manius Manilius, and Marcus Brutus dissents from them — an acute field of inquiry and not useless to the practical needs of citizens; and we both read that kind of writing and shall continue to read it gladly — while these matters, which govern life entire, are to be neglected? Granted, that kind of writing may be more saleable; but this kind is certainly more abundant. Though indeed those who have read can form their own judgment on that. For my part, I hold that this whole inquiry into the ends of good and evil has been in these pages very largely worked through by us — pursuing, to the best of our ability, not only what seemed true to us but what each school of philosophy had to say.

\ciceroSection{13} To begin with what is easiest: let Epicurus's theory come forward first, since it is most widely known. You will find it set out by me as carefully as it is usually expounded by those who belong to that school themselves — for it is truth we wish to find, not an adversary to demolish. The Epicurean position on pleasure was once argued with great precision by Lucius Torquatus, a man educated in every branch of learning, and I gave answer to him, while the young Gaius Triarius — a person of the greatest gravity and learning — was present at the discussion.

\ciceroSection{14} For when both came to visit me at my place near Cumae, we spoke at first for a little while between us about literature — in which both had an intense enthusiasm — and then Torquatus said: "Now that we have found you at leisure for once, I shall certainly hear what it is that you do not hate our Epicurus — as those who disagree with him generally do — but certainly do not approve of him either: a man whom I hold to be the one who saw the truth, who freed human minds from the greatest errors, and who handed on all that pertains to living well and happily. Though I suppose, as with our friend Triarius here, you are less pleased with him because he neglected those ornaments of style in Plato, Aristotle, and Theophrastus. For I can scarcely be brought to believe that his actual doctrines do not seem true to you."

\ciceroSection{15} "See how far you are mistaken, Torquatus," I said. "That philosopher's style does not put me off — he grasps in words what he wishes to say, he speaks plainly enough for me to understand him; and yet I would not spurn eloquence in a philosopher if he brings it, nor insist on it if he does not. It is on the substance that he does not satisfy me equally — and on quite a number of points. But there are as many opinions as there are men; so we may be mistaken." "Why then, at long last," he said, "does he not satisfy you? I think you a fair judge, provided you are really familiar with what he says."

\ciceroSection{16} "Unless you think Phaedrus lied to me," I said, "or Zeno — both of whom I heard, though the only thing about either that really earned my approval was their diligence — all of Epicurus's views are sufficiently known to me. I heard both those men I named frequently, in company with our friend Atticus, who admired them both, and had a genuine affection for Phaedrus besides; and every day we compared notes with each other on what we had heard. There was never any dispute over what I understood — only over what I approved."

\ciceroSection{17} "What is it, then?" he said — "for I am eager to hear what you do not approve." "To begin with natural philosophy, in which Epicurus takes the greatest pride — first, he borrows wholesale from Democritus, changing very little, but changing it, I think at least, in ways that make his borrowings worse rather than better. Democritus holds that what he calls atoms — that is, indivisible bodies on account of their solidity — are borne through infinite void, in which there is nothing highest, nothing lowest, nothing middle, nothing last, nothing outermost, in such a way that they cohere together through collisions, and from this are produced all existing and visible things; and that this motion of the atoms is to be understood as having no beginning, but as eternal."

\ciceroSection{18} "Now where Epicurus follows Democritus, he does not, on the whole, go wrong. Though there is much I do not approve in either of them, and above all this: that when in the study of nature two things are called for — first, what the matter is out of which each thing is made, and second, what the force is that makes each thing — they have treated the matter but left aside the force and cause of making. That, however, is a fault they share. What follows is the ruin peculiar to Epicurus: he holds that those same indivisible solid bodies travel straight downward by their own weight in a perpendicular line, and that this is the natural motion of all bodies."


\ciceroSection{19} And then that same clever man, when it struck him that if all atoms fell straight downward in a line, as I said, one could never touch another, and so … brought in a fabricated device: he declared that the atom swerves ever so slightly — by the smallest amount that can possibly be conceived — and that from this swerve arise the interlockings and joinings and adhesions of atoms among themselves, from which in turn the world is produced, and all the parts of the world, and all the things within it. Now this whole scheme is childish in its invention, and it does not even achieve what he wants it to achieve. For the very swerve is invented arbitrarily — he says the atom swerves without cause; and nothing is more disgraceful in a natural philosopher than to say that something happens without cause. He also stripped the atoms, without cause, of that natural motion by which, on his own account, all heavy bodies seek a lower place in a straight line; and yet he did not attain the very effect for the sake of which he invented this fiction.

\ciceroSection{20} For if all atoms swerve, none will ever cohere; or if some swerve and others are carried straight by their own impulse, this will in effect be assigning different provinces to different atoms — some to travel straight, some obliquely — and then that turbulent collision of atoms, in which even Democritus becomes entangled, will still prove unable to produce this ordered beauty of the world. Nor is it worthy of a natural philosopher to believe in something that is an absolute minimum — a belief he would never have formed had he preferred to learn geometry from his friend Polyaenus rather than to un-teach it to him. Democritus, a learned man, accomplished in geometry, thought the sun very large; this man thinks it perhaps a foot across — for he judges it to be just as large as it looks, or a little larger or smaller.

\ciceroSection{21} So whatever he changes, he corrupts; what he retains is entirely Democritus's — the atoms, the void, the images which they call \textit{[Greek: eidola]}, whose incursion accounts not only for our seeing but for our thinking as well; and the infinite itself, which they call \textit{[Greek: apeiria]}, is taken entirely from Democritus, as are the countless worlds that arise and perish daily. These views I find wholly unconvincing; but Democritus, praised by others, I would not have wished to see vilified by the one man who had chosen to follow him alone.

\ciceroSection{22} Now in the other part of philosophy — the part concerned with inquiry and argumentation, which is called logic \textit{[Greek: logike]} — your master, it seems to me, is entirely unarmed and stripped bare. He does away with definitions, teaches nothing about division and partition, shows no method by which an argument is constructed and brought to a conclusion, indicates no way by which sophistical snares are to be unraveled or ambiguities distinguished. He places the judgment of things in the senses; and if once anything false is admitted as true by the senses, he believes that every means of judging truth from falsehood has been swept away.

\ciceroSection{23} His central affirmation, however, is that nature herself — as he puts it — ordains and sanctions pleasure and pain. To these he refers everything we pursue and everything we shun. Now although this position is maintained more fully and more freely by Aristippus and the Cyrenaics, still I hold it to be of such a kind that nothing seems more unworthy of a human being. For nature has generated and shaped us for something rather greater, or so at least it seems to me. I may perhaps be wrong; but this is precisely my conviction: that neither the Torquatus who first won that cognomen — tearing the torc from an enemy's neck — acted as he did in order to draw from it some bodily pleasure, nor when he engaged the Latins at the Veseris in his third consulship was he driven by pleasure; and as for sentencing his own son to death by the axe, that act visibly deprived him of many pleasures — for he placed the authority of martial command and of law above nature itself and above a father's love.

\ciceroSection{24} What of T. Torquatus, consul with Cn. Octavius, who exercised that severity toward the son he had emancipated into the adoption of D. Silanus — when the Macedonian envoys brought charges against him, alleging that as praetor in the province he had taken money under color of office, Torquatus ordered him to plead his case before him, heard both sides fully, and then pronounced that the man had not shown himself in the exercise of command as his ancestors had been, and forbade him to come into his presence? Do you really think that man had his own pleasures in mind? But let us leave aside dangers, labors, and even pain, which every excellent man takes upon himself for country and kindred — not merely forgoing all pleasures in the process, but even deliberately passing them by, and preferring to endure any sufferings whatever rather than abandon any part of his duty. Let us come instead to things which make the point no less clearly but seem of lighter weight. What pleasure, Torquatus, does literary study bring to you — what to Triarius here? What pleasure do historical inquiry and the knowledge of affairs give you, what the reading through of the poets, what the retention in memory of so many thousands of lines?

\ciceroSection{25} Do not tell me: 'Why, these things are themselves a pleasure to me, and they were to the Torquati.' Epicurus never defended that position, nor did Metrodorus or anyone among them who had any understanding or had studied the matter. And as for the question often asked — why are there so many Epicureans? — there are other reasons too, but above all the crowd is attracted by the belief that Epicurus said this: that whatever is morally right and honorable is itself the source of joy, that is, of pleasure. People of the best sort fail to see that if this were so, the whole system would be overturned; for if it were granted that these pursuits are pleasant in and of themselves, apart from any reference to the body, then both virtue and the knowledge of things would have to be sought for their own sake — which is precisely what Epicurus is least willing to allow.

\ciceroSection{26} These positions of Epicurus, then, I do not accept, I said. As for the rest, I could wish either that he himself had been better furnished with the arts of learning — for it must be evident to you that he was insufficiently schooled in those studies whose masters are called educated men — or that he had not discouraged others from their pursuit. Though I see that you yourself have been discouraged not at all. When I had said this — meaning more to provoke him than to make a speech myself — Triarius smiled gently and said: You have all but driven Epicurus from the whole chorus of philosophers. What have you left him, except that, however he expresses himself, you understand what he means? He borrowed his physics from others and did not even borrow the parts of them you find convincing; where he tried to correct them, he made them worse. He had no skill in argument. When he called pleasure the highest good, he first saw too little in that very claim, and then it was not even his — for Aristippus had said it before him, and said it better. And you added, in the end, that he was also uneducated.

\ciceroSection{27} It is quite impossible, Triarius, I said, not to say what you disapprove of in a man from whom you dissent. For what would stop me from being an Epicurean, if I found his doctrines convincing — especially since mastering them would be easy work? And so criticism between those who disagree with each other is not to be condemned; what in my judgment has no place in philosophy is abuse, insult, angry outburst, and that bitter persistence in dispute that refuses to yield.

\ciceroSection{28} Then Torquatus said: I agree entirely — for debate cannot proceed without mutual criticism, nor can it proceed rightly amid anger or obstinacy. But to all of this, unless it is troublesome to you, I have something I would like to say. Do you think, I replied, that I would have said what I said unless I wanted to hear you? Shall we then, he asked, run through the whole of Epicurus's system, or confine ourselves to the question of pleasure, which is the whole point of dispute? That is entirely for you to decide, I said. Very well, then, he said: I shall lay out one thing — the most important — and take up natural philosophy another time, when I shall also demonstrate to you, point by point, both the swerve of the atoms and the size of the sun and the many errors of Democritus exposed and corrected by Epicurus. For now I shall speak of pleasure — nothing novel, to be sure, but things I am confident you yourself will be persuaded by.

\ciceroSection{29} Certainly, I said, I shall not be contentious; and if you convince me of what you say, I will gladly agree. Convince you I shall, he said — provided only you show that fairness you have just displayed. I prefer, however, to speak in a continuous discourse rather than to put questions or be questioned. As you please, I said. Then he began. First, he said, I shall proceed as the founder of this school himself prescribes: I shall establish what kind of thing it is we are inquiring about, and what its nature is — not because I think you do not know, but so that our argument may move forward by method and system. We are inquiring, then, what the highest and ultimate good is — the thing which, in the judgment of all philosophers, must be such that everything is referred to it, while it itself is referred to nothing. Epicurus places this in pleasure, which he declares the highest good; he declares pain the highest evil; and he undertook to teach this as follows:

\ciceroSection{30} Every living creature, as soon as it is born, seeks pleasure and takes delight in it as the highest good, and recoils from pain as the highest evil and repels it as far as it can — and it does this while still uncorrupted, with nature herself pronouncing judgment purely and without distortion. And so, he says, no argument, no disputation is needed to show why pleasure is to be sought and pain avoided. He holds these things are felt directly — as fire is hot, snow is white, honey is sweet. None of these facts need elaborate proofs to establish; it is enough merely to call attention to them. For there is a difference, he holds, between a formal argument with its conclusions and a simple, ordinary observation and reminder: the former unfolds something hidden and, as it were, wrapped up; the latter passes judgment on things plain and open. Since, once we have stripped away the senses, nothing remains of the human being, it follows necessarily that what is in accordance with nature and what is against it must be judged by nature itself. And what does nature perceive, what does it judge, by which it might pursue or avoid anything, except pleasure and pain?

\ciceroSection{31} There are, however, some among our own school who hold that this account needs to be stated more precisely, and who deny that it is sufficient for the good and the bad to be judged by sensation alone — they hold that the mind and reason can also grasp that pleasure itself is desirable in and of itself, and pain itself to be avoided in and of itself. Accordingly they say there dwells in our minds a kind of natural and innate conception to the effect that we feel the one as something to be pursued and the other as something to be shunned. Others, on the other hand — and I agree with these — hold that, since many arguments are offered by many philosophers why pleasure should not be counted among good things and pain among bad, we ought not to place too much confidence in our case, but should argue it and set it out carefully and reason about pleasure and pain with closely marshaled arguments.

\ciceroSection{32} But so that you may understand the whole origin of that error made by those who denounce pleasure and praise pain, I shall lay the whole matter open and expound what that very discoverer of truth, that architect, so to speak, of the happy life, actually said. No one rejects pleasure itself because it is pleasure, or hates it or flees from it; but those who do not know how to pursue pleasure by the use of reason are overtaken by great pains — and no one loves pain itself because it is pain, or seeks it out, or desires to embrace it; but there are occasions when labor and pain bring some great pleasure in their train. To take the most trivial instance: who among us undertakes any laborious exercise of the body, unless to obtain some advantage from it? And who could rightly blame a man who desired to enjoy a pleasure in which no annoyance follows, or one who avoided a pain that produces no pleasure?

\ciceroSection{33} But we do find fault, and rightly hold them deserving of contempt, those who are so seduced and corrupted by the allurements of present pleasures that they are blinded by desire and fail to foresee the pains and troubles that will follow — and equally culpable are those who abandon their duties out of weakness of spirit, that is, out of flight from toil and pain. The distinction here is easy to make and readily at hand. In a free hour, when we have the option of choosing and nothing prevents us from doing what pleases us most, every pleasure is to be embraced and every pain to be driven away. But in certain circumstances, and in accordance with our obligations or the demands of the situation, it will often fall out that pleasures are to be declined and annoyances not refused. And so the wise man maintains this principle of selection in such matters: he declines certain pleasures in order to secure greater ones in return, or he endures certain pains in order to repel greater ones. Holding this view as I do, what reason have I to fear that I cannot fit our Torquati into it?

\ciceroSection{34} The Torquati whom you called to mind a moment ago, with impressive accuracy of recollection and in a spirit kindly and generous toward us — well, the recital of my ancestors' praises did not corrupt me, nor did it slow my readiness to reply. How do you interpret what those men did? Do you really suppose they charged an armed enemy, or showed such severity toward their own children and their own blood, with no thought whatever for their own interests, their own advantage? Even wild animals do not rush and rampage so blindly that we cannot understand where their movements and impulses are tending; do you think men of such extraordinary quality performed such great deeds without a reason?

\ciceroSection{35} What that reason was I shall examine presently. For now I hold to this: if they acted as they did for some reason — and these actions are beyond doubt glorious — then virtue was not in and of itself the reason for what they did. He stripped the torc from the enemy. — And he covered himself, so as not to be killed. — But he ran a great risk. — In the eyes of the whole army. — What did he gain from it? — Praise and affection — the surest safeguards for a life lived without fear. He put his son to death. — If for no reason, I should not wish to be descended from a man so harsh and so cruel; but if his purpose was to ratify the discipline of military command at the cost of his own grief, and to hold the army in check in a most serious war through fear of punishment, then he looked to the safety of his fellow citizens, in which he understood his own safety to be contained.

\ciceroSection{36} And this reasoning reaches far and wide. In the very ground on which your school's argument loves most to parade itself — yours especially, who make a point of tracing things back to antiquity — the commemoration of illustrious and courageous men and the praising of their deeds, not for some material benefit but for the very splendor of honorable conduct itself: all of that is overturned entirely by the principle of selection I described just now, whereby pleasures are foregone in order to secure greater pleasures, or pains are taken on in order to escape greater pains.



\ciceroSection{37} But enough has been said here about the famous and glorious deeds of great men; for there will soon be a more fitting place to discuss how every virtue makes its way toward pleasure. Now, however, I shall set out the nature and character of pleasure itself, so that all error may be removed from the uninformed and so that this school, which is taken to be devoted to pleasure, comfort, and softness, may be understood for what it truly is: serious, disciplined, and austere. For we do not pursue that pleasure alone which moves our very nature by some sweetness and is received by the senses with a certain delight; the pleasure we count greatest is that which is experienced when all pain has been stripped away. For since, when we are deprived of pain, we rejoice in the very liberation and the absence of all distress, and since everything in which we rejoice is pleasure, just as everything that causes us distress is pain, the absence of all pain has rightly been called pleasure. For just as, when hunger and thirst are expelled by food and drink, the very removal of the distress brings pleasure in its train, so in every case the departure of pain produces pleasure in its wake.

\ciceroSection{38} Epicurus therefore held that there is no middle state between pain and pleasure; for that very condition which some have supposed to be intermediate, as long as it is free from all pain, is not merely pleasure but the highest pleasure. For whoever is capable of feeling must, he held, be either in pleasure or in pain. Epicurus takes the highest pleasure to be bounded by the entire removal of pain, such that pleasure may thereafter vary and be differentiated, but cannot be increased or extended.

\ciceroSection{39} There is at Athens, as I heard my father relate with amusement and wit at the Stoics' expense, a statue in the Ceramicus of Chrysippus seated with his hand extended, the posture indicating that he took delight in this little argument: ``Does your hand, as it is now disposed, desire anything?'' ``Nothing whatever.'' ``But if pleasure were a good, it would desire it.'' ``So I suppose.'' ``Pleasure, therefore, is not a good.'' Even a statue, my father used to say, would never utter such a thing if it could speak. For the argument, though reasonably sharp against the Cyrenaics, touches Epicurus not at all. For if the only pleasure were that which, so to speak, tickles the senses — to use that expression — and flows upon them with sweetness and glides in upon them, then neither the hand could be content nor any other part of us with the mere absence of pain, unaccompanied by some pleasant movement of pleasure. But if, as Epicurus holds, the highest pleasure is to feel no pain, then, Chrysippus, your first concession was rightly made — that the hand as then disposed would desire nothing — but your second was not rightly made: that if pleasure were a good, the hand would desire it. For the very reason it would not desire it is that whatever is free from pain is already in pleasure.

\ciceroSection{40} That pleasure is the ultimate good is most readily seen from the following. Let us conceive of a man enjoying great, manifold, and unbroken pleasures of mind and body alike, with no pain either present or impending: what man, I ask, could we say is in a more enviable or more desirable state? For there must be present in a man so disposed both a firmness of mind that fears neither death nor pain — since death is devoid of sensation, and pain, when prolonged, tends to be light, and when severe, tends to be brief, so that the greatness of it is relieved by its swiftness and its duration is assuaged by its yielding.

\ciceroSection{41} Add to this that he feels no terror before the divine power, that he does not let past pleasures drain away but rejoices in the steady recollection of them: what, then, could be added to this to make it better? Set against him a man worn down by agonies of mind and body as severe as any that can fall to a human being, no hope before him that they will ever be lighter, no present or anticipated pleasure beside him — what more wretched state can be named or imagined? If, therefore, a life full of pains is the thing most of all to be fled from, then assuredly to live in pain is the highest evil, and the view that follows from this is that the highest good is to live in pleasure. For our mind has no place to stand firm at its very limit, and all fears and griefs are referred to pain; nor is there any other thing that by its own nature can disturb or distress us.

\ciceroSection{42} Furthermore, the beginnings of all approach and withdrawal, and of action in general, arise from either pleasure or pain. Since this is so, it is evident that all right and praiseworthy conduct is referred to the end of living in pleasure. And since this is the highest or final or ultimate good — what the Greeks call \textit{[Greek: telos]} — to which all other things are referred while it itself is referred to nothing further, we must acknowledge that the highest good is to live pleasantly. Those who place this highest good in virtue alone and are captivated by the splendor of the name, without understanding what nature demands, will, if they are willing to hear Epicurus, be freed from the gravest error. For those distinguished and beautiful virtues of yours — unless they produced pleasure, who would count them worth praising or worth pursuing? As the science of medicine is valued not for its own art but for the sake of good health, and as the art of the helmsman, because it embodies a sound method of navigation, is praised for its usefulness rather than for its art, so wisdom, which must be regarded as the art of living, would not be sought if it produced nothing; but it is sought now because it is, as it were, the craftsman for the securing and

\ciceroSection{43} acquiring of pleasure — and what I mean by pleasure, you now see clearly, lest hostility toward the word should undermine my argument. For since human life is chiefly troubled by ignorance of what is good and bad, and through this error men are often deprived of the greatest pleasures and tormented by the sharpest pains of the mind, wisdom must be brought to bear: wisdom, which, by stripping away dreads and desires and tearing out the recklessness of all false opinions, offers itself to us as the surest guide to pleasure. Wisdom alone it is that drives grief from our minds and keeps us from shuddering with fear. Under her instruction we can live in tranquility, once the burning force of all desires has been quenched. For desires are insatiable; they overturn not only individual persons but whole households, and often shake entire commonwealths to their foundations.

\ciceroSection{44} From desires are born hatreds, quarrels, discords, seditions, and wars; and these do not merely bluster outward or rush blindly against others, but, penned within the mind, fall into conflict and discord among themselves, from which life must necessarily be made most bitter. The wise man alone, therefore, by cutting away and trimming off all emptiness and error, can live within the limits nature sets, free from distress and free from fear.

\ciceroSection{45} For what division is more useful or better suited to living well than the one Epicurus employed? He placed in one class those desires that are both natural and necessary; in a second, those that are natural but not necessary; in a third, those that are neither natural nor necessary. The principle governing them is this: the necessary ones are satisfied with no great effort or expenditure; the natural ones, though not necessary, require little, because nature herself possesses resources that content her — resources both ready to hand and bounded; but for the empty desires no measure and no limit can ever be found.

\ciceroSection{46} And if we see that all of life is thrown into confusion by error and ignorance, and that wisdom alone can free us from the assault of passions and the terror of fears, teach us to bear the injuries of fortune in moderation, and show us every path that leads to repose and tranquility — then why should we hesitate to say both that wisdom is to be sought for the sake of pleasures and that folly is to be fled for the sake of the distresses it brings?

\ciceroSection{47} By the same reasoning we shall say that temperance also is not to be sought for its own sake, but because it brings peace to our minds and soothes and calms them with a kind of harmony. Temperance is what counsels us, in things either to be pursued or avoided, to follow reason. For it is not enough to judge what ought or ought not to be done; one must also stand firm in what has been judged. Most people, because they cannot hold and keep to what they themselves have resolved, are overcome and broken by the spectacle of some pleasure thrust before them and surrender themselves to be bound by their passions, failing to foresee what will follow; and for this reason, on account of a pleasure that is small and unnecessary, one that could have been obtained in another way and that they could even have done without without pain, they fall into serious illness, into loss, into disgrace, and are often entangled in the penalties of the laws and the courts.

\ciceroSection{48} But those who intend to enjoy pleasures in such a way that no pains follow them on their account, and who keep their judgment intact so as not to be overcome by pleasure and do what they know ought not to be done — these men achieve the highest pleasure by foregoing a pleasure. In the same way they often endure pain, lest by not doing so they fall into a greater pain. From this it is understood that intemperance is not to be avoided for its own sake, and that temperance is to be sought — not because it flees pleasures, but because it wins greater ones.

\ciceroSection{49} The same account will be found to hold for courage. For neither the discharge of labors nor the endurance of pains draws us by itself and in its own right, nor does patience, nor perseverance, nor sleeplessness, nor even that quality so admired — industry — nor courage itself; rather we pursue all these so that we may live free of anxiety and fear and may relieve our mind and body, as much as we can, from distress. For as the fear of death throws the whole settled order of a tranquil life into confusion, and as it is wretched to succumb to pain and bear it with a broken and feeble spirit — a weakness of soul through which many have utterly destroyed parents, friends, not a few their country, and very many themselves — so a strong and lofty mind is free of every anxiety and anguish: it despises death, since those who have met it are in the same condition as they were before they were born; it is so prepared for pain that it remembers the greatest pains are ended by death, that slight ones have many intervals of rest, and that over moderate ones we are ourselves the masters — so that, if they are bearable, we bear them, and if not, we may depart from life with equanimity, as from a theater where the performance no longer pleases. From all this it is understood that neither cowardice nor idleness is censured in its own name, nor courage and patience praised in their own name, but the former are rejected because they produce pain, the latter pursued because they produce pleasure.

\ciceroSection{50} Justice remains, so that we may have spoken of every virtue. Much the same things can be said here. I have shown that wisdom, temperance, and courage are so bound up with pleasure that they can in no way be torn or dragged from it; and the same judgment must be passed on justice, which not only never injures anyone but, on the contrary, always acts by its own power and nature both to settle the mind and to supply the hope that none of those things nature, undistorted, desires will be lacking. And just as rashness, passion, and idleness always torment the mind, always agitate and disturb it, so injustice, once it has settled in a mind, is turbulent by the very fact of its presence; and if it has already contrived something, however secretly done, it can never trust that it will remain forever hidden. Most commonly, suspicion follows the deeds of the wicked first, then talk and rumor, then an accuser, then a judge; and many, as was the case in your consulship, have turned witness against themselves.

\ciceroSection{51} And if some men seem to themselves sufficiently fortified and defended against the knowledge of other men, they still dread the gods, and believe that those very anxieties by which their minds are eaten away night and day are inflicted on them by the immortal gods as punishment. And what accretion from wicked acts can be great enough to diminish the distresses of life, compared with how great an increase they suffer, both from the consciousness of those acts and from the penalty of the laws and the hatred of one's fellow citizens? Yet in certain men there is no limit to money, to honor, to power, to lust, to feasting, or to the remaining desires, which no prize ever won by wickedness diminishes but rather inflames, so that they seem to stand in greater need of restraint than of reform.

\ciceroSection{52} True reason, then, invites the sound-minded to justice, equity, and good faith; nor does unjust action profit the weak and ungoverned man, who can neither easily accomplish what he attempts nor retain what he has accomplished. Wealth, whether of fortune or of intellect, is more suited to generosity: those who make use of it win goodwill for themselves and, what is most fitted to living in tranquility, genuine affection — especially since there is absolutely no reason to transgress.

\ciceroSection{53} For the desires that arise from nature are easily satisfied without any injury to another; but the empty ones are not to be obeyed, since they crave nothing that is genuinely desirable, and the harm in the injustice itself outweighs the benefit in whatever is gained by it. Therefore one would not be right to call even justice something desirable in itself: it is desirable rather because it brings the greatest abundance of delight. To be loved and held dear is pleasant, because it makes life safer and pleasure more complete. And so it is not only on account of the evils that befall the wicked that we think wickedness is to be fled, but far more because, in the mind of the man in whom it lodges, it never lets him breathe, never lets him rest.

\ciceroSection{54} And if even the praise of the virtues themselves — in which the rhetoric of all other philosophers chiefly exults — can find no outlet unless it is directed toward pleasure, and if pleasure alone calls us to itself and attracts us by its very nature, then it cannot be in any doubt that this is the highest and ultimate good of all, and that to live happily is nothing other than to live with pleasure.

\ciceroSection{55} To this settled and stable doctrine I shall explain briefly what is annexed. There is no error in the ends themselves — that is, in pleasure or in pain — but people go wrong in those matters when they are ignorant of the means by which these ends are brought about. The pleasures and pains of the mind, we freely admit, arise from the pleasures and pains of the body — and so I grant the point you were making just now: if any of our people think otherwise, they fall away, and I see indeed that there are many such, though they lack judgment. Yet even though pleasure of the mind brings us delight and pain brings us distress, each of these nonetheless has its origin in the body and is referred back to the body; and for that reason the pleasures and pains of the mind are not on that account any the less, but rather far greater, than those of the body. For by means of the body we can perceive nothing except what is present and at hand; but by means of the mind we perceive things both past and future. For even when we suffer equal pain of mind while we suffer in the body, a very great addition can still be made to that mental pain if we believe some eternal and boundless evil to be hanging over us. And the same thing may be transferred to pleasure — so that it is greater when we fear nothing of that kind.

\ciceroSection{56} This much, at any rate, is perfectly clear: the greatest pleasure or distress of the mind contributes more, whether toward a happy life or toward a miserable one, than either of those same things when it is equally prolonged in the body. We do not accept, however, that when pleasure is removed, pain immediately follows in its train — unless pain should happen to take the place of pleasure; but we do hold that we ourselves rejoice in the removal of pains, even when no pleasure that stirs the senses has succeeded in their place, and from this it can be understood how great a pleasure it is to feel no pain.

\ciceroSection{57} But just as we are raised up by the goods we expect, so we rejoice in those we remember. Fools are tormented by the recollection of evils; the wise are gladdened by dwelling again on past goods with grateful remembrance. It lies within our power to bury adversities in a perpetual oblivion and to remember our prosperities with pleasure and delight. When, however, we look back on past events with keen and attentive minds, it follows that distress attends them if they were bad, gladness if they were good. O what a splendid path to the happy life — how open, how plain, how straight! For when nothing can surely be better for a man than to be free from all pain and distress and to enjoy the fullest pleasures of both mind and body, do you not see how nothing is passed over that makes life better — the easier to attain that highest good which is set before us? Epicurus himself cries out — that man whom you say is devoted too much to pleasures — that a pleasant life is impossible unless one lives wisely, honourably, and justly, and that a life of wisdom, honour, and justice is impossible unless it is also pleasant.

\ciceroSection{58} For just as a city in sedition cannot be happy, nor a household discordant among its masters, so a mind at war with itself and in discord cannot taste any portion of pure and unconstrained pleasure. And one who is always engaged with conflicting and contrary desires and counsels can see nothing of calm, nothing of tranquillity.

\ciceroSection{59} And if the pleasantness of life is impeded by the graver diseases of the body, how much more must it be impeded by the diseases of the mind! The diseases of the mind are those boundless and empty desires for wealth, glory, power, and the pleasures of appetite. Then come anxieties, troubles, griefs, which eat men's minds away and wear them down with cares — these are men who do not understand that the mind need feel no pain that is detached from present or impending bodily pain. And there is no fool who does not labour under one or another of these diseases; therefore no one is free from wretchedness.

\ciceroSection{60} There is death besides — which hangs over them like the stone over Tantalus — and then superstition, by which one who is infected can never know rest. Add to this that they take no pleasure in remembering past goods, do not enjoy what is present, but merely wait upon the future; and because the future cannot be certain, they are consumed with anguish and fear, and are most tormented when at last they feel, too late, that they have spent themselves in vain in the pursuit of money or of power or of wealth or of glory. They lay hold of no pleasures, though it was in hope of obtaining these that they had taken on so many great and heavy toils.

\ciceroSection{61} And then there are the other sort: small, narrow souls who either despair of everything always, or are spiteful, envious, disagreeable, light-shunning, abusive, monstrous; and some again are enslaved to the trifling follies of erotic passion, others are wanton, others reckless, insolent — and at the same time intemperate and cowardly, never constant in any purpose. Because of all this there is not a moment's respite from distress in their lives. So it is that none of the foolish is happy and none of the wise is other than happy. And this we hold far better and more truly than the Stoics. For they deny that anything is good except some shadow of theirs they call the \textit{honestum} — a name more brilliant than solid — and they hold that virtue resting on this \textit{honestum} requires no pleasure at all, and is in itself sufficient for the happy life.

\ciceroSection{62} These things, however, can be said in a certain sense not only without our objection but even with our approval. For Epicurus portrays his wise man as always happy in this way: his desires are bounded, he disregards death, he holds true views about the immortal gods free of all fear, and does not hesitate to leave this life if it be better for him to do so. Furnished with these things, he is at every moment in a state of pleasure. Nor is there any moment at which his pleasures do not outweigh his pains. For he remembers the past with gratitude, takes hold of the present in such a way as to perceive how great and how delightful these things are, does not hang suspended upon the future but awaits it; he enjoys the present, and is as far as possible removed from those vices I catalogued a moment ago; and when he compares the life of fools with his own, he is filled with great pleasure. As for whatever pains do come his way, they never hold enough force to prevent the wise man from having more to rejoice over than to be distressed by.

\ciceroSection{63} Epicurus is entirely right, moreover, in saying that fortune plays a small role in the affairs of the wise man, and that the greatest and most serious matters are managed by his own counsel and reason, and that no greater pleasure can be drawn from time unlimited than from this time we can see to be finite. In dialectic, your discipline, he judged there lay no path either toward living better or toward arguing more rigorously. In natural philosophy he placed the greatest weight. That study brings to light the force of words, the nature of discourse, and the relation of things that follow from or conflict with one another. By knowing the nature of all things we are relieved of superstition, freed from the fear of death, and not thrown into confusion by ignorance of things — an ignorance that often gives rise to terrible terrors. And finally we shall even be better in our conduct when we have learned what nature requires. And then, when we hold a firm knowledge of things, guided by that rule which has, as it were, descended from heaven for the understanding of all things — a rule by which all judgments concerning things are directed — we shall never, overcome by anyone's argument, be turned from our settled view.

\ciceroSection{64} For unless the nature of things has been thoroughly investigated, there is no way we shall be able to defend the judgments of the senses. Moreover, everything we discern with the mind arises entirely from the senses; and if all the senses are true — as Epicurus's account teaches — then at last something can be known and grasped. Those who remove the senses and say that nothing can be perceived cannot even set out, without the senses, the very point they are arguing. Furthermore, when knowledge and understanding are done away with, all method for conducting life and managing affairs is done away with as well. Thus from natural philosophy we draw courage against the fear of death, steadiness against the dread of religion, the settling of the mind through the removal of ignorance of all hidden things, self-mastery through the setting out of the natural kinds of desire, and — as I explained just now — the standard of understanding and the distinction of true from false established by the judgment derived from that same standard.

\ciceroSection{65} There remains the topic most necessary of all to this inquiry: friendship — which, you maintain, will not exist at all if pleasure is the highest good. On this matter Epicurus says that of all the things wisdom procures to serve the happy life, nothing is greater than friendship, nothing more abundant, nothing more delightful. He proved this not merely in words but far more by his life, his deeds, and his character. How great a thing this is the invented tales of the ancients declare: in all those stories, numerous and varied, reaching back to remotest antiquity, you can barely find three pairs of friends — even if you set out from Theseus and come all the way to Orestes. Yet Epicurus, in a single household, and that a modest one, what great throngs of friends he held together, and with what harmony of devoted affection! This practice continues among Epicureans to this day. But let us return to the matter; there is no need to speak of individuals.

\ciceroSection{66} I see, then, that our people have argued about friendship in three ways. Some, when they denied that pleasures belonging to one's friends are to be sought for their own sake as eagerly as we seek our own — a position at which some think the stability of friendship is shaken — nonetheless defend this position and extricate themselves from it easily, or so it seems to me. For just as they argue that virtue, of which I spoke earlier, cannot be separated from pleasure, so they argue the same about friendship: since a solitary life and a life without friends is full of ambushes and fear, reason itself urges us to form friendships, the acquisition of which strengthens the spirit and cannot be detached from the hope of obtaining pleasures.

\ciceroSection{67} And just as hatred, envy, and contempt stand in the way of pleasures, so friendships are not only the most loyal supporters of pleasure but its actual producers — both for friends and for oneself; and we do not only enjoy them in the present but are lifted up by hope of the time that follows and comes after. And because we can in no way maintain a firm and lasting pleasantness of life without friendship, nor indeed preserve friendship itself unless we love our friends as much as ourselves — for this very reason this result is brought about in friendship, and friendship is joined together with pleasure. For we rejoice in a friend's joy as much as in our own, and we are pained equally by his distress.

\ciceroSection{68} Therefore the wise man will be disposed toward his friend as toward himself; and whatever toils he would undertake for the sake of his own pleasure, the same toils he will undertake for the sake of his friend's pleasure. And what was said about virtues — how they are always attached to pleasures — the same must be said about friendship. Epicurus puts it admirably in almost these very words: the same understanding, he says, that has confirmed the mind against any fear of evil, whether everlasting or prolonged, has also discerned that within the very span of this life the surest protection is friendship.

\ciceroSection{69} There are, moreover, certain Epicureans who are a little too timid in the face of your attacks — though sharp enough in their reasoning — and who fear that if we hold friendship to be sought for the sake of our own pleasure, the whole of friendship may seem, as it were, to limp. And so they say that the first meetings, joinings-together, and dispositions to form acquaintances arise because of pleasure; but that when long practice has produced intimacy, love then blossoms to such a degree that friends come to be cherished for their own sake even if no utility accrues from the friendship. Indeed, if through habit we are accustomed to love places, shrines, cities, gymnasia, playing-fields, dogs, horses, and the customary exercises of sport or of the hunt, how much more readily and with how much more justification can this happen with the intimacy of human beings?

\ciceroSection{70} There are others, again, who say there is a kind of compact among the wise, whereby they love friends no less than themselves. We understand that this can happen, and we often see it happen; and it is plain that nothing better suited to a pleasant life can be found than such a union. From all of this it may be judged that the rational account of friendship is not merely not impeded if the highest good is placed in pleasure, but that without this founding principle no account of friendship can be found at all.

\ciceroSection{71} Wherefore, if the things I have said are brighter and clearer than the sun itself, if all of it has been drawn from the wellspring of nature, if our entire argument is confirmed in all its credibility by the senses — that is, by witnesses uncorrupted and intact — if even infant children, even the very beasts, in their inarticulate way, declare with nature as their teacher and guide that nothing is prosperous except pleasure and nothing harsh except pain, and if on these matters their judgment is neither depraved nor distorted — surely we ought to feel the greatest gratitude to him who, having, as it were, heard this voice of nature, grasped it with such firm conviction and such weight that he leads all persons of sound mind onto the path of a calm, tranquil, quiet, and happy life. And the fact that he seems to you insufficiently learned has this cause: he judged that no learning deserves the name except that which advances the discipline of the happy life.

\ciceroSection{72} Would he have spent time in unrolling the poets, as Triarius and I do at your urging — poets who offer no solid use and whose charm is entirely for children? Or would he have worn himself out, as Plato did, on music, geometry, arithmetic, and astronomy, pursuits that start from false premises and for that reason cannot be true, and that even if they were true would contribute nothing to make our lives more pleasant — that is, better? Would he have devoted himself to those arts, and left aside the art of living, so great, so laborious, and as productive as it is demanding? No: Epicurus is not unlearned — it is those others who are untaught, those who think that what is shameful not to have learned as children must be learned right on into old age. When Torquatus had finished speaking, he said: I have set out my view — and I have done so with the deliberate purpose of learning yours, since the opportunity to do this at my own convenience has never before been given to me.


\ciceroBook{2}

\ciceroSection{1} While both of them were looking at me and signalling that they were ready to listen, I began. First of all, I said, I ask you not to think of me as someone about to lay out a formal lecture in the manner of a schoolroom — a practice I have never greatly admired even in the philosophers themselves. For when did Socrates, who can rightly be called the father of philosophy, ever do any such thing? That was the custom of those who were then called sophists, and the first of that company to dare demand a question from an assembled audience — that is, to invite anyone present to name whatever subject he wished to hear discussed — was Gorgias of Leontini. A bold business, and one I should call shameless, had this institution not afterward been taken over by our own philosophers.

\ciceroSection{2} But both that Gorgias I have named and the other sophists — as can be understood from Plato — we see played with by Socrates. For his practice was to draw out by repeated questioning the opinions of those with whom he was conversing, so that, in response to their answers, he could say what he thought fit. This practice was not kept up by later philosophers; but Arcesilas revived it and established the rule that those who wished to hear him should not ask questions of him, but should themselves state what they thought. When they had done so, he argued against it. But his audience would defend their own view as long as they could. Among the other philosophers, on the other hand, whoever raises a question is silent thereafter; and indeed this is now the way even in the Academy. For when the man who wishes to hear has said, for instance, ``Pleasure seems to me to be the highest good,'' the refutation is conducted in a continuous discourse, so that one can easily understand that those who claim something seems so to them do not actually hold that view themselves, but wish to hear the opposing arguments. We conduct things in a more convenient fashion.

\ciceroSection{3} For Torquatus not only stated his position but also gave his reasons for it. Now I think — though I was greatly delighted by his unbroken discourse — that it is nonetheless more convenient, when one presses on each point in turn and understands what each man concedes and what he denies, to draw from the conceded points whatever conclusion one wishes and so to reach a result. For when a discourse rushes forward like a torrent, however much it sweeps along of every kind, there is nothing you can hold, nothing you can lay hands on, nowhere you can check the headlong flow. In all inquiry that is conducted by some method and system, the first requirement of the discourse is to set down, as it were in a legal formula, the terms under which the matter shall be conducted — so that those who are arguing together may agree on what it is they are arguing about.

\ciceroSection{4} This principle was laid down in Plato's \textit{Phaedrus} and approved by Epicurus, who recognized that it ought to be observed in every discussion. But what followed immediately from it he did not see. For he rejects definition as a practice he finds disagreeable — without which it is sometimes impossible for those in dispute to agree on what the subject of their debate actually is, as is the case in this very matter we are now discussing. We are inquiring into the end of goods. Can we know what that end is like unless we first compare notes on what we mean by ``end'' when we say ``end of goods,'' and what we mean by ``good'' itself?

\ciceroSection{5} Yet this uncovering of things that are hidden — when the nature of each thing is disclosed — is precisely what definition is. And you were making use of it yourself, though without noticing it, on several occasions. For you were defining this very end, or boundary, or ultimate point as that to which all right actions are referred but which is itself referred to nothing beyond it. A fine definition. You would perhaps also have defined what the good itself is, had there been need — defining it either as that which is by nature to be sought, or as that which profits, or as that which pleases, or as that which frees from pain. Now, since it does not displease you to define things, and you do so when you choose, I should like you, if it is no trouble, to define what pleasure is, since that is the whole subject of this inquiry. ``Who, pray,'' he said, ``does not know what pleasure is, or who stands in need of any definition in order to understand it more fully?''

\ciceroSection{6} I myself, I said, would have claimed that I did not know — except that I fancy I have pleasure well understood, firmly grasped and held in my mind. But as things stand, I say that Epicurus himself does not know it, and wavers on the point, and that a man who so frequently insists that we must carefully express the meaning we attach to words does not himself always understand what meaning this word ``pleasure'' carries — that is, what thing is signified by it. Then he laughed and said: ``This really is splendid — that the man who declares pleasure to be the end of all things to be sought, the ultimate and supreme good, should not know what pleasure itself is and of what quality!'' ``Well then,'' I said, ``either Epicurus does not know what pleasure is, or every mortal in the world does not know it.'' ``How so?'' he said. ``Because pleasure is something that everyone recognizes as a movement perceptible to the senses that floods them with a certain agreeableness.'' ``Well?''

\ciceroSection{7} ``Does Epicurus not know that pleasure?'' he said. ``Not always,'' I said; ``for sometimes he knows it all too well — as when he testifies that he cannot so much as conceive of where any good might be, or what it might be, apart from that which is taken in through food and drink and the delight of the ears and obscene pleasure.'' ``Are these his words?'' ``As if I were in the least embarrassed by those passages,'' he said, ``or could not make plain in what sense they are said!'' ``I have no doubt,'' I said, ``that you can do so with ease, and there is nothing you need be embarrassed about in agreeing with a wise man — who alone, so far as I know, has dared to call himself wise. For Metrodorus, I think, did not himself claim the title; he was unwilling to refuse so great a compliment when Epicurus bestowed it upon him. As for the Seven, they were named wise men not by their own vote but by the suffrage of all peoples.''

\ciceroSection{8} But at this point I grant, he said, that Epicurus knows the same force of pleasure that everyone else does. For all men call by the name of pleasure \textit{[Greek: hedone]} in Greek, \textit{voluptas} in Latin, the agreeable movement by which the senses are gladdened. ``What then is it you are looking for?'' he said. ``I will tell you,'' I said, ``and I say it more for the purpose of learning than because I wish to find fault with you or with Epicurus.'' ``And I too,'' he said, ``would rather learn something if you bring it forward than criticize you.'' ``Do you know, then,'' I said, ``what Hieronymus of Rhodes says the highest good is, to which he thinks everything ought to be referred?'' ``I know,'' he said: ``his view is that the end is freedom from pain.'' ``Well? And what does that same man think about pleasure?''

\ciceroSection{9} ``He denies,'' he said, ``that pleasure is to be sought for its own sake.'' ``He therefore considers rejoicing to be one thing, and not being in pain to be another.'' ``And indeed he is very badly wrong,'' he said; ``for, as I explained a little while ago, the end of increasing pleasure is the removal of all pain.'' ``Whether that `freedom from pain' amounts to what you claim,'' I said, ``I will consider later. But that pleasure is one thing and freedom from pain another — this you must concede, unless you are going to be very stubborn.'' ``You will find me stubborn on this point at least,'' he said; ``for nothing truer can be said.'' ``Is there pleasure,'' I asked, ``in the drinking for one who is thirsty?'' ``Who,'' he said, ``could deny that?'' ``Is it the same pleasure as when the thirst has been quenched?'' ``On the contrary, it is of a different kind,'' he said; ``for when the thirst is quenched the pleasure is one of stable rest, whereas the pleasure of quenching the thirst itself is in motion.'' ``Why then,'' I said, ``do you call things so dissimilar by the same name?'' ``What about what I said just now?'' he said. ``Do you not remember? — that when all pain is removed, pleasure is varied, not increased.''

\ciceroSection{10} ``I do remember,'' I said; ``but what you just said was good Latin and not entirely clear. For \textit{varietas} is a Latin word, properly used of different colours, but transferred to many different things: a poem of varied texture, a speech of varied style, varied characters, varied fortune — pleasure too is commonly called varied when it is derived from many dissimilar things that produce dissimilar pleasures. If you meant variety in that sense, I would follow you — as indeed I do follow you even without your saying it. But what this variety of yours is, I do not see clearly enough: your claim that when we are free from pain we are at the pinnacle of pleasure, and then that when we are partaking of things that bring a sweet movement to the senses we are in that kind of pleasure which is in motion and which produces the variety of pleasures, but does not increase that pleasure of not being in pain — why you should call that last thing pleasure, I do not understand.'' ``Is there anything,'' he said, ``that can be more agreeable than to feel no pain?''

\ciceroSection{11} ``Let nothing be better, by all means,'' I said — ``for that is not yet what I am asking — but is it therefore the same thing as pleasure? Is what I might call `painlessness' the same as pleasure?'' ``Entirely the same,'' he said, ``and the greatest pleasure at that — greater than which none can exist.'' ``Why then do you hesitate,'' I said, ``given that you have established the highest good as consisting entirely in not being in pain, to hold to that single thing, to defend it, and to maintain it?''

\ciceroSection{12} For what need is there to bring pleasure into the company of the virtues — like a harlot introduced into a gathering of matrons? The name of pleasure is invidious, disreputable, suspect. And so you are fond of saying that we do not understand what Epicurus means by pleasure. Which, whenever it has been said to me — and it has been said not a few times — I confess that, easy-tempered as I am in argument, I sometimes feel a flicker of irritation. Do I not understand what \textit{[Greek: hedone]} means in Greek, what \textit{voluptas} means in Latin? Which of the two languages is it that I do not know? Besides, how is it that I do not know and all those who have chosen to be Epicureans do? Which indeed is a point your school makes very well: that a man who is going to be a philosopher need know no literature. And so, just as our ancestors brought that celebrated Cincinnatus from his plough to be dictator, so you collect from the villages fine fellows to be sure, but assuredly not men of much learning.

\ciceroSection{13} Do those men, then, understand what Epicurus says, while I do not? That you may know I understand it: I say in the first place that \textit{hedone} \textit{[Greek: hedone]} and \textit{voluptas} are the same thing. Indeed, we often look for a Latin word that exactly matches a Greek one and expresses the same meaning; here there was nothing to look for. No word can be found that declares in Latin what \textit{hedone} declares in Greek more exactly than \textit{voluptas} declares it. To this one word, all who know Latin attach two things: gladness in the mind, and a sweet, agreeable movement of sensation in the body. For the character in Trabea speaks of the pleasure of the mind as great gladness, the same gladness as the character in Caecilius who declares that he is glad with every kind of gladness. But here is the difference: \textit{voluptas} can be used of the mind — a flawed thing, as the Stoics think, who define it as ``an irrational elation of the mind, supposing it is enjoying some great good'' — whereas \textit{laetitia} and \textit{gaudium} are not used of the body.

\ciceroSection{14} In ordinary Latin usage, however, \textit{voluptas} is set down for the agreeableness experienced when something acts upon one of the senses. This agreeableness too, if you like, carry over into the mind; for ``to please'' is said of both. Consider this: between the man who says ``My joy is so overwhelming that I am beside myself'' and the man who says ``Now at last my soul is on fire'' — one being carried away by gladness, the other tormented by pain — there is a middle condition: ``Though this acquaintance between us is quite recent,'' where the speaker is neither rejoicing nor distressed. So too between the man who enjoys the bodily pleasures he has desired and the man who is racked by extreme pain there is the man who is free of both.

\ciceroSection{15} Do I seem, then, to have a sufficient grasp of the meaning of words, or must I still be taught either Greek or Latin? And yet consider this: if I fail to understand what Epicurus is saying, despite what appears to be my perfectly good knowledge of Greek, there may be some fault in the man who speaks in such a way as not to be understood. This happens without blame in two ways: either if one does it deliberately, as Heraclitus did — who ``bears the surname Obscure \textit{[Greek: skoteinos]}, because his account of nature was too dark'' — or when the obscurity comes from the difficulty of the subject rather than of the words, as is the case in Plato's \textit{Timaeus}. Epicurus, however, in my judgment neither wishes, if he can help it, to speak otherwise than plainly and openly, nor is he dealing with a matter that is obscure in the way natural philosophers treat it, or technical in the way mathematicians do; he is speaking of something clear and plain and already spread among ordinary people. Yet you do not deny that we understand what pleasure is; you only say we do not understand what he means by it. And from this it follows, not that we fail to understand the force of the word, but that he speaks in his own peculiar way and disregards ours.



\ciceroSection{16}
For if Epicurus is saying the same thing as Hieronymus — whose verdict is that the highest good is to live free of all distress — then why does he prefer to speak of pleasure rather than of the absence of pain, as Hieronymus does, the man who actually understands what he is saying? But if he means that pleasure in the kinetic sense should be added — for that is what he calls this sweet pleasure: ``in motion,'' while the state of feeling no pain he calls ``in stable equilibrium'' — then what is he driving at? He cannot make it seem to anyone who knows himself, that is, who has looked clearly into his own nature and perception, that freedom from pain and pleasure are one and the same thing. This, Torquatus, is to do violence to the senses, to wrench from our minds those concepts of words that are woven into us from infancy. For who does not see that nature presents us with exactly three conditions? The first, when we are in pleasure; the second, when we are in pain; the third — the one I am in at this moment, and I believe you are too — neither in pain nor in pleasure. A man feasting is in pleasure; a man on the rack is in pain. Do you really not see the vast multitude of people standing between these two extremes, neither rejoicing nor suffering?

\ciceroSection{17}
``Not at all,'' he said; ``and I maintain that everyone who is free from pain is in pleasure — and indeed in the highest pleasure.'' So the man who mixes a drink for someone else without being thirsty himself, and the man who drinks it down while he is parched — you put them in the same pleasure? At this Torquatus said: ``I asked you at the start, if you recall, to stop these dialectical traps, and that is precisely why I said from the beginning that I preferred this approach — I saw this coming.'' ``Then you would rather,'' I said, ``that we argue in the manner of orators than of logicians?'' ``As if,'' he replied, ``continuous discourse belonged only to orators and not to philosophers as well.'' ``That is Zeno's point,'' I said, ``the Stoic. He held, as Aristotle had before him, that every power of speaking is distributed into two parts: rhetoric he likened to the open palm, dialectic to the closed fist — because orators speak in a wider range while logicians speak in a more compressed one. Very well, I shall comply with your preference and speak, if I can, in the manner of orators — but the rhetoric of philosophers, not that forensic kind of ours, which is obliged, when it addresses a popular audience, to be somewhat blunter.''

\ciceroSection{18}
But while Epicurus despises dialectic, Torquatus — the one discipline that holds within it all knowledge of what each thing is and all power to judge the character of each, and the whole systematic art of reasoned argument — he runs headlong in his writing, as it seems to me at least, and fails to distinguish with any method the things he means to teach, as in the very passage we have just been discussing. The highest good, on your account, is pleasure. Then you must make clear what pleasure is; otherwise the question under investigation cannot be resolved. Had he done so, he would not have been left in such difficulty. Either he would have defended the pleasure that Aristippus defended — the kind by which the senses are sweetly and agreeably moved — which even animals, if they could speak, would call pleasure; or, if he preferred to speak in his own idiom rather than as all the Greeks of Athens and Mycenae and the rest, cited together in that anapaest, would speak, he would have called only this condition of feeling no pain ``pleasure'' and dismissed the Aristippean variety; or, if he approved of both, as he does, he would have joined freedom from pain with pleasure and operated with two distinct ultimate goods.

\ciceroSection{19}
Many great philosophers have in fact made such combinations of ultimate goods: Aristotle joined the exercise of virtue with prosperity in a life lived to its full completion; Callipho added pleasure to moral worth; Diodorus added freedom from pain to that same moral worth. Epicurus would have done the same if he had combined this position — now held by Hieronymus — with the old position of Aristippus. As it is, those two are in disagreement with each other. Each therefore uses a single terminus. Both are superb in their command of Greek, and neither Aristippus, who declares pleasure the highest good, includes freedom from pain within pleasure, nor does Hieronymus, who sets the highest good as freedom from pain, ever use the word ``pleasure'' for that painlessness — since he does not even count pleasure among things to be pursued.

\ciceroSection{20}
The two things are themselves distinct, not merely in name. The one is being without pain; the other is being in pleasure. You are attempting to make not only a single name out of these two utterly dissimilar things — which I could more easily tolerate — but a single thing, which is impossible. Epicurus, who approves of both, ought to have used both, just as he does in substance, even if he fails to separate them in words. For in many passages, while praising that pleasure which we all call by the same name, he ventures to say that he cannot so much as conceive of any good separated from the Aristippean kind of pleasure; and he says this in the very work where his whole discussion is devoted to the highest good. Yet in another book — the one in which he is said to have issued the weightiest of his maxims in compressed form, as if uttering the oracles of wisdom — he writes in these words, which you surely know, Torquatus; for which of your circle has not learned by heart Epicurus's \textit{[Greek: kyriai doxai]}, that is, his ``principal doctrines,'' as one might render it, the authoritative and supreme judgments, brief formulas for a blessed life on account of their gravity? Observe, then, whether I construe this maxim correctly:

\ciceroSection{21}
``If the things that produce pleasures for the dissolute were to free them from their fear of gods and death and pain, and were to teach them the limits of desire, we should have nothing to find fault with, since from every side they would be filled with pleasures and would have no part of anything painful or distressing — and that is what constitutes the bad.'' At this point Triarius could not contain himself. ``I ask you,'' he said, ``Torquatus, does Epicurus actually say that?'' He seemed, knowing perfectly well that he did, to want nonetheless to hear Torquatus admit it. And Torquatus, entirely unruffled and with real composure, replied: ``In precisely those words; but you do not see what he means by it.'' ``If he means something different from what he says,'' I replied, ``I shall never understand what he does mean; but he states plainly what he understands. And if he means what he says — that the dissolute are not to be censured, provided they are wise — then what he says is absurd; just as it would be absurd to say that parricides are not to be censured, provided they feel no greed and no fear of gods or death or pain. And yet what is the point of granting any exception to the dissolute at all, or of imagining a category of persons who, while living in debauchery, would on those terms alone not be censured by the greatest of philosophers, provided they took care in other respects?''

\ciceroSection{22}
But surely, Epicurus, you would censure the dissolute for this very reason — that they live in pursuit of pleasures of every description — especially since, on your own account, the highest pleasure is feeling no pain? And yet we shall find that profligates are, in the first instance, not so scrupulous about religion that they eat off sacred dishes, and then not so afraid of death that they have that line from \textit{Hymnides} forever on their lips: ``Six months of life are enough for me; the seventh I pledge to Orcus.'' As for Epicurus's remedies against pain, they will produce them readily enough, like medicines from a physician's case: ``If it is severe, it is brief; if it is long, it is mild.'' The one thing I cannot see is how, if a man is dissolute, he can possibly keep his desires within bounds.

\ciceroSection{23}
What then is the point of saying: ``I should find nothing to censure, if they kept their desires within bounds''? That is equivalent to saying: ``I would not censure profligates if they were not profligates.'' By the same reasoning you would not censure the dishonest either, if they were men of good character. This stern philosopher does not consider self-indulgence reprehensible in itself — and by Hercules, Torquatus, to speak the truth, if pleasure is the highest good, he is entirely right not to. For do not present me with the kind of profligates you people usually conjure up — the ones who vomit onto the table, who are carried out of banquets, who stuff themselves again the next day in the same state of overindulgence, who, as the saying goes, have never seen either a sunset or a sunrise, who run through their inheritances and are left destitute. None of us imagines that profligates of that stamp live pleasantly. Picture instead the refined and polished sort: the best cooks and bakers, fine fish and game and hunt, all their pleasures exquisitely chosen; the wine poured from a full jar, that golden-hued wine Lucilius describes, from which no ladle and straining-cloth have yet taken anything; musicians at table and all that follows — those pleasures, I mean, without which Epicurus says he cannot so much as conceive what the good is; let there be beautiful boys to serve, and linen and silver and Corinthian bronzes to match, the room itself, the building — these profligates I will never call happy or even well off.

\ciceroSection{24}
The conclusion that follows is not that pleasure is not pleasure, but that pleasure is not the highest good. Nor was Laelius — who in his youth had heard Diogenes the Stoic and later Panaetius — called the Wise because he was unaware of what is most delightful. It does not follow that a man whose heart is discerning has no discerning palate; it is rather that he counted such things as trivial. ``O sorrel, how you flaunt yourself — and you are not sufficiently known for what you are!'' On this very score Laelius used to call out, addressing our gluttons one by one, what that famous \textit{[Greek: sophos]} applauds — the shrewd observer: ``O Publius, O Gallonius, you guzzler, you wretch! You have never dined well in your life, for all that you spend everything on prawn and on giant sturgeon.'' This from a man who, setting no value on pleasure, denies that one who places everything in pleasure has dined well, yet does not deny that Gallonius has ever dined pleasurably — that would be a lie — but well. In this way, with gravity and seriousness, he severed pleasure from the good. Which yields the conclusion: all who dine well dine pleasurably, but those who dine pleasurably do not thereby automatically dine well. Laelius always dined well.

\ciceroSection{25}
What does ``well'' mean? Lucilius will explain: ``with food well prepared and well seasoned''; but offer the centrepiece of the dinner: ``with good conversation''; and what did that produce? ``Pleasure, if you ask''; for he came to dinner to satisfy in tranquility the natural wants of the body. Rightly, then, does Laelius deny that Gallonius ever dined well, and rightly call him wretched, especially since he spent his every effort on that pursuit. No one denies he dined pleasurably. Why then not well? Because what is well is also rightly done, frugally, honorably; whereas that man dined badly, depravedly, wickedly, disgracefully; therefore not well. Laelius was not ranking the flavor of sorrel above the flavor of Gallonius's sturgeon; he was setting no store by flavor at all — which he could not have done if he placed the highest good in pleasure.

\ciceroSection{26}
Pleasure, then, must be set aside — not only so that you may follow the right, but so that you may speak decently and with restraint. Can we possibly say that pleasure is the highest good in life, when we cannot seem to say it even of a dinner? Now consider how the philosopher himself speaks. He posits three classes of desire: natural and necessary; natural but not necessary; neither natural nor necessary. This first division is inelegant: what were two classes he has made into three. That is not dividing but fracturing. Those who have learned what he holds in contempt set it out as follows: two classes of desire, the natural and the empty; and of the natural, two subdivisions, the necessary and the not necessary. The thing would then be settled. To count a subdivision as a separate class within a division is a logical error.

\ciceroSection{27}
But let us pass over that. He despises elegance of argument; he speaks loosely. We must bear with his manner, provided his meaning is sound. And there is this further thing, which I cannot fully approve of and only just tolerate — a philosopher speaking of ``limiting'' desires. Can a desire be limited? It ought to be uprooted and extracted by the roots. For anyone who harbors a desire can rightly be called a man of desire. So there will be the miser — but in moderation; the adulterer — but with restraint; the profligate — in the same measure. What sort of philosophy is it that does not eradicate depravity but contents itself with mediocrity of vice? Though on the substance of this classification I entirely approve: I only miss the elegance. Let him call these the ``natural wants'' and reserve the word ``desire'' for another context — to deploy it, with all its force, when he speaks of greed, of intemperance, of the greatest vices, as though accusing them of a capital offence.

\ciceroSection{28}
But these points he makes with rather too much freedom and too often. That, for my part, I do not find fault with: it befits a great philosopher of such renown to defend his doctrines boldly. Yet because he appears at times to embrace with considerable fervor that pleasure which all peoples call by that name, he falls into serious difficulties: with public conscience set aside, there seems to be nothing so base that he would not appear ready to do it for pleasure's sake. Then, when he blushes — and nature's force is very powerful — he retreats to the position that nothing can be added to the pleasure of one who feels no pain. But that condition of feeling no pain is not called pleasure. ``I am not troubled by the name,'' he says. What of the fact that the thing itself is entirely different? I could find many people — countless, rather — less fastidious and less contentious than you are, whom I could easily convince of my point. Why then do we hesitate to conclude that if feeling no pain is the highest pleasure, then being in pain is the greatest suffering? Because pain and pleasure are not contraries; what is contrary to pain is the removal of pain.

\ciceroSection{29}
Not to see this is the strongest possible evidence that the pleasure he speaks of — the pleasure without which he says he cannot conceive of the good at all, the pleasure he elaborates as follows: what is perceived by the palate, what by the ears; and he adds the rest, things whose very mention requires a word of apology — this pleasure, the only good that this stern and weighty philosopher recognizes, he himself fails to see is not even something to be pursued, since on his own authority we do not feel the lack of that pleasure when we are free of pain.

\ciceroSection{30}
What contradictions these are! Had he learned to define, to divide, to command the force of language, and above all to hold to the received usage of words, he would never have stumbled into such rough ground. As it is, you see what he does. The thing that no one has ever called pleasure, he calls pleasure; two distinct things he makes into one. This kinetic pleasure — for so he styles those sweet and, as it were, honeyed pleasures — he sometimes disparages so severely that you would think you were hearing Manius Curius; at other times he praises it so lavishly that he declares he cannot even conceive of anything else being good. Language of this kind ought to be suppressed not by some philosopher but by a censor. For the fault lies not in the argument alone but in the morals. He does not censure self-indulgence, provided only it is free from unlimited desire and from fear. At this point he seems to be recruiting disciples: those who wish to live as profligates should first become philosophers.

\ciceroSection{31} The origin of the highest good, I take it, is sought from the first stirring of animate creatures at birth. The moment an animal is born, it takes pleasure in pleasure and pursues it as a good, and recoils from pain as an evil. Now he holds that the best judges of good and evil are animals not yet corrupted. This is your own position, set out in your own words. How much is wrong with it! For by which pleasure will a wailing infant adjudicate the highest good and evil — the static pleasure or the kinetic? Since, if the gods approve, we are learning how to speak from Epicurus: if the static, then nature plainly demands the preservation of itself, which we grant; but if the kinetic — as you do assert — then no base pleasure will be so vile as to deserve rejection; and at the same time that newly born creature is not taking its start from the supreme pleasure, the one you placed in the absence of pain.

\ciceroSection{32} And yet Epicurus did not derive this argument from infants, nor even from beasts, which he holds to be mirrors of nature, in order to say that under nature's guidance this pleasure of painlessness is pursued by them. For this static condition of being without pain cannot stir the appetite of the soul, nor does it carry any force to impel it. It is the other pleasure that impels — the one that strokes the senses with delight. And so Epicurus invariably relies on this point to prove that pleasure is pursued by nature: that the pleasure which consists in motion draws both infants and beasts toward itself, whereas that stable pleasure, which amounts to nothing more than the absence of pain, does not. How then does it follow that nature takes its start from one pleasure, yet the highest good is placed in another?

\ciceroSection{33} As for beasts, I hold that their judgment carries no weight. For though they may not have been corrupted, they can still be defective. Just as one rod is bent and curved by deliberate art and another was born that way, so the nature of wild creatures is not corrupted by bad training but is defective by its own nature. Moreover, it is not pleasure that nature moves an infant to pursue, but only self-love — the desire to be whole and safe. For every animal, from the moment it is born, loves itself and all its parts, and embraces above all those two capacities that are greatest in it: mind and body; and next the several parts of each. For there are in both mind and body certain preeminent faculties, and once an animal has dimly recognized these it begins to distinguish, so that it pursues those things first given it by nature and repels their contraries.

\ciceroSection{34} Whether pleasure is included among these primary natural endowments is a great question. To suppose, however, that there is nothing among them except pleasure — no limbs, no senses, no movement of intelligence, no bodily integrity, no bodily health — strikes me as the height of stupidity. And from that starting-point the whole account of good and evil must necessarily flow. To Polemo, and before him to Aristotle, those things appeared primary which I mentioned a moment ago. From that conviction sprang the teaching of the old Academics and the Peripatetics, that the end of goods is to live in accordance with nature — that is, to enjoy with the aid of virtue those things first given us by nature. Callipho joined to virtue nothing but pleasure; Diodorus, the absence of pain. … All the ends I have named flow as consequences from their respective principles: for Aristippus the end is simple pleasure; for the Stoics, agreement with nature, which they take to mean life governed by virtue — that is, to live honorably — and they interpret this as: to live with understanding of those things that come about by nature's dispensation, choosing what accords with nature and rejecting the contrary.

\ciceroSection{35} There are thus three ends that make no place for moral worth: one belonging to Aristippus or Epicurus, a second to Hieronymus, a third to Carneades; and three in which moral worth is joined with some addition — Polemo's, Callipho's, Diodorus's; and one that is simple and entire, with Zeno as its author, resting wholly on the fitting and the beautiful, that is, on moral worth alone. As for Pyrrho, Aristo, and Erillus, they were long since discarded. The rest held consistently, so that their final ends harmonized with their starting-points: thus for Aristippus the end is pleasure, for Hieronymus absence of pain, for Carneades the enjoyment of the primary natural endowments. But Epicurus, having placed pleasure in the primary commendation of nature, was bound, if he meant the pleasure Aristippus meant, to hold the same final good as Aristippus; or if he meant the pleasure Hieronymus meant, he ought not then to have done what he did — place the pleasure of Aristippus in the primary commendation.

\ciceroSection{36} For his claim that pleasure is judged a good and pain an evil by the senses themselves concedes more to the senses than our laws permit us when we sit as judges in private litigation. We can judge nothing except what falls within our jurisdiction — and for this reason judges are accustomed, when they deliver their verdict, to add in vain the formula ``if this is within my jurisdiction''; for if it was not within their jurisdiction, the verdict is no more given for the omission of that phrase. What do the senses judge? Sweet and bitter, smooth and rough, near and far, at rest and in motion, square and round.

\ciceroSection{37} Reason, then, when brought to bear — reason armed first with the knowledge of divine and human things, which may rightly be called wisdom, and then accompanied by the virtues, which reason would have as mistresses of all things, while you would have them as servants and attendants of the pleasures — reason will pronounce its just verdict, and its first pronouncement will be that pleasure has no place here: not merely that it should not be set alone in the seat of the highest good that we are seeking, but that it should not even be joined to moral worth in any way.

\ciceroSection{38} The same verdict will apply to the absence of pain. Carneades too will be dismissed, and no account of the highest good that shares in pleasure or painlessness, or that is without a share in moral worth, will win approval. Thus two positions will remain for further and repeated consideration. Either reason will determine that nothing is good except what is honorable and nothing evil except what is base, and that everything else either carries no weight at all or so little that it is not to be sought or shunned but merely selected or set aside; or it will prefer that end which is seen to be most richly adorned by moral worth and at the same time enriched by the primary endowments of nature and by the perfection of a whole life. This determination it will make the more cleanly once it has seen whether the dispute between these two positions concerns the facts or merely the words.

\ciceroSection{39} Following the authority of that principle, I shall do the same. I shall, as far as I can, reduce the contentiousness, and I shall think that all the simple positions — those in which there is no admixture of virtue — ought to be excluded from philosophy altogether: first, those of Aristippus and all the Cyrenaics, who had no fear of placing the highest good in that pleasure which stirs the senses with the greatest sweetness, making light of this very absence of pain.

\ciceroSection{40} These men failed to see that, just as the horse is born for running, the ox for ploughing, the hound for tracking, so man — as Aristotle says — is born for two things, understanding and action, as a mortal god; and that conversely they wished to have this divine animal come into being as some sluggish and languid creature of the field, born for grazing and the pleasures of procreation — than which nothing seems to me more absurd.

\ciceroSection{41} So much against Aristippus, who counts that pleasure not merely the highest good but the only one — the pleasure that everyone, using a single term, calls pleasure. Your school sees it differently. But he, as I said, errs badly. For neither the bodily form of man nor the surpassing power of his intelligence indicates that he was born for this one purpose alone — the enjoyment of pleasures. Nor should Hieronymus be given a hearing, whose highest good is the very thing that you sometimes — or rather all too often — say: the absence of pain. For if pain is an evil, freedom from that evil is not sufficient for the good life. That is better left for Ennius to say: \textit{Nimium boni est, cui nihil est mali} — ``He has too much of good who has nothing of evil.'' Let us judge the happy life by the gaining of good, not by the driving away of evil; and let us seek it not in idleness, whether in the enjoyment of pleasure after Aristippus's fashion, or in the absence of pain after this man's, but in doing something and reflecting on something.

\ciceroSection{42} The same arguments can be made against that Carneadean highest good — a position he advanced not so much to defend it as to set it against the Stoics, with whom he was at war. It is, however, of such a character that, if joined to virtue, it seems as though it would carry real authority and could fully supply the happy life, which is the whole question before us. For those who attach to virtue either pleasure — the one thing virtue holds most cheap — or the absence of pain, which, even if it is free from evil, is still not the highest good, are availing themselves of an addition that is not particularly convincing; nor do I understand why they make the addition so sparingly and grudgingly. It is as if they were being charged for what they add to virtue — and first they add the cheapest of commodities, then they add single items one at a time, rather than joining to virtue in a body all those things which nature first approved.

\ciceroSection{43} Because Aristo and Pyrrho regarded all such things as utterly worthless — declaring that there was no difference whatsoever between perfect health and the gravest illness — argument against them was rightly abandoned long ago. In their effort to confine everything within virtue alone, they stripped virtue of all power of selection and gave it neither a source from which it might spring nor a ground on which it might stand; and so in the very act of embracing virtue they destroyed it. Erillus, by referring everything to knowledge, glimpsed a single good, but not the highest good and not one capable of governing life. And so he too was dismissed long ago; after Chrysippus came along, no one seriously argued against him. You two, then, are what remain; for the contest with the Academics is uncertain — they affirm nothing and are content, as if despairing of certain knowledge, to follow whatever appears probable.

\ciceroSection{44} With Epicurus the contest is more demanding, for two reasons: he conflates two kinds of pleasure, and he personally, along with his friends and the many subsequent defenders of his view, has had, I know not how, an influence entirely disproportionate to his authority — the popular crowd is with them. Unless we refute them, all virtue, all decency, all true glory must be surrendered. So, with all other positions cleared away, what remains is not my contest with Torquatus, but virtue's contest with pleasure. And that contest, the sharp and careful Chrysippus did not regard as trivial; he thought the whole difference between the highest goods lay precisely in the comparison of these two. For my own part, I am persuaded that once I have demonstrated that moral worth is something pursued for its own power and for its own sake, your entire position collapses. And so I shall establish briefly, as the occasion demands, what moral worth is, and then address every one of your arguments, Torquatus — unless my memory should fail me.

\ciceroSection{45} By the honorable, then, we understand something of such a kind that, with all advantage stripped away and with no rewards or gains whatever attached, it can rightly be praised for itself alone. How to apprehend what this is — not so much through the definition I have given, though that helps somewhat, as through the common judgment of all mankind and through the pursuits and deeds of the best men, who do a great many things for this one reason alone: because it is fitting, because it is right, because it is honorable, even when they see that no gain will follow. Human beings differ from beasts in many respects, but above all in this one: they have been given by nature the power of reason, a mind keen and vigorous and capable of working through many things at once — a mind that is, so to speak, keen-scented, one that sees the causes of things and their consequences, draws analogies, connects what is separate, links the future to the present, and takes in the whole sweep of a life to come. This same power of reason also made man a being that reaches out toward other men, connatural with them in fellowship of nature, speech, and way of life, so that setting out from the love of those nearest to him he extends his care further and further, entangling himself first in the fellowship of his fellow citizens and then of all mankind; and, as Plato wrote to Archytas, he must remember that he was not born for himself alone, but for his country, for those dear to him — leaving to himself only the smallest portion.


\ciceroSection{46} And since that same nature has implanted in man a desire to see the truth — as is most plainly evident in the fact that when we are free from care we long to know even what passes in the heavens — we are led by these first impulses to love all that is true: that is, the faithful, the candid, the consistent. Falsehood, deceit, imposture we correspondingly abhor, as fraud, perjury, malice, injustice. That same rational faculty holds within itself something vast and magnificent, fitted more to command than to obey; it counts all human things not merely endurable but slight, it is lofty, elevated, fearing nothing, yielding to no one, always unvanquished.

\ciceroSection{47} When these three classes of the honourable have been marked out, a fourth follows, one that belongs to the same beauty and is fitted together from the other three, and in which order and restraint reside. When its likeness was perceived in the beauty and dignity of outward forms, men passed from that recognition to the beauty of deeds and of speech. For from those three virtues I have described, this fourth holds recklessness in check, does not venture to harm anyone by insolent word or act, and is afraid to do or say anything that would appear insufficiently worthy of a man.

\ciceroSection{48} There you have, Torquatus, the form of moral worth set out complete and fully rounded, a form sustained by these four virtues, the very virtues you yourself named. Your Epicurus declares that he has no idea at all what those who measure the highest good by moral worth mean, or what sort of thing they are proposing. For if they refer everything to moral worth and deny that pleasure is contained within it, he says they are making an empty noise with words — he uses exactly these phrases — and neither comprehend nor perceive what meaning is to be attached to this name of moral worth. It is, he says, only what common usage means: moral worth is that which common report and fame hold glorious. Yet even this, he concedes, though it is often sweeter than certain pleasures, is nonetheless sought for the sake of pleasure. Do you see what a vast disagreement this opens?

\ciceroSection{49} A distinguished philosopher — a man who set not only Greece and Italy but all the world of nations in motion — declares that he cannot understand what moral worth is, unless perhaps it is what the crowd praises by repute. I, on the contrary, hold that this is frequently disgraceful; and when it is not disgraceful, even then a thing is not called morally worthy because many people praise it, if that thing is in itself right and praiseworthy; it is called worthy because it is of such a nature that, even if men were ignorant of it or had fallen silent about it, it would still be praiseworthy by its own beauty and splendour. And so Epicurus himself, overcome by nature — which cannot be resisted — says elsewhere (and it is what you too said a little while ago) that one cannot live pleasantly unless one also lives honestly.

\ciceroSection{50} But what does he mean by "honestly" there? The same as pleasantly? If so, the statement resolves to this: one cannot live honestly unless one lives honestly. Or does he mean unless one lives in the good repute of the crowd? Without that repute, then, he denies that a pleasant life is possible? What can be more shameful than to make the life of a wise man depend on the chatter of the foolish? What, then, does he understand by "honest" in this passage? Surely nothing except what can rightly be praised for its own sake. For if it is praised for the sake of pleasure, what sort of praise is that — the kind that can be fetched from a market stall? He is not the kind of man who, when he rates moral worth so highly as to say that without it a pleasant life is impossible, should mean by that worth only what is popular, or should deny that a pleasant life is possible without popular esteem; he cannot mean by "honest" anything other than what is right and praiseworthy by its own force, its own nature, of itself.

\ciceroSection{51} And so, Torquatus, when you were saying that Epicurus cries out that one cannot live pleasantly unless one lives honestly, wisely, and justly, you yourself seemed to me to be exulting. Such was the power that resided in those words, because of the dignity of the things they signify, that you grew more erect, you paused at times, you looked toward us as though bearing witness that moral worth and justice were at last receiving some praise even from Epicurus. How becoming it was to you to use words — words without which, if philosophers would not use them, we should have no need of philosophy at all! It is love of these words, so rarely invoked by Epicurus — wisdom, courage, justice, temperance — that has drawn men of the finest minds to devote themselves to the study of philosophy.

\ciceroSection{52} "The sense of sight," says Plato, "is the keenest of all our senses; yet through it we do not see wisdom." \textit{[Greek: ommatōn to oxytaton]} What ardent loves would wisdom kindle if we could behold her! Why, indeed? Because she is so clever at devising pleasures to best advantage? Why is justice praised? And what is the origin of that proverb, worn smooth by long use: "a man to share the dark with" \textit{[Greek: ho en skotō]}? The saying bears the widest application to every sphere: that in all our dealings we should be moved by the reality of the thing, not by any witness.

\ciceroSection{53} The arguments you were urging are feeble and quite without force: that the wicked are tormented by the testimony of their own conscience, and then by the fear of punishment — whether they are actually suffering it or are constantly in dread of suffering it one day. We ought not to picture that man of good character as a timid or weak-spirited person who tortures himself over everything he has done and goes in fear of all things; we should picture him rather as a shrewd man who refers everything adroitly to self-interest — sharp, crafty, an old hand at the game — one who easily contrives how to deceive in secret, without a witness, without any accomplice.

\ciceroSection{54} Surely you do not suppose I am speaking of Lucius Tubulus? He, while serving as praetor presiding over the court of inquiry into murder, took bribes to pervert justice so openly that the following year the tribune Publius Scaevola brought before the assembly the question whether they wished an inquiry to be conducted into the matter. By that plebiscite the senate decreed that the investigation be entrusted to the consul Gnaeus Caepio. Tubulus left for exile at once and did not dare to answer the charge; the case was plain. What I am asking about, then, is not the straightforward criminal but the cunning criminal — such as Quintus Pompeius was when he repudiated the treaty with Numantia; and not a man who is afraid of everything, but first of all one who pays no heed to the testimony of conscience, which is in any case the easiest thing in the world to suppress. For the man who is described as secretive and close-lipped is so far from giving himself away that he will even contrive to appear distressed when someone else acts dishonestly. What else, after all, is it to be cunning?

\ciceroSection{55} I recall being present when Publius Sextilius Rufus laid the matter before his friends in this form: that he was the heir of Quintus Fadius Gallus, in whose will it was written that Fadius had asked him to ensure that the entire inheritance should pass to his daughter. This Sextilius denied had taken place. He was able to do so with impunity — for who was going to refute him? None of us believed him; it was far more plausible that this man was lying, since he had an interest in lying, than that the testator, having written that he made a request he was bound to make, had not made it. He added that he had sworn to observe the Lex Voconia and could not violate it without perjury, unless his friends judged otherwise. We were present as young men, but many men of the highest standing were there too, and not one of them thought Fadia should receive more than the Lex Voconia allowed her to inherit. Sextilius kept possession of a very large estate; had he followed the advice of those who place the right and honourable above every gain and advantage, he would not have touched a single coin. Do you suppose he was afterwards a man of anxious or troubled mind? Not in the least; on the contrary, he was made wealthy by that inheritance and rejoiced in it accordingly. He set great store by money, not merely when acquired without violating the laws, but even when acquired by the laws' own means. And money, on your account, is certainly to be pursued even at some risk; for it is the producer of many great pleasures.

\ciceroSection{56} Just as those, therefore, who hold that right and honourable conduct is desirable for its own sake must often face danger for the sake of honour and integrity, so those on your side, who measure everything by pleasure, must face danger in order to win great pleasures. If a great matter is at stake, a great inheritance, and money is the parent of very many pleasures, then your Epicurus, if he wishes to follow his own definition of the good, must do what Scipio did when the glory of drawing Hannibal back to Africa was set before him. How great a danger he accepted! For he was referring his whole effort to moral worth, not to pleasure. So your sage, moved by some substantial gain, will fight to the death with hemlock if need be.

\ciceroSection{57} If the crime has been committed in secret, he will rejoice; if he is caught, he will despise every penalty. For he will have schooled himself to despise death, exile, and even pain itself. Pain, which you render intolerable when you threaten it as punishment for the wicked, you render tolerable when you insist that the wise man always has more good than bad. But picture your shrewd wrongdoer as not only shrewd but powerful as well — as Marcus Crassus was, a man who, for all that, was accustomed to make good use of what was his — or as our Pompey is today, who deserves gratitude for acting rightly, since he could have acted unjustly with impunity as much as he pleased. How many wrongs can be done that no one is in a position to reproach! Suppose a dying friend asks you to restore his inheritance to his daughter, and has written nothing to this effect — as Fadius did write — and told no one: what will you do?

\ciceroSection{58} You, certainly, would restore it. Epicurus himself would perhaps restore it — as Sextus Peducaeus, son of Sextus, did: a man who left behind in his son a living image of his own humanity and integrity, himself as learned as he was the best and most just of men. He knew that no one had heard of any request made to him by Gaius Plotius, a distinguished Roman knight from Nursia; yet he went of his own accord to the woman, told her of her husband's commission — which she had no expectation of hearing — and restored the inheritance. But what I am asking you is this: since you would certainly have done the same, do you not recognize that the force of nature is all the greater for the fact that even you yourselves, who refer everything to your own advantage and — as you put it — to pleasure, yet perform acts from which it is plain that you are following not pleasure but duty, and that upright nature prevails over your perverse theory?

\ciceroSection{59} "Suppose," says Carneades, "you know that a viper is hidden somewhere, and that a man who does not know this is about to sit down upon it, and his death would be to your profit: you will act dishonestly if you give no warning — yet with impunity; for who can prove that you knew?" But I have said enough on this. It is transparent that unless equity, good faith, and justice proceed from nature — and if all these things are referred back to utility alone — no man of genuine goodness can be found. On these matters there is a sufficiently full treatment in my work \textit{On the Commonwealth}, given in the words of Laelius.

\ciceroSection{60} Apply the same argument to self-restraint, or temperance, which is the governance of the appetites in obedience to reason. Is a man who yields to lust without a witness taking adequate account of decency? Or is there something that is in itself shameful, even when no disgrace accompanies it? And again: do brave men enter battle by casting up a ledger of pleasures, and shed their blood for their country — or are they driven by a certain ardour and impetus of the spirit? Come, Torquatus — which do you suppose that famous Imperiosus, had he heard both our accounts of him, would have preferred to listen to: your account of him, or mine? Mine says he did nothing for his own sake and everything for the commonwealth; yours says he did nothing except for his own. And if you were to make this even more explicit and say openly that he did everything for the sake of pleasure, how in the world do you suppose he would have taken that?



\ciceroSection{61} Very well — grant that Torquatus acted as he did for his own advantage, if you insist (I prefer that word to ``pleasures,'' especially for so great a man). But did his colleague Publius Decius, the first in that family to hold the consulship, think of his own pleasures when he had devoted himself to death and was spurring his horse headlong into the thick of the Latin line? Where could he have seized any pleasure, or when? He knew he must die in the instant, and sought that death with a fiercer passion than Epicurus thinks any pleasure worth pursuing. If his deed had not been rightly praised, his son would not have imitated it in his own fourth consulship; nor in turn would the grandson born of him, while waging war with Pyrrhus as consul, have fallen in battle and offered himself as the third victim of that family to the commonwealth.

\ciceroSection{62} I will restrain myself from examples. The Greeks can manage with a modest supply: Leonidas, Epaminondas, three or four others perhaps. If I begin to collect our own men, the day will fail me before I exhaust them — and I shall accomplish, to be sure, what the project demands: that pleasure be delivered bound and constrained to virtue. But, as Aulus Varius — who had the reputation of a rather harsh judge — used to say to a fellow juror when witnesses had been given and more were still being called forward, ``Either this testimony is enough, or I do not know what enough means'' — so I say: enough witnesses have been produced by me. What then? Was it pleasure that moved you yourself — most worthy of your ancestors — when, still a young man, you snatched the consulship from Publius Sulla? When you brought news of this to your father, that bravest of men — what a consul, what a citizen he was always, and above all after his consulship! It was under his lead indeed that I myself acted as I did, choosing the good of all rather than my own.

\ciceroSection{63} But how fine your argument seemed, when you placed on one side a man heaped up with pleasures great and many, with no pain present or in prospect, and on the other a man racked by the severest torments throughout his body with no pleasure attached or anticipated — and then asked who could be more wretched than the second or more blessed than the first; and from this you concluded that pain is the highest evil and pleasure the highest good. There was a man, Lucius Thorius Balbus of Lanuvium, whom you cannot remember. He lived in such a way that no pleasure however refined could be found in which he did not abound. He was both eager for pleasures and a connoisseur of them and well supplied with them; so far from being superstitious that he held in contempt the many sacrifices and temples of his own native town; so far from fearing death that he was killed in battle for the commonwealth.

\ciceroSection{64} He did not measure his appetites by Epicurus's classification but by his own satiety. He maintained his health, however: he used such exercise as would bring him to dinner both thirsty and hungry; he chose food that was at once the most agreeable and the easiest to digest; wine, too, for pleasure but in quantity short of harm. He employed all those other accessories without which Epicurus declares he cannot conceive what good even is. Every pain was absent; and if any had been present, he would not have borne it softly, yet would have turned to doctors rather than to philosophers. His complexion was excellent, his health unimpaired, his standing in society great, and his life in short filled with a variety of every pleasure.

\ciceroSection{65} You call this man happy — your own system compels you to. But I do not venture to name the man I set above him; virtue itself will speak for me, and will not hesitate to set Marcus Regulus above your happy voluptuary. Regulus returned of his own will to Carthage, compelled by no force but by the pledge he had given to the enemy; and at the very moment when he was being tormented by sleeplessness and starvation, virtue cries aloud that he was more blessed than Thorius drinking among his roses. He had waged great wars, had twice been consul, had celebrated a triumph; and yet he did not count those earlier glories as great or as illustrious as that final ordeal he had embraced for the sake of honor and constancy — an ordeal that to us, hearing of it, seems pitiable, but to him who endured it was a matter of \textit{voluptas}. For it is not by merriment and wantonness, not by laughter and jest — companions of levity — but often amid sorrow, by firmness and constancy, that men are blessed.

\ciceroSection{66} Lucretia, violated by force by the king's son, called her fellow citizens to witness and killed herself. That outrage, with Brutus for its leader and instigator, was the cause of liberty for the state; and in memory of that woman, her husband and her father were made consuls in the first year of the Republic. Lucius Verginius, a man of modest means and one among many, sixty years after liberty was won, killed his own virgin daughter with his own hand rather than surrender her to the lust of Appius Claudius, who then held supreme power.

\ciceroSection{67} Either these acts must be condemned by you, Torquatus, or the advocacy of pleasure must be abandoned. And what advocacy, what case, is this that pleasure can plead — one that can produce neither witnesses nor admirers from among the great men of history? For just as I summon from the records of the annals those witnesses whose whole lives were consumed in glorious toil, men who could not endure even to hear the word ``pleasure'' — so in your school's discussions history falls silent. I have never heard Lycurgus, Solon, Miltiades, Themistocles, Epaminondas named in the school of Epicurus, men who are on every other philosopher's lips. As it is, now that we too have begun to handle this subject, our friend Atticus will supply us from his treasury — what men, and how great! Is it not better to say something about them than to fill so many volumes on the subject of Themista?

\ciceroSection{68} Let those things be Greek matters — though from them we have received both philosophy and every liberal discipline; yet there is still something that is permitted to them but not to us. Stoics and Peripatetics disagree: the one school denies that anything is good except what is honorable, the other grants the very greatest weight — far the greatest — to honor, but holds nevertheless that there are certain goods in the body and in things external. The contest is honorable and the debate is splendid, for the whole dispute is about the dignity of virtue. But when you dispute with your own people, much must be heard about pleasures of the most degraded kind, which Epicurus introduces very frequently.

\ciceroSection{69} You cannot, then, defend these positions, Torquatus — trust me in this — if you look carefully into yourself, your thoughts, and your pursuits. You will blush, I tell you, for that picture which Cleanthes used to paint in words with remarkable skill. He would bid his audience imagine a painting in which Pleasure sat enthroned in full regalia, dressed in the finest attire, with the Virtues standing about her as handmaids, charged with no other duty and regarding nothing else as their proper office except to minister to Pleasure and merely to whisper in her ear — if indeed a painting could make this intelligible — that she should take care to do nothing rash that might give offense to men's minds, or anything from which some pain might arise. ``For it is we, the Virtues, who were born for this: to serve you; we have no other business.''

\ciceroSection{70} But Epicurus — for this is your school's beacon of light — declares that no one who does not live honorably can live pleasantly. As if I cared what he asserts or denies. What I ask is what is consistent for a man who places the highest good in pleasure to say. What reason can you give why Thorius, why Gaius Postumius, why their master in all this, Orata, did not live most pleasantly? Epicurus himself says, as I mentioned earlier, that the life of the dissolute is not to be faulted — provided they are not complete fools, that is, provided they neither crave nor fear. When he promises a remedy for both of these conditions, he is promising license for debauchery. Strip out craving and fear, and he declares he finds nothing in the life of the profligate to censure.

\ciceroSection{71} You cannot, therefore, while directing everything by pleasure, either uphold or retain virtue. For neither should a man be counted good and just who refrains from injustice only to avoid suffering some harm himself. You know the verse, I imagine: \textit{No man is pious who\ldots\ piety\ldots} — take care to think nothing truer than that. For a man is not just so long as he acts out of fear; and assuredly, if he ceases to fear, he will cease to be just: he will not fear, moreover, if he can either conceal what he does or make good his position by the weight of resources — and he will certainly prefer to have the reputation of a good man without being one, rather than to be one without being thought so. Thus you hand us, in place of genuine and assured justice, a counterfeit of justice; and you drive us, in a manner, to disregard our own settled conscience and go hunting for the mistaken opinion of others.

\ciceroSection{72} The same can be said of the other virtues, all of which you ground in pleasure as though in standing water. What then? Can we call that same Torquatus brave? — for I delight in your family and its name, even if I cannot, as you say, corrupt you; I delight in them, I say. And, by Hercules, Aulus Torquatus, that excellent man and warmest of my friends, comes before my eyes — how devoted his enthusiasm was to me, how conspicuous his service in those times known to everyone, both of you must know well. These things would bring me no gratitude — I who wish to be and to seem grateful — unless I saw clearly that he was my friend for my sake and not his own; unless you mean by ``his own'' what is the interest of all men to act rightly. If that is what you mean, we have won. That is precisely what we assert and press: that the reward of doing one's duty is the duty itself.

\ciceroSection{73} But that man of yours refuses this and demands pleasure as the payment for everything he does. To return to the case in hand: if Torquatus fought the Gaul in single combat by the Anio at the Gaul's own challenge, and dressed himself from his spoils in the torque and the cognomen — if he did all this for no reason other than that such deeds seemed to him worthy of a man — I do not count him brave. Again, if modesty, discretion, chastity, in a word temperance, are kept in check only by fear of punishment or disgrace and not by their own sanctity, then what act of adultery, what defilement, what lust will not break out and hurl itself forward, given the prospect of concealment, or impunity, or license?

\ciceroSection{74} And consider this, Torquatus: with your name, your ability, your distinction — what you do, what you contemplate, what you strive for — to what end do you refer it all? For what purpose do you wish to accomplish what you attempt? What, in the end, do you judge best in life? Do you not dare say it in public? What would you deserve to hear, when you enter on your magistracy and mount the platform to address the assembly — for you will have to announce by edict what rules you intend to observe in the administration of justice, and perhaps too you will speak, following ancestral custom, something about your forebears and yourself — what, I ask, would you deserve to hear if you declared that in that magistracy you would do everything for the sake of pleasure, and that in your whole life you had done nothing except for the sake of pleasure? ``Do you think me,'' you say, ``so foolish as to speak that way before an audience of laymen?'' Then say the same thing in court, or — if you shrink from a crowd — say it in the Senate. You never will. Why, unless because the speech is shameful? Do you then think that Triarius and I deserve to have shameful things said in our presence?

\ciceroSection{75} But let it pass: the word ``pleasure'' itself, you will say, lacks dignity, and perhaps we do not understand it. This, indeed, is what you repeat again and again — that we do not understand what pleasure you are speaking of. A difficult and obscure matter, evidently! When you speak of indivisible bodies and the spaces between worlds, which neither exist nor can exist, we understand; but pleasure, which every sparrow knows, cannot be understood by us? If I demonstrate that I know not only what pleasure is — it is an agreeable motion felt in the senses — but also what you want it to be, then at one moment you want the thing itself that I just described, and you give it the name pleasure in the sense that it involves movement and produces some variety of sensation; but at another moment you want a certain supreme pleasure, to which nothing can be added — the pleasure that is present when all pain is absent, which you call ``stable.''

\ciceroSection{76} Very well, call it pleasure. But try saying in any public assembly that everything you do is to avoid pain. Or if you think even that is not said with sufficient fullness or dignity, say that in your magistracy and throughout your whole life you will do everything for your own advantage, nothing except what is expedient, nothing, in the end, except for your own sake. What do you imagine the roar of the assembly would be, or what prospect for the consulship that lies so ready in your grasp? Is this then the principle you actually follow — one you use privately, with yourself and your household, yet dare not profess or put before the public? But the doctrines of the Peripatetics and the Stoics are always on your lips in the courts and in the Senate: duty, fairness, dignity, trust, rectitude, honour, what befits empire, what befits the Roman people, to face every danger for the commonwealth, to die for the fatherland. When you speak these words, we plain men are struck dumb with admiration, while you, of course, are laughing to yourself.

\ciceroSection{77} For among words as magnificent and splendid as those, pleasure finds no place at all — not merely that kinetic pleasure which all men, in city and countryside alike, all who speak Latin, call by the name pleasure, but not even that stable condition which no one but you people calls pleasure. Take care, then, lest you have no right to use our words while holding your own opinions. If you were putting on a certain expression or walk to appear more serious, you would not be your true self; are you then to put on a mask of words and say what you do not believe? Or are you, like a man who keeps one set of clothes for home and another for the Forum, to hold one opinion in private and another in public — to display one thing on your face while truth is hidden within? Ask yourself, please, whether that is right. It seems to me that those opinions are the true ones which are honourable, praiseworthy, and glorious, which can be professed in the Senate, before the people, in every gathering and assembly — so that one need not be ashamed to believe what one would be ashamed to say. And where, in any case, can friendship find a place, or who can be anyone's friend, if a man does not love that friend for the friend's own sake?

\ciceroSection{78} What is it to love — the word from which friendship itself draws its name — if not to wish a person the greatest goods, even when nothing of those goods flows back to oneself? ``It benefits me,'' he says, ``to be disposed in that way.'' Perhaps to \textit{seem} so disposed — but you cannot \textit{be} disposed so unless you actually are. And how will you be, unless love itself has taken hold of you? That is not something normally brought about by a calculated reckoning of advantage; it arises of itself and is born of its own accord. But if, you say, I am following what is useful, then friendship will last only as long as utility follows, and if utility is what constitutes friendship, the same utility will destroy it.

\ciceroSection{79} But what will you do in the end when utility, as so often happens, has deserted friendship? Will you abandon it? What sort of friendship is that? Will you keep it? How does that square with your principle? You can see what you have committed yourself to by holding that friendship is to be sought for the sake of utility. ``I shall not abandon a friend lest I fall into hatred.'' But why in the first place is that action worthy of hatred, unless because it is base? And if your reason for not deserting a friend is merely to avoid some personal disadvantage, you will nonetheless wish him dead, so as not to be tied down without profit. What then if it is not merely no advantage to you, but your private fortune must be spent, labours undertaken, your life put at risk? Will you not even then look to yourself and reflect that every man is born for himself and for his own pleasures? Will you stand surety to a tyrant for your friend's life unto death, as that Pythagorean did for the Sicilian tyrant? Or, if you were Pylades, will you declare yourself Orestes so that you may die in your friend's place? Or, if you were Orestes, will you contradict Pylades, name yourself, and if that is not accepted, at least beg that you not both be put to death together?

\ciceroSection{80} You would do all these things, Torquatus; for I hold nothing worthy of great praise that I believe you would pass by for fear of death or pain. But the question is not what accords with your nature; it is what accords with your doctrine. The principle you defend, the precepts you have learned and approved, overturn friendship from its foundations — however lavishly Epicurus, as he does, raises it to the sky in praise. ``But he himself kept friends.'' Who, I ask, denies that he was a good man, agreeable, and humane? In these discussions it is his intellect, not his character, that is under scrutiny. Let that sort of perversity remain with the Greeks in their fickleness — those who vilify with abuse the people from whom they differ on a point of truth. But however warmly he may have tended his friendships, still, if these doctrines are true — for I assert nothing absolutely — he was not quite acute enough. ``But he commended himself to many.''

\ciceroSection{81} And perhaps with justice; yet the verdict of the crowd is not the weightiest testimony. In every art, every pursuit, every branch of knowledge, and in virtue itself, the best of anything is the rarest. And yet for my part — because Epicurus himself was a good man, and there are and have been many Epicureans who are faithful in their friendships, consistent in their whole manner of life, grave in character, and who regulate their conduct by duty rather than by pleasure — this seems to me to be the greater power of honour and the lesser power of pleasure at work. For some men live in such a way that their life gives the lie to their doctrine. And whereas other men are generally reckoned to speak better than they act, these seem to me to act better than they speak.

\ciceroSection{82} But none of this matters much to the point. Let us look at what you said about friendship. Of the arguments you made, one I recognized as coming from Epicurus himself: that friendship cannot be separated from pleasure, and that therefore it is to be cultivated, because without it one cannot live in security and without fear, and therefore cannot live pleasantly either. That answer is sufficient. But you also adduced something more humane from the more recent Epicureans — something which, as far as I know, was never said by Epicurus himself — namely, that friendship is sought at first for the sake of utility, but that once intimacy has grown, a friend comes to be loved for his own sake, independent of any hope of pleasure. This can be criticized in many ways; but I will accept what they grant. For they are saying that sometimes a right action can be done with no expectation of pleasure sought or awaited — and that is enough for me, though it is not enough for them.

\ciceroSection{83} You also put forward the view of some who say that the wise make a compact among themselves that they will be disposed toward their friends exactly as they are toward themselves — that this is both possible and has often come about, and that it is supremely relevant to the attainment of pleasures. If they can forge that compact, let them forge this one too: that they love fairness, moderation, and all the virtues for their own sake, without any reward. But if we cultivate friendships for their fruits, their gains, and their utilities, if there will be no affection of the heart that makes friendship desirable in and of itself, by its own force, from itself and for itself, then surely there is no question that we would place farms and islands above friends.

\ciceroSection{84} You are free at this point to rehearse again the words in which Epicurus speaks in praise of friendship — and they are excellent words. I am not asking what he says, but what he can say that is consistent with his own principle and doctrine. Friendship, on his account, is sought for utility's sake. Can you seriously think this Triarius here could be more useful to you than your own granaries at Puteoli? Gather all your usual arguments: ``The protection of friends.'' You have enough protection in yourself, enough in the laws, enough in ordinary friendships. You will not, after all, be exposed to contempt. And you will easily avoid hatred and envy — Epicurus himself gives precepts for that. Besides, with revenues large enough to support generosity, you can, even without this Pylades-and-Orestes intensity of friendship, secure and fortify yourself admirably with the goodwill of many. ``But with whom will you share your jokes and your serious concerns, as the saying goes, with whom your secrets, your most private thoughts?''

\ciceroSection{85} With yourself, best of all — and then with a friend of ordinary attachment. But grant that those needs are real and pressing: what have they to do with the utility of such enormous wealth? You see, then, that if you measure friendship by its own warmth of affection, nothing is more excellent; but if you measure it by profit, the most intimate bonds of friendship are outstripped by the rent from productive estates. I must be loved for myself, not for my possessions, if we are to be true friends. But we are dwelling too long on matters that are perfectly plain. For once we have established and concluded that virtue and friendship find no place anywhere if everything is referred to pleasure, there is not much more that needs to be said. And yet, so that no part of your argument may seem to have gone without reply, I will add a few more words on the rest of what you said.

\ciceroSection{86} Since, then, the whole sum of philosophy is referred to living happily, and it is in pursuit of this one thing that men have given themselves to this study, while some place the happy life in one thing and some in another — you in pleasure — and correspondingly place all misery in pain, let us first examine what your happy life amounts to. And you will grant me this, I think: if there is any such thing as being happy, it must reside entirely within the power of the wise man. For if a happy life can be lost, it cannot be happy. Who can feel confident that something fragile and perishable will remain stable and secure for him? And whoever lacks confidence in the permanence of his goods must of necessity fear that he will at some point lose them and become miserable. But no one can be happy in fear of losing the greatest things. Therefore no one can be happy at all.

\ciceroSection{87} For a happy life is not generally said to reside in some portion of time, but in the whole extent of it; nor is a life called complete in any sense unless it is finished and fulfilled; nor can anyone be happy at one time and miserable at another — for a man who considers that he might become miserable will not be happy. Once a happy life has been taken up, it endures as surely as that wisdom which is itself the author of the happy life, and it does not wait for the final hour of a man's years, as Herodotus writes was the counsel Solon gave to Croesus. ``But,'' you will say — and as you yourself observed — Epicurus denies that length of time adds anything to living happily, and holds that no less pleasure is received in a brief span than if it were everlasting.

\ciceroSection{88} These statements are thoroughly inconsistent. When he places the highest good in pleasure, he denies that a longer span of life produces greater pleasure than a finite and moderate one. A man who places all good in virtue can say that the happy life is made complete by the perfection of virtue, for he denies that a day adds anything to the highest good. But a man who holds that life is made happy by pleasure — how will he be consistent with himself if he denies that pleasure grows with length of time? By the same token, he should say the same of pain. Is it not the case that the most protracted pain is the most wretched, and that duration makes pleasure more desirable? Why, then, does Epicurus forever call god happy and eternal? Strip away the eternity and Jupiter is no more blessed than Epicurus; both enjoy the highest good, that is, pleasure. ``But Epicurus also has pain.'' He counts it as nothing, for he says that if he were being burned alive he would declare, ``How pleasant this is!''

\ciceroSection{89} In what respect, then, is he surpassed by god, if he is not surpassed in eternity? And what good is there in eternity beyond the highest pleasure, and that pleasure everlasting? What is the point, then, of speaking grandly unless you speak consistently? ``The happy life resides in bodily pleasure'' — add ``and in pleasure of the mind too,'' if you like, so long as that mental pleasure is itself, as you hold, derived from the body. Very well: who can guarantee that pleasure in perpetuity to the wise man? For the things by which pleasures are produced are not within the wise man's control. The happy life does not reside in wisdom itself, but in those things which wisdom secures with a view to pleasure. But the whole of that is external; and what is external is at the mercy of chance. So it comes about that fortune is the mistress of the happy life — the very fortune which Epicurus says makes only slight incursions on the wise man.

\ciceroSection{90} ``Come,'' you will say, ``these are trifling objections. Nature itself makes the wise man rich, and Epicurus has taught that nature's wealth is easily obtained.'' This is well said, and I do not dispute it; but these very claims are at war with one another. For he denies that a smaller pleasure is taken from the plainest fare — the most ordinary food and drink — than from the most exquisite table. If he were denying that there is any difference for living happily in what kind of food one uses, I would grant it and even praise it; he would be saying something true. I hear Socrates, who takes no account of pleasure at any point, saying that hunger is the best seasoning for food, and thirst for drink. But a man who refers everything to pleasure and lives like Gallonius yet talks like Piso the Frugal — him I do not heed, nor do I think he is saying what he really believes.

\ciceroSection{91} He declared that natural wealth was easy to come by, on the ground that nature is content with little. Certainly — if only you did not set pleasure at so high a price. The pleasure one derives from the cheapest things, he says, is no less than that derived from the most costly. This is not merely to lack good sense but to lack even a palate. Those who despise pleasure itself may fairly say they prefer a herring to a sturgeon; but for the man whose highest good resides in pleasure, all things must be judged by sensation and not by reason, and those things must be called best that are the most agreeable.

\ciceroSection{92} Very well; let him obtain the highest pleasures not only cheaply but, for all I care, at no cost at all, if he can manage it. Let the pleasure in that cress which Xenophon records the Persians habitually ate be no less than in the tables of Syracuse that Plato so sharply condemns. Let the procurement of pleasure be, I say, as easy as you wish. But what shall we say about pain? Its torments are so severe that a happy life, if pain is truly the highest evil, cannot exist within them. For Metrodorus himself — virtually a second Epicurus — describes the happy man in words roughly to this effect: \textit{when the body is in good condition and the prospect of its remaining so is assured.} But can that prospect be assured for anyone — I will not say for a year, but for the evening? If pain, which is the highest evil, is always to be feared even when it is not present — since at any moment it may arrive — how can the highest evil's fear take up residence in a happy life?

\ciceroSection{93} Epicurus, he says, has handed down a method for disregarding pain. That itself is already absurd: that the greatest evil should be disregarded. But what, precisely, is this method? \textit{The most intense pain, he says, is brief.} First, what do you mean by brief? And then, what pain do you call most intense? What of it — can the most extreme pain not persist for several days? See that it does not persist for months! Unless you mean the kind that kills the instant it seizes its victim. Who fears that kind of pain? I would rather you had relieved the pain by which I watched a man of the finest character and the most humane disposition, Gnaeus Octavius, son of Marcus, a friend of mine, laid low — not once, and not briefly, but frequently and for an extended time without respite. What torments he endured, immortal gods, when all his limbs seemed to be on fire! And yet he did not appear wretched, because that was not the highest evil, but only burdened and hard-pressed; whereas he would have been wretched had he been swimming in pleasures while living a degraded and vicious life.

\ciceroSection{94} Your claim that intense pain is brief and prolonged pain mild — I do not grasp what it amounts to. For I observe pain that is both intense and of considerable duration, and there is a truer way of bearing it, a way you cannot employ, because you do not love honor as a good in its own right. Courage has its own precepts, its own rules, which forbid a man to become unmanned in the face of pain. For this reason it should be thought a shameful thing — not to feel pain, for that is sometimes unavoidable, but to defile with Philoctetes' cries that Lemnian rock, which \textit{resounds with wailing, lamentation, groaning, and roaring, giving back his woeful voice.} Let Epicurus whisper his consolations — if he can — to the man from whose \textit{veins, soaked in viper's venom, the poison rouses hideous torment.} This is Epicurus's reply: Philoctetes, hush — the pain is brief. But the man has lain ten years already in his cave. Then if it is long, it is mild; for it grants intervals and lets up.

\ciceroSection{95} First, it does not often grant them; and second, what kind of respite is it when the memory of recent pain is fresh and the fear of the approaching pain already overhead torments the mind? \textit{Let him die}, someone says. Perhaps that would be best; but then what of the principle that \textit{pleasure always outweighs}? If that is true, consider whether you commit a wrong in counseling death. Let us rather say this: it is shameful, it is unworthy of a man, to be broken by pain, crushed, overwhelmed. Your maxims — \textit{if it is severe, it is brief; if it is long, it is mild} — are slogans, nothing more. Pain is ordinarily assuaged by the remedies of courage, greatness of soul, endurance, and fortitude.

\ciceroSection{96} Hear — so that I need not go far afield — what Epicurus says as he dies, so that you may see how his actions conflict with his words. \textit{Epicurus to Hermarchus, greetings. In the course of what is at once a blissful and my final day, I write these words.} The afflictions of bladder and bowel were present in such severity that nothing could exceed their violence. Wretched man! If pain is the highest evil, that is the only word for it. But let us hear him further. \textit{Yet against all these,} he says, \textit{was set the joy in my soul that I drew from the recollection of my arguments and discoveries. But do you, as befits the goodwill toward me and toward philosophy which you have maintained since your youth, take care of Metrodorus's children.}

\ciceroSection{97} I will not now set the death of Epaminondas or of Leonidas above his. The former, having defeated the Lacedaemonians at Mantinea, saw that a grievous wound was killing him; as soon as he could see clearly, he asked whether his shield was safe. When his weeping men told him it was, he asked whether the enemy had been routed. When he had heard that too as he desired, he ordered the spear by which he was transfixed to be drawn out. So he died in joy and victory, with much blood flowing. Leonidas, the Lacedaemonian king, faced with the choice of shameful flight or glorious death, posted himself and the three hundred he had led out from Sparta against the enemy at Thermopylae. Those are deaths worthy of commanders; philosophers as a rule die in their own beds. Yet it matters how they die. He counts himself blessed as he dies. High praise indeed. \textit{Against my most intense pains,} he says, \textit{was set the joy of my soul.}

\ciceroSection{98} I hear the voice of a philosopher, Epicurus; but you have forgotten what you are bound to say. For in the first place, if the things whose recollection you claim gives you joy are true — that is, if your writings and doctrines are sound — then you cannot feel joy at all. You have nothing left to refer to the body; and you have always maintained that no one feels joy or pain except on account of the body. \textit{I take joy}, he says, \textit{in what is past.} In what past things? If they are things that had to do with the body, then I observe that you are setting your doctrines against these pains as a counterweight — not the memory of bodily pleasures. But if they are things that had to do with the mind, your claim that no joy of the mind exists which is not referred to the body is false. And why, further, do you commend to Hermarchus the children of Metrodorus? What is there, in that noble act of devotion, in that great faithfulness of yours — for I judge it to be great — that you refer to the body?

\ciceroSection{99} Turn this way and that, Torquatus — you will find nothing in that celebrated letter consistent or coherent with Epicurus's own doctrines. Thus he is refuted by himself, and his writings are convicted by his own integrity and conduct. For that commendation of children, the remembrance and warmth of friendship, the keeping of the highest obligations in the final breath — all of this shows that probity is inborn in human beings as a thing given freely, not solicited by pleasures nor summoned by the wages of reward. What greater testimony do we need that the things that are honorable and right are desirable for their own sake, when we behold such solemn duties performed by a dying man?

\ciceroSection{100} But while I judge that letter worthy of praise — and I have just rendered it almost word for word — it is in no way consistent with the whole of his philosophy; and I judge his will to be at odds not only with the dignity proper to a philosopher, but with his own stated convictions. For he wrote, often and at length, and also briefly and precisely in the book I have just mentioned, that death is nothing to us. His argument was that whatever has been dissolved is without sensation, and whatever is without sensation is nothing to us at all. The thesis could have been set out more elegantly and more neatly. As it stands — \textit{whatever has been dissolved is without sensation} — the statement does not say clearly enough what is meant by dissolution. Yet I understand what he intends.

\ciceroSection{101} I ask, then, why it is that when, at dissolution — that is, at death — all sensation is extinguished, and when absolutely nothing remains that pertains to us, he takes such careful and precise precautions in his will: directing that Amynomachus and Timocrates, his heirs, shall provide, with the approval of Hermarchus, sufficient funds to celebrate his birthday every year in the month of Gamelion, and that every month on the twentieth day of the moon they shall likewise provide for a feast of those who have philosophized alongside him, so that his own memory and the memory of Metrodorus may be honored.

\ciceroSection{102} I cannot say that these provisions are unbecoming a man of charm and warmth; but they are in no way becoming a philosopher — least of all a natural philosopher, which he claims to be — who supposes that any person has a birthday. Why? Can the same day recur that has come once? Surely it cannot. Can one of the same kind recur? Not even that, unless many thousands of years elapse so that all the heavenly bodies return at the same moment to their original positions. There is therefore no such thing as anyone's birthday. So it is observed! Yes, and I certainly did not know that! But even if it exists — will it be observed after death? And will the man who proclaimed to us as if by oracle that nothing after death pertains to us be making provision for it in his will? These were not the thoughts of a man who had traversed in his mind unnumbered worlds and boundless regions with no shore and no boundary. Would Democritus have done anything of the sort? To pass over others, I cite him: the one man Epicurus himself followed.

\ciceroSection{103} And if a day was to be marked, should it not rather be the day on which he was born, but the day on which he became wise? You will say he could not have become wise had he not first been born. By that reasoning, not even had his grandmother not been born. The whole business, Torquatus, is unworthy of educated men — this wish to have the memory of one's name celebrated at feasts after death. I will not say how you conduct those days or what ridicule you incur from people of wit and taste — there is no need for a quarrel — I will say only this: it would have been more in keeping with your character to observe Epicurus's birthday of your own accord than to have him direct in his will that it be observed.

\ciceroSection{104} But to return to our subject — we had been speaking about pain when we were drawn off to that letter — the whole argument may now be closed as follows: the man who is in the grip of the highest evil is not happy at the time he is in its grip; the wise man, however, is always happy and is at some times in pain; pain, therefore, is not the highest evil. Now what is to be made of the claim that the wise man's past goods do not slip away from him and that he ought not to dwell on past misfortunes? First, is memory of what we choose to remember within our power? When someone promised Themistocles the art of memory, or whatever man it was, Themistocles replied: \textit{I would rather have the art of forgetting.} For, he said, I remember even what I wish to forget, and cannot forget what I wish to forget.

\ciceroSection{105} A man of great genius — yet the truth of the matter is that it is the act of a philosopher too imperious to forbid remembering. Take care lest these commands of yours be like the Manlian orders, or more severe still, if you command what I have no power to perform. And what if the memory of past misfortunes is itself pleasant? So that certain proverbs prove truer than your doctrines. For the common saying runs: \textit{labors once past are sweet to recall,} and Euripides put it well — I shall render it in Latin if I can; you all know the Greek verse: \textit{Sweet is the memory of labors past.} \textit{[Greek: hēdeia mnēmē ponōn parelthontōn]} But let us return to past goods. If such goods were of the kind that Gaius Marius could call upon — so that, driven into exile, destitute, sunk in a marsh, he might lighten his suffering by the recollection of his trophies — I would listen and approve unreservedly. For the blessed life of the wise man can neither be completed nor brought to its conclusion if each good decision and each good action is buried under the oblivion of the man who performed them.


\ciceroSection{106}
But it is the recollection of pleasures already enjoyed, you say, that makes a life happy — and specifically pleasures enjoyed through the body. For if there are any others, then it is simply false that all pleasures of the mind derive from their partnership with the body. Now if the pleasure of the body delights us even in retrospect, I cannot understand why Aristotle ridicules so sharply the epitaph of Sardanapallus — that king of Syria who boasted in it that he had carried away with him into death every gratification of his lusts. For what, Aristotle asks, could he have gone on feeling even while alive for any longer than the moment of his enjoyment — and how then could it have persisted after he was dead? Bodily pleasure, then, flows away; every instance of it takes flight almost as it arrives, and more often it leaves behind a cause for regret than for fond recollection. Hence Africanus is far happier, addressing his country in those famous words: \textit{Cease, Rome, your enemies to fear} — and the rest, spoken so magnificently: \textit{For my labours have won you monuments of glory.} He rejoices in past \textit{labours}; you command men to rejoice in past pleasures. He calls himself back to deeds from which he never drew any return to his body; you are absorbed entirely in the body. And further: how is the claim you advance even coherent — that all pleasures and all pains of the mind relate to the pleasures and pains of the body?

\ciceroSection{107}
Is there nothing, then, that ever gives you pleasure in itself alone? I see who I am speaking to — you, Torquatus: does nothing give you pleasure purely on its own account? I pass over dignity, moral worth, the very beauty of the virtues, which have been discussed already. Let me put forward things of less weight: a poem, a speech in prose, whether in writing or in reading; the history of all great deeds, of all the regions of the earth, which you are forever pursuing; a work of sculpture, a painting, a pleasant landscape; the games; the hunt; the villa of Lucullus — for if I named yours, you would have a way out and would say it relates to the body — but do you refer the things I have listed back to the body? Is there nothing that gives you pleasure for its own sake? Either you will be most obstinate if you persist in referring what I named to the body, or you will have abandoned the whole of Epicurus's account of pleasure if you deny it. And as for your argument that the pleasures and pains of the mind are greater than those of the body, because the mind participates in all three tenses of time while the body feels only the present — how can it be demonstrated that someone who is glad on my account rejoices more than I do myself?

\ciceroSection{108}
The pleasure of the mind arises because of the pleasure of the body, and the pleasure of the mind is greater than that of the body. It follows that the one who offers congratulations is more joyful than the one being congratulated. But while you are trying to show that the wise man is happy — since he receives the greatest pleasures of the mind, greater in every dimension than those of the body — you do not see what difficulty stands in your way. For he will also receive pains of the mind greater in every dimension than those of the body. And so the man whom you insist must always be happy will necessarily be wretched on occasion — and you will never achieve that result, not for so long as you refer everything to pleasure and pain.

\ciceroSection{109}
Therefore, Torquatus, some other highest good for man must be found. Let us concede pleasure to the animals — the same animals whose testimony you habitually invoke on the question of the highest good. But what of the fact that even animals do many things, each guided by its own nature, partly with a kind of tenderness and partly with effort — in giving birth, in rearing their young — so that it is perfectly apparent that something quite different from pleasure is what they have in view? Others again take delight in running and in wide-ranging movement; others imitate in their flocking something resembling the fellowship of a civic community.

\ciceroSection{110}
We observe in certain kinds of birds unmistakable signs of dutiful affection, of recognition, of memory; and in many creatures we observe longing and attachment. Can it be, then, that in animals there will be found certain counterparts of human virtues, distinct from pleasure, while in human beings themselves there will be no virtue except for the sake of pleasure? And shall we say that man, who surpasses all other living creatures by so wide a margin, has been given by nature no particular prerogative of his own?

\ciceroSection{111}
If indeed everything comes down to pleasure, we are far and away surpassed by the animals, for whom the earth herself spontaneously pours forth from her own substance varied and abundant pasture without any labour on their part, while for us such nourishment is scarcely or barely sufficient even when we seek it with great toil. Nor can I in any way accept that the highest good of a beast and of a man is the same. Why is there need for such elaborate equipment in the acquisition of the finest arts? Why for such a confluence of the most honourable studies, such an accompanying retinue of virtues — if all of these are being gathered for no other end than pleasure?

\ciceroSection{112}
Consider: if Xerxes, having bridged the Hellespont and tunnelled through Athos, having marched over the sea and sailed across the land with such fleets and such cavalry and infantry forces, had come into Greece with such an impetuous advance — and if someone were to ask him the cause of so great an armament and so great a war, and he were to answer that he had wanted to carry off honey from Mount Hymettus — he would clearly seem to have undertaken such vast efforts without sufficient cause. In the same way, if we say that the wise man, equipped and adorned with the most numerous and most weighty arts and virtues — a man who encompasses not, like Xerxes, the seas on foot and the mountains by fleet, but the whole sky and the entire earth together with the sea in his understanding — is pursuing pleasure, we shall be saying that he has undertaken his great enterprise for the sake of honey.

\ciceroSection{113}
We are born for something loftier and more magnificent than that, believe me, Torquatus — and this is plain not only from the capacities of the mind: in your mind there dwells memory for things without number (in yours, indeed, without limit), there dwells the power of conjecture about what follows, differing little from divination, there dwells a sense of shame as the governor of desire, there dwells a faithful guardianship of justice for human society, there dwells a firm and steadfast contempt for pain and death in the endurance of hardships and the confronting of dangers — these things dwell in the mind; but consider also the limbs themselves and the senses, which will seem to you, like the rest of the parts of the body, not merely companions of the virtues but their very ministers and servants.

\ciceroSection{114}
Consider further: if in the body itself many things are to be ranked above pleasure — strength, health, swiftness, beauty — what, pray, do you suppose is the case with the mind? Those ancient masters of the highest learning believed that something celestial and divine resides within it. Now if the highest good were pleasure, as you claim, it would be desirable to spend one's days and nights absorbed in the greatest possible pleasure, without interruption, with all the senses suffused and stirred by every sweetness. Yet who is there, worthy of the name of man, who would wish to spend even a single day wholly in that kind of pleasure? The Cyrenaics make no scruple of wishing it; your own school expresses this more modestly — the Cyrenaics, perhaps, with greater consistency.

\ciceroSection{115}
But let us pass in review not those supreme arts, the absence of which caused men to be called idle and inert by our ancestors, but let me ask you this: do you suppose that Homer, Archilochus, Pindar — to say nothing of them — or Phidias, Polyclitus, Zeuxis directed their arts toward pleasure? Shall a craftsman set more before himself in the way of beauty of form than an outstanding citizen in the beauty of his actions? And what other cause is there for an error so great, spread so far and wide, if not that the man who pronounces pleasure the highest good deliberates not with that part of the mind in which reason and judgment reside, but with appetite — that is, with the most frivolous part of the mind? I ask you plainly: if the gods exist, as you yourselves believe, and if they can be happy without receiving bodily pleasures, or if without that kind of pleasure they are happy — then why do you refuse to allow that the wise man can enjoy the same use of his mind?

\ciceroSection{116}
Read the funeral eulogies, Torquatus — not those of men praised by Homer, not that of Cyrus, not those of Agesilaus, Aristides, or Themistocles, not those of Philip or Alexander; read those of our own people, read those of your own family. You will find no one praised in such terms as to be called a skilled contriver of pleasures. That is not what the inscriptions on monuments declare — as, for instance, this one, near the city gate:

\begin{quote}
\textit{This one man the most numerous peoples agree to have been the foremost man of his people.}
\end{quote}

\ciceroSection{117}
Do we suppose that the most numerous peoples agreed that Calatinus was the foremost man among the people because he had been supremely excellent at assembling pleasures? And shall we say that young men show good promise and great natural ability when we expect them to serve their own convenience and do whatever suits their own interest? Do we not see what vast disorder of all things follows, what a sweeping confusion? Beneficence is destroyed; gratitude is destroyed — and these are the bonds of harmony. For when you oblige someone for your own benefit, that cannot be counted a benefit but a loan put out at interest; and gratitude does not seem to be owed to one who has done a favor purely for his own sake. The greatest virtues without exception are of necessity laid low when pleasure holds the mastery. There are also many base actions which, unless honorable conduct carried the greatest weight by its own nature, it would not be easy to defend as things a wise man could not fall into.

\ciceroSection{118}
And so as not to range more widely — for the instances are without number — virtue, once it is properly praised, must of necessity bar all approaches to pleasure. Do not wait any longer for me to demonstrate this. Look into your own mind yourself; turning it over in every direction of thought, ask yourself whether you would rather spend your whole life in uninterrupted pleasures, in that tranquility you were so fond of invoking, free from pain — and with the addition, which you on your side are in the habit of appending, though it is in fact impossible, of freedom from any fear of pain — or whether you would prefer, while serving all peoples with the greatest possible benefit, while carrying relief and salvation to those in need, to endure even the labors of Hercules. For that is precisely how our ancestors described burdensome toil as something not to be fled: they even applied the word \textit{aerumnae} — a word of the bleakest cast — to what a god endured.

\ciceroSection{119}
I would draw this answer out of you and compel you to reply, were it not that I feared you would say that Hercules himself performed, for the salvation of mankind, his mighty labors — for the sake of pleasure. When I had said this, Torquatus replied: "I have people to whom I can refer all of that — and though I could answer something myself, I would rather find them better prepared." "You mean our mutual friends Siro and Philodemus," I said, "men as excellent in character as they are supremely learned." "You understand me exactly," he said. "Very well," I said; "but it would have been more fitting for Triarius to arbitrate our disagreement." "I recuse myself," said Triarius with a smile, "as an improper judge on this question at least; for you press us gently, while this man here harries us in the Stoic fashion." Then Triarius said: "From now on I shall engage more boldly; for I shall have close at hand the very arguments I have just heard, and I will not come to grips with you until I have seen you instructed by the men you mention." When this had been said, we brought both our walk and our discussion to an end.


\ciceroBook{3}

\ciceroSection{1} Pleasure herself, Brutus, were she to plead her own case and not have such stubborn advocates, would, I think, yield to dignity, convicted as she was in the previous book. For surely she would be shameless to keep on fighting against virtue, or to set what is agreeable above what is honourable, or to maintain that the body's sweetness, and the gladness born of it, weighs more than firmness and steadiness of mind. And so let us dismiss her, and bid her keep within her own bounds, lest the rigour of our argument be hampered by her coaxings and enticements.

\ciceroSection{2} For we must seek where that highest good lies which we mean to find, since pleasure has been removed from it, and much the same can be said against those who held the absence of pain to be the end of goods; nor indeed ought any highest good to win approval that is without virtue, than which nothing can be more excellent. And so, although in that conversation which was held with Torquatus we were not slack, still this present contest with the Stoics is the sharper one to be readied. For what is said about pleasure is argued with no great keenness and no great depth; neither are those who defend it crafty in debate, nor do those who speak against it have a difficult case to repel.

\ciceroSection{3} Epicurus himself even says that there is no need to argue about pleasure at all, since the judgment of it lies in the senses, so that it is enough for us to be reminded of it, and useless to be taught about it. For us, therefore, the argument on both sides was straightforward. For there was nothing tangled or twisting in Torquatus' discourse, and our own speech, as it seems to me, was clear. But you are not unaware how subtle, or rather how thorny, the Stoics' manner of arguing is — and this both for the Greeks and, still more, for us, who must coin words as well, and lay down new names for new things. And no one even moderately learned will wonder at this, if he reflects that in every art whose practice is not common and widely shared there is much novelty of terms, since the words for those matters that each art handles have to be established.

\ciceroSection{4} And so the logicians and the natural philosophers use words that are not known even to Greece itself; the geometers, the musicians, and the grammarians too speak after a certain fashion of their own. The very arts of the rhetoricians, which are wholly of the forum and the public square, use words nonetheless that are, so to speak, private and their own in their teaching. And to leave aside these refined and liberal arts, not even craftsmen could maintain their crafts unless they used words unknown to us but familiar to themselves. Indeed even farming, which shrinks from all finer elegance, has nonetheless marked with new names the matters in which it is engaged. So much the more must the philosopher do this; for philosophy is the art of life, and one who discourses on it cannot snatch his words from the marketplace.

\ciceroSection{5} And yet, of all the philosophers, the Stoics have made the most innovations, and Zeno, their founder, was not so much a discoverer of things as of new words. But if, in that language which most people think the richer, the Greeks were granted leave that their most learned men should use unfamiliar words about matters out of the common run, how much more is it to be granted to us, who now venture to handle these things for the first time? And since I have often said — and that with some complaint, not only against the Greeks, but against those too who wish to be reckoned Greeks rather than our own — that we are not merely not outdone by the Greeks in copiousness of words, but are actually their superiors in it, we must work to achieve this not only in our own arts but in those very arts of theirs. And yet those words which, on the precedent of the ancients, we use as Latin — like \textit{philosophia} itself, like \textit{rhetorica}, \textit{dialectica}, \textit{grammatica}, \textit{geometria}, \textit{musica} — although they could have been rendered in Latin, still, since they have been taken up by usage, let us reckon them our own.

\ciceroSection{6} So much, then, about the names of things. But about the things themselves I often fear, Brutus, that I may be faulted, since I write these matters to you, who have gone so far both in philosophy and in the best kind of philosophy. Were I doing this as if to instruct you, I should be rightly faulted. But I am far from that, and I send these things to you not that you may learn what is already most familiar to you, but because I find rest most easily in your name, and because I have in you the fairest assessor and judge of those studies which I share with you. You will attend, then, as you are wont, with care, and will decide that controversy which I had with your uncle, a man godlike and beyond compare.

\ciceroSection{7} For when I was at my place near Tusculum and wished to use certain books from the library of young Lucullus, I came to his villa to fetch them myself, as was my custom. When I had come there, I saw M. Cato, whom I had not known to be present, sitting in the library surrounded by a great many books of the Stoics. For there was in him, as you know, an avidity for reading, and he could not be sated — he who, not even dreading the empty censure of the crowd, used often to read in the senate house itself while the senate was being assembled, withdrawing no effort from the commonwealth. So much the more, then, in fullest leisure and with so great a store at hand, did he seem to gorge himself on books, if that word may be used in so glorious a matter.

\ciceroSection{8} And when it happened that each of us caught sight of the other unawares, he rose at once. Then those first words we use at meeting: 'What brings you here?' he said. 'From your villa, I suppose,' and: 'Had I known you were there, I would have come to you myself.' 'Yesterday,' I said, 'once the games had begun, I set out from the city and arrived toward evening. My reason for coming here was to fetch certain books from this place. And indeed, Cato, this whole store ought soon to be known to our Lucullus; for I would rather he took his delight in these books than in the rest of the villa's adornment. For it is a great care to me — though this, to be sure, is properly your charge — that he be so schooled as to prove worthy of his father, and of our Caepio, and of you, who are so near to him in blood. And I take this trouble not without cause; for I am moved by the memory of his grandfather — you are not unaware how highly I valued Caepio, who, as my judgment runs, would already be among our leading men were he alive — and Lucullus too is ever before my eyes, a man both pre-eminent in every virtue and joined to me by friendship and by complete agreement of will and judgment.'

\ciceroSection{9} 'You do nobly,' he said, 'in that you both keep the memory of those men, each of whom commended his children to you by his will, and love the boy. As for the charge you call mine, I do not indeed refuse it, but I take you on as my partner. I add this too: that the boy already gives me many signs both of modesty and of natural gift — but you see his age.' 'I see it indeed,' I said, 'but he ought even now to be steeped in those arts which, if he drinks them in while he is still tender, he will come to greater things the better prepared.' 'So it shall be, and indeed we shall speak of these matters together more carefully and more often, and act on them in common. But let us sit down, if you please,' he said. And so we did.

\ciceroSection{10} Then he said: 'But you, who have so many books yourself, what after all are you seeking here?' 'Certain commentaries of Aristotle's,' I said, 'which I knew to be here, I came to carry off, that I might read them while I was at leisure — which does not often fall to us.' 'How I wish,' he said, 'that you had inclined toward the Stoics! For surely, if it belonged to anyone, it belonged to you to count nothing among goods but virtue.' 'Take care it has not been rather yours,' I said, 'since you held in substance the same view as I, not to fasten new names upon things. For our reasoning agrees; it is our language that wars.' 'By no means does it agree,' said he. 'For whatever you call worth seeking, and reckon among goods, apart from what is honourable, you will both put out the very light of virtue, as it were, and overturn virtue from its foundations.' 'These things are grandly said, Cato,' I replied, 'but do you not see that you share the glory of your words with Pyrrho and with Aristo, who make all things equal?

\ciceroSection{11} On these men I am eager to know what you think.' 'You ask me what I think?' he said. 'Those good men, brave, just, and temperate, whom we have either heard to have lived in the commonwealth or have ourselves seen — men who, with no doctrine, following nature herself, did many praiseworthy deeds — those men, I hold, were better fashioned by nature than they could have been fashioned by philosophy, had they approved any other than the one that counts nothing among goods but the honourable, and nothing among evils but the base. As for the rest of the philosophers' systems — one more so than another, to be sure, but all of them — those that count anything devoid of virtue either among goods or among evils, I judge that they not only help us not at all, nor make us firmer toward becoming better, but actually corrupt nature herself. For unless this is held fast, that the only good is what is honourable, in no way can it be proved that the happy life is brought about by virtue. And if that be so, I do not know why effort should be given to philosophy. For if any wise man could be wretched, then truly I should think that vaunted and memorable virtue not much to be valued.'

\ciceroSection{12} 'What you have so far said, Cato,' I replied, 'you could say in the very same way if you followed Pyrrho or Aristo. For you are not unaware that to these men the honourable seems not only the highest good but, as you would have it, the only good. And if that is so, it follows — the very thing I see you wishing — that all the wise are always happy. Is it these men, then, you praise, and do you judge that we ought to follow this view of theirs?' 'By no means theirs, at any rate,' he said. 'For since this is proper to virtue, to exercise choice among those things that are in keeping with nature, those men who have so levelled all things as to render them equal on either side, so that they made no use of selection — those men have done away with virtue itself.'

\ciceroSection{13} 'That, indeed, you put excellently,' I said; 'but I ask whether you must not do the same, you who call nothing good that is not right and honourable, and who take away every distinction among the rest of things.' 'If indeed I were taking it away,' he said; 'but I leave it standing.' 'In what way?'

\ciceroSection{14} I said. 'If there is one virtue, one single thing — this which you call honourable, right, praiseworthy, seemly — for it will be the better known what its quality is when it is marked by several terms declaring the same thing — if, then,' I said, 'this is the only good, what will you have besides that you might pursue? Or if nothing is evil except what is base, dishonourable, unseemly, crooked, shameful, foul — that we may make this too conspicuous by several names — what besides will you say is to be shunned?' 'To you, who are not ignorant what I am about to say,' he said, 'but who, as I suspect, wish to snatch something from my brief reply, I will not answer point by point; I will rather unfold — since we are at leisure, unless you think it out of place — the whole view of Zeno and the Stoics.' 'Not at all out of place,' I said, 'and that unfolding of yours will do much toward the matters we are seeking.'

\ciceroSection{15} 'Let us try, then,' he said, 'although this Stoic system has in it something rather difficult and obscure. For when, in the Greek tongue, these very names of new things once … did not seem strange, names which now a long-standing usage has worn smooth, what do you suppose will be the case in Latin?' 'That, surely, is most easy,' I said. 'For if Zeno was permitted, when he had discovered some unfamiliar thing, to fasten upon it a name unheard of as well, why should it not be permitted to Cato? Nor yet will it be necessary to render word for word, as unskilled interpreters are wont, when there is a more usual word that declares the same thing. For my part, I am even accustomed, where the Greeks use a single word, if I cannot do otherwise, to set out the same thing in several. And yet I think we ought to be granted leave to use a Greek word, whenever a Latin one should less readily come — so that this licence not be granted to ``saddle-cloths'' and ``wine-coolers'' rather than to \textit{[Greek: proēgmena]} and \textit{[Greek: apoproēgmena]}; though these, to be sure, it will be right to call ``the preferred'' and ``the rejected.'''

\ciceroSection{16} 'You do well,' he said, 'in helping me; and for those terms you have just mentioned I shall use the Latin instead, while in the rest you will come to my aid if you see me stuck.' 'I shall do so eagerly,' I said. 'But ``fortune favours the bold''; so make the attempt, I beg. For what more godlike thing can we do than this?' 'It is the view of those,' he said, 'whose reasoning I approve, that as soon as an animal is born — for here is where we must begin — it is reconciled to itself and commended to itself, toward its own preservation, and toward its own state and the things that preserve that state; while it is estranged from its own destruction and from those things that seem to bring destruction on. They prove this to be so on the ground that, before pleasure or pain has touched them, the young seek what is wholesome and spurn the contrary — which would not happen unless they loved their own state and feared destruction. And it could not be that they should seek anything unless they had a sense of themselves, and thereby loved themselves. From which it ought to be understood that the starting-point is drawn from self-love.'

\ciceroSection{17} Among the first natural objects of self-love, most Stoics hold that pleasure should not be reckoned. With them I emphatically agree, lest, if nature should seem to have placed pleasure among the things first sought after, many shameful consequences should follow. As for why we love those things first taken up by nature, there seems to be argument enough in this: that there is no one who, when either course is open to him, would not rather have all the parts of his body fit and whole than, with the same use available, maimed or twisted. And as for the apprehensions of things — which we may call comprehensions, or perceptions, or, if these words please less or are less understood, \textit{[Greek: katalēpseis]} — these, then, we judge are to be taken up for their own sake, because they hold within them something that, as it were, embraces and contains the truth. This can be discerned in small children, whom we see take delight, even where nothing turns on it for them, if by their own reasoning they have found something out for themselves.

\ciceroSection{18} The arts too we think are to be taken up for their own sake, both because there is in them something worthy of being taken up, and because they consist of apprehensions and contain within them something established by reason and method. From a false assent, moreover, they hold that we are more estranged than from any other thing that runs counter to nature. As for the limbs — that is, the parts of the body — some seem to have been bestowed by nature for the sake of their use, such as the hands, the legs, the feet, and those parts that lie within the body, the usefulness of which is even debated among physicians; while others were given for no usefulness at all, but as a kind of adornment — like the tail to the peacock, the iridescent plumage to doves, and to men the breasts and the beard.

\ciceroSection{19} These things are perhaps stated rather drily; for they are, as it were, the first elements of nature, to which a richness of language can scarcely be applied, nor do I for my part intend to chase after it. Yet still, when you speak of grander matters, the matter itself carries the words away; and so the discourse becomes at once weightier and more brilliant. ``It is as you say,'' I replied. ``But still, everything stated lucidly about a good subject seems to me splendidly stated. To wish to speak of matters of this kind ornately would be childish; to be able to set them out plainly and perspicuously is the mark of a learned and intelligent man.''

\ciceroSection{20} ``Let us go forward, then,'' he said, ``since we have set out from these first principles of nature, with which what follows must agree. And this is the first division. They say that a thing is `estimable' — for so, I think, let us call it — which is either itself in keeping with nature or brings about something such that it is worthy of selection on this account, that it possesses some weight worthy of esteem, which they call \textit{[Greek: axia]}; and on the other hand `inestimable,' which is the contrary of the foregoing. With the first principles thus laid down, that the things in keeping with nature are to be taken up for their own sake, and their contraries likewise to be rejected, the first proper function — for so I render \textit{[Greek: kathēkon]} — is that one should preserve oneself in the state of nature; next, that one should hold to the things in keeping with nature and drive off their contraries. Once this selection and likewise rejection has been discovered, there follows next selection joined with proper function, then that selection continuously, and at last selection steadfast and in accord with nature; and in this there first begins to be present, and to be understood, what may truly be called good.

\ciceroSection{21} For the first appropriation of a human being is to the things in keeping with nature. But as soon as he has grasped understanding — or rather a notion, which they call \textit{[Greek: ennoia]} — and has seen the order of things to be done and, so to speak, their concord, he then esteems that order far more highly than all those things he had first loved; and so by knowledge and reason he gathers the conclusion that herein is set that highest good of a human being, to be praised and sought for its own sake. And since this good is set in what the Stoics call \textit{[Greek: homologia]} — let us call it concord, if you please — since, then, in this lies the good to which all things are to be referred, honourable deeds and the honourable itself, which alone is reckoned among goods, although it arises later, nonetheless is the only thing to be sought by its own force and worth; whereas of the things that are first in nature, none is to be sought for its own sake.

\ciceroSection{22} But since those things that I called proper functions set out from the first principles of nature, they must necessarily be referred to these, so that it may rightly be said that all proper functions are referred to the end of our attaining the first principles of nature — yet not so that this should be the ultimate good, since honourable action is not present in the first appropriations of nature; for it is a consequence, and arises later, as I said. It is nonetheless in keeping with nature, and it urges us to seek it for its own sake far more than all the foregoing things do. But here, first of all, an error must be removed, lest anyone suppose it to follow that there are two ultimate goods. For just as, if someone has set himself to aim a spear or an arrow at something, then — as we speak of the ultimate good — so he must do everything he can to take aim, and yet, while he does everything to attain his purpose, his taking aim is, as it were, the ultimate, what we call the highest good in life, whereas his actually striking the mark is, as it were, to be selected, not to be sought.

\ciceroSection{23} But since all proper functions set out from the first principles of nature, wisdom itself must necessarily set out from these same principles. Yet just as it often happens that a man recommended to another esteems the one to whom he is recommended more than the one by whom he was, so it is by no means strange that at first we are recommended to wisdom by the first principles of nature, but afterwards wisdom itself becomes dearer to us than the things by which we came to it. And just as our limbs have been given to us in such a way that they plainly appear to have been given for a certain manner of living, so the impulse of the mind, which in Greek is called \textit{[Greek: hormē]}, seems to have been given not for any sort of life whatever, but for a certain form of living; and the same holds of reason, and of perfected reason.

\ciceroSection{24} For just as to the actor his gesture, to the dancer his movement, is given not as any sort whatever but as something definite, so life must be lived in a certain definite manner, and not in any manner at all — the manner we call fitting and consonant. For we do not judge wisdom to be like navigation or medicine, but rather like that acting I just mentioned, and like dancing, so that the ultimate thing — that is, the performance of the art — lies within the art itself, and is not sought from outside. And yet there is also a certain unlikeness between wisdom and these very arts, in that the things rightly done in those arts do not nonetheless contain all the parts of which they consist; whereas the things we call right, or rightly done — if you please, but they call them \textit{[Greek: katorthōmata]} — contain all the measures of virtue. For wisdom alone is wholly turned in upon itself, which is not the case with the other arts.

\ciceroSection{25} But it is clumsy to compare the ultimate end of medicine and of navigation with the ultimate end of wisdom. For wisdom embraces greatness of soul and justice, and the power to judge that all that befalls a human being lies beneath it — which does not fall to the other arts. And no one will be able to hold those very virtues I just mentioned unless he has settled that there is nothing that matters or differs one from another, except the honourable and the base.

\ciceroSection{26} Let us now see how splendidly these things follow from what I have already laid down. Since this is the ultimate — for you perceive, I believe, that for some time now I have been rendering what the Greeks call \textit{[Greek: telos]} sometimes by `the ultimate,' sometimes `the last,' sometimes `the highest'; and it will be allowed to say `the end' too in place of `the ultimate' or `the last' — since, then, this is the ultimate, to live in agreement and in keeping with nature, it follows of necessity that all the wise always live happily, completely, and prosperously, hindered by nothing, prevented by nothing, lacking nothing. Now the thing that holds together not only the system of which I speak but our life and fortunes too — namely, that we should judge the honourable alone to be good — can indeed be amplified and adorned rhetorically, at length and copiously, with the choicest words and the weightiest sentiments; but the brief, sharp inferences of the Stoics delight me. Their arguments, then, are drawn to a conclusion thus.

\ciceroSection{27} Whatever is good is praiseworthy; whatever is praiseworthy is honourable; therefore whatever is good is honourable. Does this seem soundly enough concluded? Surely; for you see that what was produced from the two premises that were assumed is what stands concluded. Now of the two premises from which the conclusion was produced, the first is commonly disputed: that not every good is praiseworthy. For that the praiseworthy is honourable is granted. But it is utterly absurd that there should be something good which is not to be sought, or to be sought which is not pleasing, or, if pleasing, not also to be loved; and so also to be approved; and thus also praiseworthy; and that is honourable. So it comes about that whatever is good is also honourable.

\ciceroSection{28} Next I ask: who could boast either of a wretched life or of one not blessed? Of a blessed life alone, then. From which it is concluded that the blessed life is worthy of boasting, so to speak — which could not rightly fall to any but an honourable life. So it comes about that the honourable life is the blessed life. And since the man to whom it falls to be justly praised has something distinguished, marking him for honour and glory, so that on account of these things, great as they are, he can justly be called blessed, the same will most rightly be said of the life of such a man. Thus, if the blessed life is discerned by honour, the honourable alone must be held to be good. And what of this?

\ciceroSection{29} Could it in any way be denied that no one can ever be rendered of steadfast and firm and great soul — the man we call brave — unless it has been settled that pain is no evil? For just as one who reckons death among evils cannot but fear it, so no one can fail to be anxious about, and to despise, the thing he has decreed to be an evil. With this laid down and approved by everyone's assent, this further point is taken up: that the man of great and brave soul looks down upon all that can befall a human being and counts it as nothing. And since this is so, it has been established that nothing is an evil which is not base. And that lofty and surpassing man, of great soul, truly brave, counting all things human beneath him — he, I say, whom we wish to fashion, whom we are seeking — he certainly ought to have confidence in himself and in his life, both as it has been led and as it is to come, and ought to judge well of himself, settling that no evil can befall the wise man. From which is understood that same point: that the good alone is the honourable, and that this is to live blessedly — to live honourably, that is, with virtue.

\ciceroSection{30} I am not, indeed, unaware that the philosophers have held varied opinions — I mean those who placed the highest good, which I call the ultimate, in the mind. Although some have followed these opinions wrongly, nonetheless I set above all of them, of whatever kind they may be, those who placed the highest good in the mind and in virtue — above not only those three who severed virtue from the highest good, when they placed in the highest goods either pleasure or freedom from pain or the first principles of nature, but also those other three who thought virtue would be defective without some addition, and for that reason each added one of those three things I named above.

\ciceroSection{31} But there are nonetheless some who are utterly absurd: those who said that the ultimate good is to live with knowledge, and those who said there is no difference among things — and that thus the wise man will be blessed, preferring nothing to anything else by any weight at all — and those who, as certain Academics are said to have held, made the ultimate good and the highest task of the wise man to be resistance to presentations and the firm withholding of his assent. To each of these a copious reply is commonly made, but matters that are evident need not be long. For what is plainer than this: that if there were no selection of the things in keeping with nature out of those that are contrary to nature, all that prudence which is sought and praised would be done away with? Once those opinions, then, that I have laid down are circumscribed — and any that are like them — what remains is that the highest good is to live applying knowledge of those things that happen by nature, selecting what is in keeping with nature and rejecting what is contrary to it; that is, to live in keeping with nature and in agreement with it.

\ciceroSection{32} But in the other arts, when a thing is said to be done with skill, it must be reckoned, in a manner of speaking, as something later and consequent — what they call \textit{[Greek: epigennēmatikon]}; whereas in the case of what we call done with wisdom, it is from the very first most rightly so called. For whatever proceeds from wisdom must straightway be complete in all its parts; for in that lies what we say is to be sought. For just as it is a sin to betray one's country, to do violence to one's parents, to plunder temples — sins that lie in the act — so to fear, so to grieve, so to be in the grip of lust, is a sin even apart from any act. And just as these latter are sins not in their later and consequent stages but straightway in their first, so the things that proceed from virtue are to be judged right by their first undertaking, not by their completion.

\ciceroSection{33} The good, moreover, a word that has been used so often in this discussion, is also unfolded by definition. But their definitions differ among themselves only very slightly, and even so they look toward the same thing. I, for my part, agree with Diogenes, who defined the good as that which is by nature complete. And following from this, he said that what is beneficial too — let us call it a ``benefit,'' for so the word stands \textit{[Greek: ōphelēma]} — is a motion or a state proceeding from that which is by nature complete. And since notions of things arise in our minds whenever something has been grasped through experience, or through combination, or through resemblance, or through the comparison made by reason, it is by this fourth means, which I placed last, that the notion of the good has come to be. For when the mind ascends, by the comparison made by reason, from those things that are in keeping with nature, then it arrives at the notion of the good.

\ciceroSection{34} But this very good we both feel and call good not by any addition, nor by growth, nor by comparison with other things, but by its own proper force. For just as honey, though it is exceedingly sweet, is felt to be sweet by its own proper kind of flavour and not by comparison with other things, so this good of which we are speaking is indeed to be valued most highly, yet that valuation holds good by its kind, not by its magnitude. For since valuation — which is called \textit{[Greek: axia]} — has been counted neither among goods nor, again, among evils, however much you add to it, it will remain within its own kind. The proper valuation of virtue, then, is something different, which holds good by its kind and not by growing.

\ciceroSection{35} Nor, indeed, are the disturbances of the mind — which make the life of the foolish wretched and bitter, and which the Greeks call \textit{[Greek: pathē]} (I might have rendered the word itself and called them ``diseases,'' but that would not fit them all; for who is in the habit of calling pity, or anger itself, a disease? but they call it \textit{[Greek: pathos]}) — let it be called, then, a ``disturbance,'' which by its very name seems to be declared a flaw, and these disturbances are not set in motion by any natural force. And they are all four in kind, more numerous in their parts: distress, dread, craving, and what the Stoics call by a name common to body and mind \textit{[Greek: hēdonē]}, which I prefer to call ``gladness,'' the pleasurable elation, as it were, of an exultant mind. The disturbances, moreover, are stirred by no force of nature, and they are all opinions and the judgments of an unsteady mind. And so the wise man will always be free of them.

\ciceroSection{36} That everything honourable is to be sought for its own sake — this we hold in common with the views of many other philosophers. For apart from the three schools that shut virtue out from the highest good, this view must be upheld by all the rest of the philosophers, and most of all by these Stoics, who held that nothing whatever should be counted among goods except the honourable. But this defence, certainly, is very easy and very ready to hand. For who is there, or who has there ever been, of so burning a greed or so unbridled in his cravings, that the very thing he would be willing to obtain by any crime whatever he would not far rather come to without an outrage — even with total impunity laid before him — than in that other way?

\ciceroSection{37} And again, what advantage, or what profit, are we seeking when we desire to know those things that are hidden from us — how they are moved, and from what causes, those things that turn in the heavens? Who, moreover, lives by manners so boorish, or who has so violently hardened himself against the pursuits of nature, that he shrinks from things worth knowing, and does not seek them out for their own sake, apart from any pleasure or advantage, and counts them as nothing? Or who is there, learning of the deeds, the sayings, the counsels of our forefathers — whether of the Africani, or of that great-grandfather of mine whom you forever have on your lips, or of the rest of those brave men outstanding in every virtue — who is touched by no pleasure in his mind?

\ciceroSection{38} And who, brought up and reared honourably in a respectable family, is not offended by baseness itself, even when it is not going to harm him? Who looks with an unruffled mind upon a man he thinks lives foully and disgracefully? Who does not hate the sordid, the empty, the fickle, the worthless? And what will it be possible to say, if we lay it down that baseness is not to be fled for its own sake, to keep men, once they have found darkness and solitude, from holding back from no shameful thing — unless baseness itself deters them by its own foulness? Countless things can be said to this effect, but there is no need. For there is nothing about which there can be less doubt than that the honourable is to be sought for its own sake, and in the same way the base is to be fled for its own sake.

\ciceroSection{39} Now once that point has been established, of which we spoke before — that the honourable alone is the good — it must be understood that the honourable is to be valued more highly than those intermediate things that are procured from it. But when we say that folly and cowardice and injustice and intemperance are to be fled on account of the consequences that come of them, we do not say it in such a way that this language should seem to clash with what was laid down, that the base alone is evil; for they are not referred to bodily harm, but to base actions, which arise out of the vices. For what the Greeks call \textit{[Greek: kakiai]} I prefer to name ``vices'' rather than ``vicious deeds.''

\ciceroSection{40} ``Why, Cato,'' I said, ``you do it with brilliant words, and words that declare just what you mean! And so you seem to me to teach philosophy in Latin and to grant it, as it were, citizenship. Until now it appeared to be a foreigner sojourning at Rome and not to offer itself to our conversations — and this above all on account of a certain refined fineness both of subjects and of words. For I know there are some who can philosophize in any language whatever; for they use no divisions, no definitions, and they say themselves that they approve only those things to which nature silently assents. And so, in matters not at all obscure, the labour of disputation among them is not great. For which reason I attend to you closely, and whatever names you impose on the things this discussion is about, I commit to memory; for perhaps I too shall have to use those same terms. To the virtues, then, you seem to me to have set their contrary vices most rightly, and in keeping with the usage of our own speech. For whatever is blameworthy in and of itself, I think it is for that very reason named a vice — or even that to be blamed is said from the vice. But if you had called \textit{[Greek: kakia]} ``viciousness,'' Latin usage would carry us over to some other single particular vice. As it is, to every virtue a vice is opposed under a contrary name.''

\ciceroSection{41} Then he said: ``These things, then, being thus laid down, there follows a great contest. The Peripatetics handled it rather softly — for their manner of speaking is not acute enough, owing to their ignorance of dialectic — but your Carneades, with a certain remarkable training in dialectic and with the highest eloquence, brought the matter to the sharpest crisis, since he did not cease to contend, in this whole question that is called the question of goods and evils, that the Stoics had no dispute with the Peripatetics over substance, but only over names. To me, however, nothing seems so plain as that these views of those philosophers differ from one another more in substance than in words; I declare that there is far greater discrepancy in substance between the Stoics and the Peripatetics than in words, seeing that the Peripatetics say that everything they themselves call good has a bearing on living happily, whereas our school does not think the happy life is made complete out of everything that is worthy of some valuation.

\ciceroSection{42} Or can anything be more certain than this — that on the reasoning of those who place pain among evils, the wise man cannot be happy when he is racked on the rack; whereas on the reasoning of those who do not count pain among evils, it surely follows of necessity that the happy life is preserved for the wise man in the midst of every torment? For if men endure the same pains more tolerably when they take them on for their country than when the cause is a slighter one, it is opinion, not nature, that makes the force of the pain either greater or less.

\ciceroSection{43} Nor is even this consistent: that if, since there are three kinds of goods — which is the Peripatetics' view — each man is the happier the fuller he is of bodily or external goods, this same thing should be approved by us as well, so that whoever has more of those things that are highly valued in the body is the happier. For they think the happy life is made complete out of bodily advantages, but our school thinks nothing of the kind. For since it is our doctrine that the happy life is not made happier, nor more to be sought, nor more highly to be valued, even by an abundance of those goods that we may truly call good, surely a multitude of bodily advantages bears still less upon the happy life.

\ciceroSection{44} For indeed, if both to be wise and to be healthy were to be sought, then the two joined together would be more to be sought than wisdom alone; and yet, even if both are worthy of valuation, the two joined would not be worth more than wisdom by itself. For we who judge health worthy of some valuation, and yet do not place it among goods, hold all the same that there is no valuation so great that it should be set before virtue. The Peripatetics do not hold this same thing; for they must say that an action that is both honourable and free of pain is more to be sought than the same action would be if it were attended with pain. To us it seems otherwise — whether rightly or not, later; but can there be a greater dissension over substance?

\ciceroSection{45} For just as the light of a lamp is dimmed and overwhelmed by the light of the sun, and just as a drop of honey is lost in the vastness of the Aegean sea, and just as a farthing added to the riches of Croesus, or a single step on the road that runs from here to India, amount to nothing — so, since the end of goods is what the Stoics declare it to be, all that valuation of bodily things must of necessity be dimmed and buried by the splendour and magnitude of virtue, and must perish away. And just as opportuneness — for so let us call good timing \textit{[Greek: eukairia]} — is not made greater by the lengthening of time (for things that are called opportune have their own due measure), so right performance — for thus I render \textit{[Greek: katorthōsis]}, since a right deed is \textit{[Greek: katorthōma]} — right performance, then, and likewise conformity, and finally the good itself, which consists in this, that it accords with nature, has no increase by addition.

\ciceroSection{46} For just as that opportuneness, so these things of which I have spoken are not made greater by the lengthening of time; and for this reason the happy life does not seem to the Stoics more desirable or more to be sought if it is long than if it is short. And they use a comparison: that as, if it were the merit of a buskin to fit the foot snugly, neither would many buskins be set before few, nor larger ones before smaller, so those whose every good is bounded by conformity and opportuneness will set neither more before fewer nor longer-lasting before briefer. Nor, indeed, do they argue acutely enough who say:

\ciceroSection{47} if good health is to be valued more highly when long than when short, then the longest exercise of wisdom too is worth the most. They do not understand that the valuation of health is judged by its span, that of virtue by its opportuneness — so that those who say the one would seem bound to say this too, that a good death and a good birth are better when long than when short. They do not see that some things are valued more highly for their brevity, others for their long duration.

\ciceroSection{48} And so it is consistent with what has been said that, on the reasoning of those who think that end of goods which we call the uttermost, the ultimate, can grow, it should likewise be their view that one man is wiser than another, and that, in the same way, one man can either sin or act rightly more than another — which it is not permitted to us to say, who do not think the end of goods can grow. For just as those who are sunk in the water can breathe no more, even if they are not far from the surface, so as to be on the very point of emerging, than if they were still in the depths; and as that puppy which already draws near to seeing sees no more than the one just born — so the man who has advanced somewhat toward the disposition of virtue is no less in wretchedness than the man who has advanced not at all. I am aware that these things seem astonishing; but since the foregoing points are surely firm and true, and these are consistent with them and follow from them, there is no room to doubt the truth of these either. And yet, although they deny that either virtues or vices can grow, they nonetheless think that each of the two can in a manner be poured out and, as it were, spread wider. Riches, however, in Diogenes' judgment, have only this force: that they serve as guides, as it were, toward pleasure and toward good health;

\ciceroSection{49} —but even granting that they hold such things together, they do not do the same for virtue as for the other arts, toward which money can serve as a guide, though it cannot hold them together. And so, if pleasure or sound health belongs among goods, then riches too must be placed among goods; but if wisdom is a good, it does not follow that we should call riches a good as well. Nor can that which belongs among goods be held together by anything that does not itself belong among goods; and for this reason, since the cognitions and graspings of things, out of which the arts are produced, set our impulse in motion, then—seeing that riches do not belong among goods—no art can be held together by riches.

\ciceroSection{50} And even if we were to concede this point about the arts, the account of virtue would still not be the same, because virtue requires a great deal of reflection and exercise, which is not the case in the arts, and because virtue embraces the stability, the firmness, and the constancy of one's whole life—and we do not see these same qualities in the arts. Next there is unfolded the difference among things; for if we were to say there was no such difference, the whole of life would be thrown into confusion, as it is by Aristo, and no function or work of wisdom would be found, since among the very things that bear on the conduct of life there would be no difference at all, and no choice would need to be exercised. And so, once it had been firmly established that the only good is the honourable and the only evil the base, they nevertheless wished there to be, among those things which carry no weight toward living happily or wretchedly, something that still made a difference, such that some of them are valuable, others the reverse, others neither.

\ciceroSection{51} And of those that are valuable, in some there is reason enough why they should be preferred to certain others—as in health, in soundness of the senses, in freedom from pain, in glory, in riches, and in things of that sort—while others are not of this kind; and likewise of those that are worth no valuation, some have reason enough why they should be rejected—such as pain, sickness, the loss of the senses, poverty, disgrace, and the like—while others are not so. From this arose what Zeno called \textit{[Greek: proēgmenon]} and, on the contrary, what he called \textit{[Greek: apoproēgmenon]}; for he was using, in a copious language, words nonetheless coined and new—something not granted to us in this poor language of ours, although you are in the habit of calling even this one the more copious of the two. But it is not out of place, so that the force of the word may be more easily understood, to set out Zeno's reasoning in coining it.

\ciceroSection{52} For just as, he says, no one in a king's court says that the king himself has been, so to speak, advanced to high station (for that is what \textit{[Greek: proēgmenon]} means), but rather those who hold some office, whose rank ranks next, so as to be second to the royal sovereignty—so too in life it is not the things that hold first place, but the things that hold second place, that are to be named \textit{[Greek: proēgmena]}, that is, "advanced." And these we may call by that very name—it will be word for word—or "promoted" and "demoted," or, as we said a moment ago, "preferred" and "preeminent," and the others "rejected." For once the matter is understood, we ought to be accommodating in our use of words.

\ciceroSection{53} But since we say that everything which is good holds first place, it follows necessarily that this, which we name "preferred" or "preeminent," is neither good nor evil. And we define it thus: that which is indifferent, with a moderate valuation; for what they call \textit{[Greek: adiaphoron]}, this came to me in such a way that I would render it "indifferent." For it could in no way be the case that nothing should be left among the intermediates which was either in keeping with nature or contrary to it; nor, once that was left, that nothing should be placed among these things which was valuable enough; nor, once this was granted, that there should not be certain things "preferred." Rightly, then, has this distinction been drawn; and to make the matter more easily grasped, they even propose this comparison:

\ciceroSection{54} For just as, they say, if we were to imagine it to be, so to speak, the end and the ultimate goal to cast a knucklebone so that it stands upright, then a knucklebone that has been thrown so as to fall upright will have something "preferred" with respect to the end, and one that falls otherwise the reverse—and yet that "preference" of the knucklebone will have no bearing on the end I spoke of—so the things that are "preferred" are indeed referred to the end, but have no bearing on its force and nature.

\ciceroSection{55} There follows that division by which some goods bear on that ultimate end (for so I call the things that are termed \textit{[Greek: telika]}; for let us make it our practice, as we agreed, to say in several words what we cannot say in one, so that the matter may be understood), while others are productive, which the Greeks call \textit{[Greek: poiētika]}, and others are both. Among those that bear on the end, nothing is good except honourable actions; among the productive, nothing except a friend; but wisdom they hold to be both bearing on the end and productive. For since wisdom is fitting action, it falls within that bearing-on-the-end class I mentioned; but inasmuch as it brings about and produces honourable actions, it can be called productive.

\ciceroSection{56} These things that we call "preferred" are in part preferred in and of themselves, in part because they bring something about, in part on both grounds: in themselves, like a certain bearing of the mouth and the face, like one's posture, like one's movements, in which there are some things to be preferred and some to be rejected; others will be called preferred for this reason, that they produce something out of themselves, like money; and others on both grounds, like sound senses, like good health.

\ciceroSection{57} As for good repute—for it is more fitting in this place to call what they term \textit{[Greek: eudoxian]} "good repute" rather than "glory"—Chrysippus indeed and Diogenes used to say that, with its usefulness stripped away, not even a finger should be stretched out for its sake; and with these men I emphatically agree. But those who came after them, since they could not withstand Carneades, said that this good repute I have spoken of is to be preferred and taken up for its own sake, and that it belongs to a man of free birth and liberal upbringing to wish to be well spoken of by his parents, by his kin, and by good men too, and that for the sake of the thing itself, not for any use of it; and they say that, just as we wish our children's interests provided for, even if they are to be born after our death, for their own sake, so we ought also to provide for our reputation after death—for the sake of the thing itself, even with its usefulness stripped away.

\ciceroSection{58} But although we say that the honourable alone is good, it is nonetheless consistent to perform a duty, even though we place that duty neither among goods nor among evils. For there is something in these matters that is reasonable, and reasonable in such a way that an account of it can be given—therefore, such that an account can be given of having acted reasonably. And a duty is something done in such a way that a reasonable account of the deed can be given. From this it is understood that duty is something intermediate, placed neither among goods nor among their contraries. And since, among those things which are neither among the virtues nor among the vices, there is nonetheless something which can be of use, that thing is not to be done away with. And there is also a certain action of this kind, and indeed such that reason demands that we do and perform some of these things. Now what is done by reason we call a duty. A duty, then, belongs to that class which is placed neither among goods nor among their contraries.

\ciceroSection{59} And this too is plain, that the wise man does something in those intermediate matters. He judges, then, when he acts, that this is a duty. And since he is never deceived in his judging, there will be a duty in intermediate matters. This is established also by the following chain of reasoning: since we see that there is something which we should call a right action, and that this is a perfected duty, there will also be a duty begun—so that, if to return a deposit justly is among right actions, then to return a deposit is to be placed among duties; for with the addition of "justly" it becomes a right action, while the mere returning of it is in itself placed among duty. And since there is no doubt that, among the things we call intermediate, one is to be taken up and another rejected, whatever is so done or said is wholly contained within duty. From this it is understood that, since by nature all men hold themselves dear, the foolish man no less than the wise will take up what is in keeping with nature and reject its contraries. Thus there is a certain duty common to the wise man and the fool, from which it follows that duty is concerned with what we call intermediate matters.

\ciceroSection{60} But since all duties proceed from these things, it is not without reason that all our deliberations are said to be referred to them—among them, both departure from life and remaining in life. For the man in whom there are more things in keeping with nature, his duty is to remain in life; while the man in whom there are, or seem likely to be, more contrary things, his duty is to depart from life. From this it is clear both that it is sometimes the wise man's duty to depart from life though he is happy, and the fool's to remain in life though he is wretched.

\ciceroSection{61} For that good and that evil, as has often already been said, follow afterward; but those first things of nature, whether favourable or contrary, fall under the wise man's judgment and choice, and they form, as it were, the underlying material of wisdom. And so the whole account of remaining in life and of departing from it is to be measured by the things I spoke of above. For neither is that man held in life by virtue, nor is death to be sought by those who are without virtue. And it is often the wise man's duty to forsake life though he is supremely happy, if he can do so opportunely—which is to act in keeping with nature. For so they hold: that happy living is a matter of opportunity. And so wisdom prescribes that the wise man should, if it is to his purpose, forsake wisdom itself. For which reason, since the force of the vices is not such as to furnish a cause for voluntary death, it is plain that even fools, who are likewise wretched, have a duty to remain in life, if they possess the greater part of those things which we say are in keeping with nature. And since the man who departs from life and the man who remains are equally wretched, and longer duration does not make life more to be fled from for him, it is not without reason said that those who can enjoy more of the things of nature ought to remain in life.

\ciceroSection{62} They judge it relevant to the matter, moreover, to understand that it comes about by nature that children are loved by their parents. From this beginning we trace the common fellowship of the human race. This ought first to be understood from the shape and the members of our bodies, which of themselves declare that nature has held to a plan of procreation. Nor indeed could these things agree with one another—that nature should both will the begetting of offspring and yet not care that what is begotten be loved. The force of nature can be discerned even in the beasts; for when we observe their toil in bearing and in rearing their young, we seem to hear the very voice of nature. And so, just as it is plain that we shrink from pain by nature, so it is apparent that we are driven by nature herself to love those whom we have begotten.

\ciceroSection{63} From this is born the further fact that there is also a natural bond uniting man with man, such that one human being ought not to seem alien to another for this very reason, that he is a human being. For just as some of our members are, as it were, born for themselves—like the eyes, like the ears—while others also assist the use of the remaining members—like the legs, like the hands—so certain monstrous beasts are born for themselves alone; but that creature called the pinna, in its gaping shell, and the one that swims out of the shell and, because it guards it, is called the pinna-guard, and which is shut in when it has withdrawn back into the same shell—so that it seems to have given warning to beware—and likewise ants, bees, and storks do certain things for the sake of others as well. Far more closely joined are human beings. And so by nature we are fitted for gatherings, assemblies, and communities.

\ciceroSection{64} They hold, moreover, that the world is governed by the divine will of the gods, and that it is, as it were, a common city and community of men and of gods, and that each one of us is a part of this world; from which it follows by nature that we set the common advantage before our own. For just as the laws set the safety of all before the safety of individuals, so the good and wise man, obedient to the laws and not ignorant of his civic duty, provides for the advantage of all more than for that of any one man or for his own. Nor is the betrayer of his country more to be blamed than the man who deserts the common advantage or safety for the sake of his own advantage or safety. From this it comes about that the man who meets death for the commonwealth deserves praise, because it is fitting that our country should be dearer to us than we are to ourselves. And since that saying is held to be inhuman and wicked—the saying of those who declare that they would not object if, once they themselves were dead, the conflagration of all the earth should follow (a thing usually expressed in a certain common Greek verse)—it is surely true that we ought also to provide, for their own sake, for those who are one day to exist.

\ciceroSection{65} From this disposition of our minds were born the wills and the last charges of the dying. And because no one would be willing to pass his life in utter solitude, not even amid a boundless abundance of pleasures, it is easily understood that we were born for the joining and gathering together of human beings and for a natural community among them. We are driven by nature, moreover, to wish to be of use to as many as we can, and above all by teaching and by handing on the principles of practical wisdom.

\ciceroSection{66} And so it is not easy to find a man who does not pass on to another what he himself knows; so inclined are we not only to learning but to teaching as well. And just as it has been given to bulls by nature to fight with all their force and onset against lions in defence of their calves, so those who are strong in resources and able to do it — as we have been told of Hercules and of Liber — are roused by nature to the preserving of the human race. Indeed, when we call Jupiter Best and Greatest, and when we call the same god the Preserver, the God of Hospitality, the Stayer of the Line, we mean it to be understood that the safety of human beings rests in his protection. But it is altogether unfitting that we, while cheap and neglectful of one another among ourselves, should demand to be dear to the immortal gods and beloved by them. Therefore, just as we use our limbs before we have learned for the sake of what use we have them, so we are by nature joined and bound together with one another for the civic community. And were this not so, there would be no place for justice and none for goodness.

\ciceroSection{67} And just as they hold that there are bonds of law between human beings and their fellows, so they hold that a human being has no bond of law with the beasts. For Chrysippus said it splendidly: that everything else was brought to birth for the sake of human beings and of the gods, but that they were brought to birth for the sake of their own community and fellowship, so that human beings may make use of beasts for their own benefit without injustice. And since the nature of a human being is such that a kind of civic law passes between him and the human race, the man who keeps that law will be just, the man who departs from it unjust. But just as, though a theatre is common property, it can still rightly be said that the place each man has taken is his own, so in the city or in the universe — common though they are — the law does not stand in the way of each thing's belonging to its several owner.

\ciceroSection{68} And since we see that a human being was born to watch over and preserve other human beings, it is in keeping with this nature that the wise man should be willing to manage and administer the commonwealth, and, in order to live according to nature, to take a wife and to wish to have children by her. Not even chaste loves do they reckon foreign to the wise man. As for the reasoning and the way of life of the Cynics, some say that it does fall to the wise man, should some chance of that kind happen to befall him such that it must be done, while others say it does so in no way at all.

\ciceroSection{69} But in order that all fellowship, all conjunction, all affection of one human being toward another may be preserved, they held that both gains and losses — what they call benefits and harms \textit{[Greek: ōphelēmata kai blammata]} — are common; the one set of these does good, the other does harm. And they said that these are not only common but equal as well. Advantages, however, and disadvantages — for so I render what they call things well-used and things ill-used \textit{[Greek: euchrēstēmata kai dyschrēstēmata]} — they held to be common, but they would not have them equal. For those things that do good or do harm are either goods or evils, and these must of necessity be equal; but advantages and disadvantages belong to that class which we have called the preferred and the rejected, and these can fail to be equal. And so gains are said to be common, while right acts and wrongs are not held to be common.

\ciceroSection{70} They judge, moreover, that friendship is to be cultivated, because it is of that class which does good. Yet although in friendship some say that the wise man holds his friend's interest as dear as his own, while others say that each man holds his own dearer, still even these latter confess that it is foreign to justice — the justice for which we seem to have been born — to take from another anything that one may claim for oneself. By no means at all is it approved by this school of which I speak that either justice or friendship should be taken up or valued for the sake of advantages. For those very advantages will be able to undermine and overthrow them. Indeed, neither justice nor friendship will be able to exist at all, unless they are sought for their own sake.

\ciceroSection{71} Law, moreover — that which can be so named and so called — is by nature; and it is foreign to the wise man not only to do an injustice to anyone but even to do harm. Nor is it right to share or join in an injustice with friends or with those who have served us well; and it is defended most weightily and most truly that fairness can never be severed from advantage, and that whatever is fair and just is also honourable, and conversely, that whatever is honourable will also be just and fair.

\ciceroSection{72} And to these virtues, of which we have been speaking, they join also dialectic and natural philosophy, and they call both of them by the name of virtue: the one because it possesses a method by which we may neither assent to anything false nor ever be deceived by a treacherous plausibility, and by which we may hold fast and keep safe what we have learned about good things and bad. For without this art they think that anyone can be led astray from the truth and deceived. Rightly, then, if in all matters recklessness and ignorance are faulty, the art that does away with these has been named a virtue.

\ciceroSection{73} To natural philosophy too the same honour is assigned, and not without reason, because the man who is to live in agreement with nature must set out from the whole universe and from its governance. Nor indeed can anyone judge truly about good things and bad except when the whole system of nature and the life even of the gods has been understood, and whether or not the nature of a human being agrees with the nature of the whole. And those old precepts of the sages, who bid us yield to the occasion, follow god, know ourselves, and do nothing to excess — what force these have, without natural philosophy, no one can see; and they have the greatest force. Furthermore, what power nature has toward the cultivating of justice, toward the keeping of friendships and the rest of our affections — this knowledge alone can teach. Nor indeed can piety toward the gods, nor how great a gratitude is owed to them, be understood without an unfolding of nature.

\ciceroSection{74} But by now I feel that I have been carried further than the plan I set out would require. The truth is, the admirable composition of the system and the incredible order of its matter has drawn me on. For — by the immortal gods! — do you not marvel at it? What, after all, can be found, whether in nature, than which nothing is more fitting, nothing more precisely laid out, or in the works made by hand, so well composed, so closely framed and fitted together? What that comes after fails to accord with what came before? What follows that does not answer to what went above? What is not so linked, each thing from another, that, if you should move a single letter, the whole would totter? And yet there is nothing that can be moved.

\ciceroSection{75} How weighty, indeed, how magnificent, how steadfast is the character of the wise man that is brought to completion! For he, once reason has taught him that what is honourable is alone the good, must of necessity be happy forever, and truly possess all those titles that the unlearned are wont to mock. For he will more rightly be called a king than Tarquinius, who could rule neither himself nor his own; more rightly the master of the people — for that is what the dictator is — than Sulla, who was the master of three ruinous vices, extravagance, greed, and cruelty; more rightly rich than Crassus, who, had he not been in want, would never have wished to cross the Euphrates with no cause for war. All things will rightly be said to be his, who alone knows how to use all things; rightly too he will be called beautiful — for the features of the mind are more beautiful than those of the body; rightly the only free man, obeying no one's domination and bending to no desire; rightly unconquered, since, even if his body should be bound fast, no chains at all can be cast upon his mind; and he need not wait for any season of life, that only then at last he may be judged whether he has been happy, when he has finished the last day of his life by death — which is what that one of the Seven Sages, with no wisdom, counselled to Croesus.

\ciceroSection{76} For if Croesus had ever been happy, he would have carried his happy life unbroken all the way to that pyre heaped up by Cyrus. But if it is so — that no one save the good man is happy, and that all good men are happy — then what is more to be cultivated than philosophy, and what is more divine than virtue?

\ciceroBook{4}

\ciceroSection{1} When he had said this, he made an end. But I said: 'Indeed, Cato, you have set those matters out so well — so much of it from memory, and matters so obscure, yet so lucidly — that we must either give up wanting to say anything against it at all, or else take some time to think it over. For a doctrine so carefully built — not merely founded, but raised up storey by storey — is no easy thing to master to the bottom, even if it should be less than true (and that I do not yet venture to say).' Then he said: 'Do you really mean it? When I see you, under this new law, answering the prosecutor on the same day and finishing your whole case in three hours, do you suppose that in this present matter I will let you put the time off? And yet your case here will be argued no better than those you sometimes win. So set about this one too — all the more since it has often been handled, both by others and by you yourself, so that words can scarcely fail you.'

\ciceroSection{2} Then I said: 'By Hercules, I am not in the habit of arguing rashly against the Stoics — not because I altogether agree with them, but because shame holds me back; for they say so much that I can scarcely understand.' 'Some of it, I admit, is obscure,' he said; 'yet it is not put obscurely by them on purpose — the obscurity is in the things themselves.' 'Then why is it,' I said, 'that when the Peripatetics say the very same things, there is not a word that fails to be understood?' 'The very same things?' he said. 'Or have I argued too feebly that the Stoics differ from the Peripatetics not in words, but in the whole matter and the entire substance of their view?' 'Well, Cato,' I said, 'if you make that good, you may carry me over to your side, all of me.' 'For my part,' he said, 'I thought I had said enough. So address yourself to those points first, if you please; or, if you would rather have something else, after that.' 'No,' I said, 'I shall take each thing in the place where it stands — by my own judgment, unless that is an unfair demand.' 'As you please,' he said; 'though that other way were more fitting — to grant each man his own choice.'

\ciceroSection{3} 'I hold, then, Cato,' I said, 'that those old pupils of Plato — Speusippus, Aristotle, Xenocrates, and after them Polemo and Theophrastus — had set up their doctrine fully enough, and richly enough, and with enough elegance, so that there was no reason for Zeno, once he had heard Polemo, to part company both with Polemo himself and with his predecessors. Their teaching ran as follows; and in it I should be glad if you would mark what you think needs changing, without waiting for me to deal with everything you said. For I think the whole reasoning of their school must be set against the whole of yours, and joined in battle entire.

\ciceroSection{4} They saw that we are so born as to be fitted, in common, for those virtues that are well known and illustrious — I mean justice, temperance, and the rest of the same kind, all of which resemble the other arts and differ only in their material and in the handling of it for the better; they saw, too, that we reach after these very virtues with a higher spirit and a keener ardour, and that we have implanted in us — or rather inborn — a certain longing for knowledge, and that we are born for the gathering of men together and for the partnership and community of the human race, and that these things shine out most in the greatest natures. And so they divided the whole of philosophy into three parts — a partition that we see has been kept by Zeno.

\ciceroSection{5} Of these, since one is the part by which conduct is held to be shaped, I postpone that part, which is, so to speak, the very root of this inquiry. What the end of goods may be I will treat presently; here I say only this much: that what we may rightly call the civil part — the Greeks call it the political \textit{[Greek: politikon]} — was handled gravely and copiously by the old Peripatetics and Academics, who, agreeing in substance, differed in their words. How much they wrote about the commonwealth, how much about the laws! How much they left behind not only of precepts in their treatises but of models of fine speaking in their orations! For, in the first place, the very things that had to be argued out with subtlety they expressed in a polished and apt manner, now defining, now dividing, as your school does too; but you do it more squalidly, while you see how their style shines.

\ciceroSection{6} And then the things that called for ornate and weighty speech — how magnificently they said them, how splendidly! On justice, on temperance, on courage, on friendship, on the conduct of one's life, on philosophy, on taking up public affairs — the work of men who did not pluck thorns, as the Stoics do, nor strip bones bare, but wished to say grand things ornately and lesser things clearly. And so what consolations are theirs, what exhortations, what counsels too and pieces of advice written to the greatest of men! For their training in speaking, like the very nature of the subject, was twofold. For whatever is inquired into involves a controversy either of the general kind, apart from persons and times, or of fact, or of right, or of name, with these added. In both, then, they trained themselves, and that discipline gave them their great abundance in each kind of speaking.

\ciceroSection{7} This whole field Zeno and his followers either could not maintain or would not — at any rate they abandoned it. To be sure, Cleanthes wrote a treatise on rhetoric, and Chrysippus too; but they wrote in such a way that, if anyone should set his heart on being struck dumb, he need read nothing else. And so you see how they speak. They coin new words; they desert the words in use. Yet what great things they attempt! "This whole world is our city!" You see, then, how they set their hearers on fire. How great a matter you are taking in hand — to make a man who lives at Circeii reckon this whole world to be his own municipality! "What of it? Does he set them ablaze?" He will sooner put them out, if he gets them already burning. Even those very points you stated so briefly — that the wise man alone is king, dictator, rich man — from you, indeed, neatly and roundly; and well he might, for you have it from the rhetoricians; but with them, how meagre are those very claims about the power of virtue! — a power they wish to be so great that it can make a man happy by itself. They prick, as it were, with little stings, with cramped little questions, by which even those who give their assent are not changed at all in mind, and go away the same men who came. For matters perhaps true, certainly weighty, are not handled as they ought to be, but rather more minutely than they deserve.

\ciceroSection{8} Next comes the method of reasoning and the knowledge of nature; for of the highest good we shall see presently, as I said, and shall bring our whole discussion to bear on setting it out. In these two parts, then, there was nothing that Zeno was eager to alter. For the matter stood splendidly, and that in each part. For what, among the things that bear on reasoning, was passed over by the ancients? They both defined a great many things and left behind arts of definition; and what is attached to definition — that a thing be divided into its parts — this both is done by them and is handed down as it ought to be done; the like holds for contraries, from which they came to genera and to the forms of genera. As for the method of drawing a conclusion by argument, they make the starting-point those things they call self-evident; then they follow an order; and last, the final conclusion, of what is true in each case.

\ciceroSection{9} And how great a variety of arguments concluding by reason there is among them, and how unlike these are to captious questions! And what of this — that in very many places they all but proclaim that we must look for the trustworthiness neither of the senses without reason nor of reason without the senses, and that we must not separate the one from the other? And what of this — that the things the dialecticians now hand down and teach, were they not established or discovered by those men? On these subjects, though Chrysippus laboured most, still Zeno did far less than the ancients; and of this man's work, some things were done no better than by the old, and some left undone altogether.

\ciceroSection{10} And whereas there are two arts by which reasoning and discourse are made complete — one of finding, the other of arguing — this latter both the Stoics and the Peripatetics handed down, but the former those men handed down superbly, while these did not so much as touch it. For the places from which arguments might be drawn forth, as from treasuries, your people did not so much as suspect, whereas their predecessors handed them down by craft and by method. And this is just what makes it unnecessary always to chant out the same things by rote, so to speak, and never to part from one's own little notebooks. For he who knows where each thing is laid up, and by what road one comes to it, will be able, even if something has been buried, to dig it out, and always to be his own master in argument. And though certain men endowed with great gifts attain abundance in speaking without method, still art is a surer guide than nature. For it is one thing to pour out words in the manner of the poets, another to distinguish what you say by reason and by art.

\ciceroSection{11} The like may be said of the unfolding of nature, which both these men and yours make use of — and not merely for the two reasons that Epicurus thinks, that the fear of death and of religion may be driven off, but also because the knowledge of things heavenly brings a certain modesty to those who see how great is the moderation even among the gods, how great the order; and greatness of soul, to those who behold the works and deeds of the gods; and justice too, once you have come to know what the will is of the supreme ruler and lord, what his purpose, what his pleasure — to whose nature the true and highest law is said by the philosophers to be conformed in its reason.

\ciceroSection{12} There is, in this same unfolding of nature, a certain insatiable pleasure drawn from coming to know things — the one pursuit in which, once the necessary business is dispatched and we are free of affairs, we may live honourably and like free men. Therefore in this whole reckoning, on nearly the greatest matters, the Stoics followed those men: in saying that the gods exist, and that all things are composed of four elements. But when the question was raised — a matter of no small difficulty — whether there might be some fifth nature, out of which reason and intelligence should arise (in which connection it was asked also of what kind the soul is), Zeno said that this was fire; then some things otherwise, but only a few; on the greatest matter, however, in the same way — that the universe and its greatest parts are administered by a divine mind and nature. But the material of things, and the abundance of it, we shall find meagre in these men, most rich in those.

\ciceroSection{13} How much was searched out and gathered by them about every kind of living creature — its origin, its members, its ages! How much about the things that are brought forth from the earth! How many, and on how many various matters, both causes — why each thing comes to be — and demonstrations of the manner in which each thing comes to be! And from all this abundance the most numerous and most certain arguments are drawn for setting out the nature of each thing. So far, then, so far as I at least understand, there seems to have been no cause for changing the name. For it does not follow that, because he did not follow them in everything, his school did not for that reason take its rise from there. Indeed, I hold that even Epicurus, in physics at any rate, is a Democritean. He changes a few things — or many, granted; but when he says the same as the other on most points, then surely he does so on the greatest. And since your people do the very same, they do not grant gratitude enough to the discoverers.

\ciceroSection{14} But enough of this. Now let us see, I beg, what at last he brought to bear on the highest good, which is the keystone of philosophy — what reason he had for dissenting from the discoverers as from his own parents. On this point, then — although it was carefully set out by you, Cato, this end of goods which is the keystone of philosophy, and how, and in what terms, it was stated by the Stoics — still I too shall expound it, so that we may discern, if we can, what at all of anything new was brought in by Zeno. For when his predecessors, of whom most plainly Polemo, had said that the highest good is to live in agreement with nature, the Stoics say that three things are signified by these words. The first is of this kind: to live applying knowledge of the things that come about by nature. They say that this very thing is Zeno's end — the thing meant when you said "to live consistently with nature."

\ciceroSection{15} The second is that the same thing is signified, as if it were said: to live keeping all the intermediate duties, or most of them. This, set out thus, is unlike the former. For the former is right — what you used to call \textit{[Greek: katorthōma]} — and falls to the wise man alone; whereas this is a matter of a certain inchoate duty, not a perfected one, and can fall to some who are not wise. The third, again, is to live enjoying all, or the greatest, of those things that are in keeping with nature. This is not placed within our own action; for it is made complete both out of that kind of life which has the enjoyment of virtue, and out of those things which are in keeping with nature but are not within our power. But this highest good, which is understood in the third signification, and that life which is led from the highest good — because virtue is conjoined to it — falls to the wise man alone; and that end of goods, as we see it written by the Stoics themselves, was established by Xenocrates and by Aristotle. And so by them that first constitution of nature, from which you too set out, is expounded in nearly these words:

\ciceroSection{16} Every nature wishes to be the preserver of itself, so as to be both safe and kept in being according to its own kind. To this end, they say, the arts too were sought out, to help nature; and among these is reckoned, first of all, the art of living, so as to guard what has been given by nature and to acquire what is lacking. And the same men divided the nature of man into mind and body. And since they had said that each of these is to be sought for its own sake, they said too that the virtues of each are to be sought for their own sakes; and since they set the mind, with a certain boundless praise, before the body, they set the virtues of the mind too before the goods of the body.

\ciceroSection{17} But since they would have wisdom be the guardian and steward of the whole human being, as nature's companion and helper, they said that this was the office of wisdom: that, in watching over a creature compounded of mind and body, it should aid and hold him together in both. And so, after the matter had been laid down thus simply at the outset, they pursued the rest with greater subtlety, and judged that the goods of the body were governed by some easy account; but the goods of the mind they searched out more exactly, and first of all they found that the seeds of justice were present in them, and they were the first of all the philosophers to teach that it was granted by nature that those who had been begotten should be loved by their begetters; and — what in the order of time is more ancient still — that by nature the unions of husbands and wives were joined together, the stock from which spring the friendships of kinsfolk. And setting out from these beginnings, they traced both the origin and the growth of all the virtues. From this there arose, too, greatness of mind, by which one could readily resist fortune and stand against it, since the greatest things lay within the wise man's power. The shifts and the wrongs of fortune, moreover, a life ordered by the precepts of the old philosophers readily overcame.

\ciceroSection{18} And from the beginnings given by nature there were kindled certain expansive growths of goods, some springing from the contemplation of more hidden things, since there was implanted in the mind a love of knowledge, and from this there followed also a desire of unfolding reason and of arguing; and since this is the only creature born with a share of modesty and a sense of shame, longing for the fellowship of human beings and taking heed, in all that it did or said, that nothing be done by it save honourably and becomingly — from these beginnings, as I said before, and from these seeds given by nature, temperance, moderation, justice, and all that is honourable were brought to full completion.

\ciceroSection{19} There you have, Cato, the form of those philosophers of whom I speak. And now that it has been set out, I am eager to know what the reason is why Zeno broke away from this ancient constitution, and which of these tenets was not approved by him. Was it that they said all nature is the preserver of itself? Or that every creature is commended to itself, so that it should wish to be both safe in its kind and unharmed? Or that, since the end of all the arts was that which nature most sought, the same had to be laid down for the art of life as a whole? Or that, since we are compounded of mind and body, both these themselves and their virtues were to be taken up for their own sakes? Or did so great a pre-eminence assigned to the virtues of the mind displease him? Or what is said of prudence, of the knowledge of things, of the bond of the human race; or what is said by these same men of temperance, of moderation, of greatness of mind, of all that is honourable? The Stoics will confess that all these things were nobly said, and that this was not Zeno's reason for breaking away.

\ciceroSection{20} They will allege, I suppose, certain other errors of the ancients, great ones, which that man, eager to track down the truth, could in no way endure. For what is more perverse, what more intolerable, what more foolish, than to place good health, freedom from all pains, the soundness of the eyes and the rest of the senses among goods, rather than to say that there is no difference whatever between those things and their contraries? For all those things which the ancients called goods, he said, are ``preferred,'' not goods; and likewise that the ancients spoke foolishly when they said that those things which excel in the body are to be sought for their own sakes — they are to be taken up rather than sought for their own sakes. And finally that a whole life consisting in virtue alone is not more to be sought than that life which abounds also in the other things that are in keeping with nature, but rather more to be taken up. And though virtue itself makes life so happy that it cannot be happier, still certain things are wanting to the wise even when they are happiest; and so they busy themselves with this, to drive off from themselves pains, diseases, and infirmities.

\ciceroSection{21} O the great force of genius, and the just cause, why a new school should arise! Go on further. For there follow those things which you embraced so learnedly: that the folly of all men, their injustice, and the other like vices are equal, and that all wrongs are equal; and that those who by nature and instruction had advanced far toward virtue, unless they had plainly attained it, are utterly wretched, and that between their life and the life of the most depraved there is no difference at all — so that Plato, that so great man, if he was not wise, lived no better than any utter scoundrel, nor more happily. This, you see, is the correction and the amendment of the old philosophy — a correction that can find no entrance at all into the city, into the Forum, into the Senate house. For who could endure a man speaking thus, one who professed himself an authority on living gravely and wisely, while he changed the names of things, and, though he held the same view as everyone else, fastened other names upon the things to which he assigned the same force, altering the words only, withdrawing nothing from the substance of the opinions?

\ciceroSection{22} Would an advocate, pleading for his client in the peroration, deny that exile, that confiscation of goods, is an evil? That these things are to be rejected, not fled from? That a juror ought not to feel pity? And if he were speaking in the assembly — supposing Hannibal had come to the gates and hurled a javelin over the wall — would he deny that to be captured, to be sold, to be killed, to lose one's country, is among evils? Or could the Senate, when it decreed a triumph to Africanus, speak of ``his virtue'' or ``his good fortune,'' if neither virtue can be truly ascribed to any but the wise man, nor good fortune? What sort of philosophy is this, then, which speaks in the common way in the Forum, but in its own way in its pamphlets? — especially since, in what those men signify by their own words, nothing is made new, nothing is changed about the things themselves, but the same things remain in a different guise.

\ciceroSection{23} For what difference does it make whether you call riches, resources, health goods, or ``preferred,'' when the man who calls them goods assigns no more to them than you do, who name those same things ``preferred''? And so Panaetius — a man pre-eminently noble and grave, worthy of that intimacy with Scipio and Laelius — when he wrote to Q. Tubero on the bearing of pain, nowhere set down what ought to have been the chief point, if it could be made good: that pain is not an evil; but rather what it was and of what kind, how much there was in it that was alien to us, and then what the method of enduring it was. And since he was a Stoic, his verdict, it seems to me, condemns that emptiness of words.

\ciceroSection{24} But to come closer, Cato, to the things you have said, let us proceed more pressingly, and let us compare what you have just said with the things I set before yours. The points, then, which you hold in common with the ancients, let us use as conceded; the points that come into dispute, on these, if you please, let us argue. ``For my part,'' he said, ``I am content that the matter be handled with more subtlety and, as you yourself put it, more pressingly. For what you have brought forward so far is of the popular sort, but from you I look for something more refined.'' ``From me?'' I said. ``Still, I will strive; and if not very many points occur to me, I will not shy away from those popular ones.''

\ciceroSection{25} But first let it be laid down that we ourselves are commended to ourselves, and that we have this first appetite from nature, that we preserve ourselves. On this we agree; the next thing follows, that we take heed of what we ourselves are, so that we may keep ourselves such as we ought to be. We are, then, human beings. We are compounded of mind and body, which are of a certain sort, and we ought, as the first natural appetite requires, to cherish these and to establish out of them that end of the highest and uttermost good. And this, if the first principles are true, must be established thus: to attain as many and as great as possible of those things that are in keeping with nature.

\ciceroSection{26} This end, then, they held; and what I have put in more words, they put more briefly — to live in keeping with nature: this seemed to them the uttermost of goods. Come now, let those men teach — or rather you, for who could do it better? — in what way, setting out from those same first principles, you bring it about that to live honourably (that is, to live by virtue, or in agreement with nature) is the highest good, and in what way, or at what point, you suddenly abandoned the body and all those things which, though they are in keeping with nature, are not within our power — and finally duty itself. I ask, then, how these so great commendations, springing from nature, were suddenly forsaken by wisdom.

\ciceroSection{27} But if we were seeking the highest good not of a human being, but of some living creature, and that creature were nothing but a mind — for let it be allowed to imagine something of this sort, that we may more easily find the truth — even so, this end of yours would not be the end for that mind. For it would still long for health, for freedom from pain, and would seek too the preservation of itself and the safeguarding of those things, and would set up for itself as its end to live in keeping with nature, which is, as I said, to have those things that are in keeping with nature, either all of them or the most and the greatest.

\ciceroSection{28} For of whatever sort you should frame a living creature, it must needs be that, even if it were without a body, as we are imagining, there should still be in the mind certain things resembling those that are in the body, so that in no way could the end of goods be established except as I have set it out. Chrysippus, however, in setting out the differences among living creatures, says that some of them excel in body, others in mind, and some have strength in both; then he disputes what end befits each kind of creature to have laid down for it. And when he had placed the human being in that kind, so as to assign to him excellence of mind, he established the highest good in such a way that it seemed not that the mind excels, but that there is nothing besides the mind. But in one way alone could the highest good be rightly placed in virtue alone: if there were some creature consisting wholly of mind, and that mind, even so, were such as to have in itself nothing that was in keeping with nature, as health is.

\ciceroSection{29} But it cannot even be thought what such a thing would be like, without its being at war with itself. But if he says that certain things are obscured and do not appear because they are very small, we too concede this; Epicurus says it of pleasure as well, that the smallest pleasures are often obscured and overwhelmed. But the conveniences of the body are not, in that respect, so great, nor so prolonged in their seasons, nor so many. And so, in those cases where, on account of their slightness, an obscuring follows, it often happens that we confess it makes no difference to us whether they exist or not — as, in the case of the sun, which you spoke of, it makes no difference to bring a lamp to it, or to add a farthing to the wealth of Croesus.

\ciceroSection{30} But in those matters where no such obscuring takes place, it may still be that the very thing that makes a difference is not great. As, to a man who has lived ten years pleasantly, if a month of equally pleasant life be added, since the addition has some weight toward pleasantness, it is a good; but if this be not granted, the happy life is not on that account at once destroyed. The goods of the body, however, are more like this latter case which I put; for they have an addition worth labouring for — so that the Stoics seem to me at times to be making a joke of this, when they say that, if to that life which is passed with virtue there should be added a flask or a strigil, the wise man would take up rather the life to which these things had been added, yet would not on that account be the happier.

\ciceroSection{31} Is this, after all, a fitting comparison? Is it not to be cast out with laughter rather than with argument? For whether the flask exists or not, who would not most justly be laughed at if he troubled himself over it? But if someone should relieve another of the deformity of his limbs and the torment of pains, he would earn great gratitude; nor, if that wise man were forced by a tyrant to go to the torturer's rack, would he wear the same face as if he had lost a flask, but, entering as it were upon a great and difficult contest, when he saw that he had to fight it out to the end with a deadly adversary, pain, he would rouse all the resources of fortitude and endurance, under whose protection he might enter that difficult and, as I said, great battle. Then, too, we are not asking what is obscured or perishes because it is exceedingly small, but what is of such a kind as to fill out the sum. One pleasure out of many is obscured in that life of pleasure, yet, however small it be, it is a part of that life which is set in pleasure. A coin is obscured in the riches of Croesus, yet it is a part of the riches. Wherefore let these things too, which we say are in keeping with nature, be obscured in the happy life — provided only they be parts of the happy life.

\ciceroSection{32} And yet if, as ought to be agreed between us, there is a certain natural appetite seeking after those things that are in keeping with nature, then of all those things some sum total must be made. And this once established, it will then be possible to inquire at leisure into those matters — into the magnitude of things, into their excellence, how much there is in each toward living happily, into those very obscurings which, on account of their slightness, scarcely appear, or scarcely even that. What of the point on which there is no disagreement? For there is no one who has said otherwise than that the thing to which all things are referred — which is the uttermost of things to be sought — is alike for all natures. For every nature is cherishing of itself. For what nature is there that ever forsakes itself, or any part of itself, or the disposition of that part, or its force, or any of those things that are in keeping with nature, or its motion, or its state? And what nature is forgetful of its own first constitution? Surely there is none that does not keep its own force from first to last. How, then, did it come about that human nature alone should be the one to abandon the human being, to forget the body, to place the highest good not in the whole human being, but in a part of the human being?

\ciceroSection{33} But how, on their own admission and by everyone's agreement, will it be kept true that this final good in question is the same for every nature? For it would be the same only if, in the other natures too, the final good for each were that in which each excelled. That, after all, is how the Stoics' final good is conceived.

\ciceroSection{34} Why, then, do you hesitate to change your starting-points in nature? For why do you say that every living creature, the moment it is born, is fastened to loving itself and is taken up with preserving itself? Why do you not rather say that every living creature is fastened to that which is best in it, and is taken up with guarding that one thing alone, and that the other natures do nothing else but preserve whatever is best in each? But how can a thing be best, if there is no good besides it? And if, on the other hand, the rest are to be desired, why is the final good among things to be desired not arrived at from the desiring of all of them, or of the greatest and the most? Phidias can set up a statue from the beginning and bring it to completion, and he can also take one begun by another and finish it; wisdom is like the latter. For it did not itself give birth to man; it received him from nature already begun. With its eyes fixed on her, then, it ought to finish the statue, so to speak, that she set up. What sort of man, therefore, did nature begin?

\ciceroSection{35} And what is the task, what is the work of wisdom? What is there that ought to be finished and perfected by it? If there is nothing in that which is to be perfected except a certain motion of the mind, that is, reason, then its final good must necessarily be to act in accordance with virtue; for virtue is the perfecting of reason. If there is nothing but body, the supreme goods will be those things: health, freedom from pain, beauty, and the rest.

\ciceroSection{36} As it is, the inquiry concerns the highest good of man; why, then, do we hesitate to inquire into his whole nature, to see what is brought about there? For though it is agreed among all that every duty and task of wisdom is taken up with the cultivation of man, some thinkers — and so that you may not suppose I am arguing against the Stoics alone — bring forward views that place the highest good in a kind of thing beyond our power, as though they were speaking of something inanimate; others, on the contrary, as if man had no body at all, care for nothing besides the mind — and this although the mind itself is not some empty thing, I know not what (for I cannot understand that), but belongs to a certain kind of body, so that it too is not content with virtue alone, but craves freedom from pain. And so both parties do the same thing as a man would who, neglecting the left side, looked after the right; or who, like Erillus, embraced the mind's act of knowing and abandoned its acting. For of all those who pass over much, while they single out some one thing to follow, the view is, so to speak, cut short; whereas that view is whole and complete which belongs to those who, when they inquired into man's highest good, left no part of him, of mind or of body, without its guardianship.

\ciceroSection{37} But you, Cato — because virtue, as we all admit, holds the highest and most pre-eminent place in man, and because we count those who are wise to be finished and perfect — you dazzle the keenness of our minds with the splendour of virtue. For in every living thing there is something supreme and best — in horses, in dogs, which nonetheless need both to be free of pain and to be in good health; so too, then, in man this perfection of yours is praised above all in that which is best, that is, in virtue. And so you seem to me not to consider sufficiently what the path of nature is, and what her progress. For what she does with crops — when she has brought them from the blade to the ear, she leaves the blade behind and counts it as nothing — she does not do the same with man, when she has brought him to a settled state of reason. For she always takes on something in such a way that she does not abandon the things she gave first.

\ciceroSection{38} And so she joined reason to the senses, and, once reason was achieved, did not leave the senses behind. It is as if the tending of vines, whose task is to bring it about that the vine, with all its parts, should be in the best possible condition — but let us understand it thus (for it is permitted, as you too are accustomed to do, to invent something for the sake of teaching): if, then, that tending of vines were present in the vine itself, it would, I suppose, want the other things that bear on the cultivation of the vine, just as before, but would set itself ahead of all the parts of the vine and would judge that nothing in the vine is better than itself. In like manner sense, when it has come to nature, watches over her, indeed, but watches over itself as well; whereas when reason has been taken on, it is set in such a position of mastery that all those first things of nature are made subject to its guardianship.

\ciceroSection{39} And so it does not withdraw from the care of those things over which it has been set, and which it ought to govern throughout the whole of life — so that I cannot sufficiently marvel at the inconsistency of these men. For natural desire, which they call \textit{[Greek: hormē]}, and likewise duty, and virtue itself, they will have to be guardians of the things that are in accordance with nature. But when they want to reach the highest good, they leap over everything and leave us with two tasks instead of one, so that we take up some things and seek after others, rather than enclosing both within a single end.

\ciceroSection{40} But now you say that virtue cannot be established if the things that lie outside virtue bear upon living happily. Yet the whole truth is the contrary. For virtue can in no way be brought in unless everything it chooses and everything it rejects is referred to a single sum. For if we neglect them altogether, we shall fall into the faults and errors of Aristo, and shall forget the very principles we have given to virtue itself; whereas if we do not neglect them, and yet do not refer them to the end of the highest good, we shall stray not far from Erillus's frivolity. For we should then have to take up the rules of two lives. For he makes two separate final goods, which, to be true, ought to have been joined together; as it is, they are so set apart as to be disjoined — and nothing can be more perverse than that.

\ciceroSection{41} And so the contrary is true of what you say; for virtue can in no way be established unless it holds the first things of nature as bearing upon the sum. For virtue was sought as something that would not abandon nature but would watch over her. Yet virtue, on your account, watches over a certain part and deserts the rest. And if the very making of man could speak, it would say this: that its first beginnings of desire, so to speak, were aimed at preserving itself in the nature in which it had come to be. But it had not yet been made clear enough what nature most wanted. Let it be made clear, then. What else, then, will be understood but that no part of nature is to be neglected? And if there is nothing in it besides reason, let the end of goods lie in virtue alone; but if there is body too, then this clarification of nature will surely have brought it about that we abandon the very things we held before the clarification. To live in agreement with nature, then, on this showing, is to depart from nature.

\ciceroSection{42} As certain philosophers, when, setting out from the senses, they had glimpsed some greater and more divine things, abandoned the senses, so these men, when out of the desire for things they had caught sight of the beauty of virtue, threw away everything they had seen apart from virtue itself — forgetting that the whole nature of things to be desired ranges so widely that it runs through from the beginnings to the ends, and not understanding that they are pulling away the very foundations of those beautiful and admirable things.

\ciceroSection{43} And so it seems to me that all those did indeed err who said that the end of goods is to live honourably, but that one erred more than another — Pyrrho, of course, most of all, who, once virtue was established, left nothing whatever that might be desired; then Aristo, who did not dare to leave nothing, but brought in, as the things by which the wise man is moved to desire something, whatever might fall into his mind and whatever might present itself by chance. He is better than Pyrrho in that he allowed some kind of desiring, worse than the rest in that he withdrew utterly from nature. The Stoics, in placing the end of goods in virtue alone, are like those men; in seeking a starting-point of duty, they are better than Pyrrho; in not making the objects of desire matters of chance occurrence, they outdo Aristo; but in not joining to the end of goods the things they say are suited to nature and to be taken up for their own sake, they fall away from nature and are in a certain way not unlike Aristo. For he was inventing chance occurrences, I know not what; while these men do indeed lay down the first things of nature, but sever them from the ends and from the sum of goods. When they give them preference, so that there may be some selection among things, they seem to follow nature; but when they deny that these things bear at all upon the happy life, they once more abandon nature.

\ciceroSection{44} And up to now I have stated the case for there being no reason why Zeno should have departed from the authority of his predecessors. Now let us look at the rest, unless either you wish to say something to this, Cato, or we are already running too long. Neither, surely, said he. For I want this discussion to be carried through by you, and your speech cannot seem long to me. Splendid, said I. For what could be more to my wish than to argue about the virtues with Cato, the patron of every virtue?

\ciceroSection{45} But first observe this: that gravest tenet of yours, which leads the household — that what is honourable is the only good, and that to live honourably is the end of goods — will be shared by you with all who place the end of goods in virtue alone; and what you say, that virtue cannot be given shape if anything is reckoned a good but the honourable, will be said also by those whom I have just named. To me, however, it seemed fairer to see Zeno, when he disputed with Polemo, from whom he had received what the first things of nature are, set out from common starting-points and so find where he first took his stand and whence the cause of the controversy arose, rather than to see him stand alongside men who would not so much as say that their highest goods sprang from nature, and use the same arguments they used and the same opinions.

\ciceroSection{46} But least of all do I approve of this: that, when you have taught — as you suppose — that the only good is what is honourable, you then turn round and say that there must be set before us starting-points fitted and suited to nature, out of the selection of which virtue can come to be. For virtue ought not to have been placed in the selection, so that the very thing which was the final good should acquire something else besides. For everything that is to be taken up and everything that is to be chosen or wished for ought to be contained within the sum of goods, so that one who has attained it desires nothing further. Do you see how, for those whose supreme good lies in pleasure, it is plain what they must do and not do? How no one is in doubt about the aim of all their duties, what they ought to follow, what to flee? Let this final good be what I now defend; at once it is clear what the duties are, what the actions. But you, who have nothing else set before you but the right and the honourable — you will not find out whence the starting-point of duty, whence the starting-point of action, is to arise.

\ciceroSection{47} In seeking this, then, you will all come back to nature: both those who will say that they follow whatever comes into their mind or whatever presents itself, and you yourselves. To them nature would justly reply that it is not true that the end of living happily is sought from one source, and the starting-points of conduct from herself; for there is a single account by which both the starting-points of things to be done and the final goods are contained. And just as Aristo's view was hooted off the stage when he said that one thing differs in no way from another, and that there are no things at all, besides the virtues and the vices, between which there is any difference whatever, so Zeno errs when he says that in nothing except virtue or vice is there an inclination of even the slightest weight toward attaining the highest good, and yet says that, although the other things have no weight toward the happy life, there is nonetheless in them weight toward the desiring of things — as though, indeed, this desiring did not bear upon the attainment of the highest good! But what is less consistent than what they say: that, once the highest good is known, they go back to nature in order to seek from her the starting-point of action, that is, of duty?

\ciceroSection{48} For the account of action or of duty does not drive us to seek the things that are in accordance with nature; rather, both desire and action are set in motion by those things. Now I come to those brief points of yours, which you called corollaries — and first to that than which nothing can be briefer: Every good is praiseworthy; what is praiseworthy is honourable; therefore every good is honourable. What a leaden dagger! For who would grant you that first premise? — and once it is granted, there is no need of the second; for if every good is praiseworthy,

\ciceroSection{49} everything honourable is good—but who, then, will grant you that premise, except Pyrrho, Aristo, or men of their stamp, whom you do not approve? Aristotle, Xenocrates, that whole household will not grant it, seeing that they call health, strength, riches, glory, and many other things goods, though they do not call them praiseworthy. And these men, indeed, while they hold that the end of goods is not bounded by virtue alone, nonetheless set virtue before all other things; what, then, do you suppose those will do who have cut virtue off altogether from the end of goods—Epicurus, Hieronymus, and those too, if there be any, who wish to defend the end Carneades proposed?

\ciceroSection{50} And again, how will either Callipho or Diodorus be able to concede that to you, when they join to honour something else that is not of the same kind? Are you pleased, then, Cato, once you have taken up premises that are not granted, to draw from them whatever you like? There is, besides, that sorites of yours, than which you Stoics think nothing more vicious: what is good is to be wished for; what is to be wished for is to be sought; what is to be sought is praiseworthy—and then the remaining steps. But here I plant my feet; for in the same way no one will grant you that what is to be sought is praiseworthy. As for that other conclusion, it is anything but a proper inference—it is positively dull, and it belongs, of course, to them and not to you—that the happy life is worthy of being boasted of, since without honour it cannot come about that anyone should rightly boast.

\ciceroSection{51} Polemo will grant this to Zeno, and so will Polemo's own master, and that whole line, and the rest who, while setting virtue far above all other things, nonetheless join something to it in fixing the highest good. For if virtue is worthy of being boasted of, as it is, and so far excels all other things that it can scarcely be told, then a man endowed with virtue alone, lacking everything else, will be able to be happy—and yet you will not on that account be granted that nothing but virtue is to be reckoned among goods. Those, on the other hand, for whom the highest good lies outside virtue will perhaps not grant that the happy life contains anything in which a man might rightly boast—though they do at times make even pleasures a ground for boasting.

\ciceroSection{52} You see, then, that you are taking up either premises that are not granted, or premises that, even when granted, do you no good. For my part, in all these chains of reasoning I should have thought this worthy of philosophy and of us—especially when we were inquiring after the highest good—that our life, our counsels, our wills be set right, not our words. For who can give up his conviction on hearing those terse and pointed arguments that delight you, as you say? When men are waiting and longing to hear why pain is not an evil, the Stoics tell them that to feel pain is harsh, troublesome, hateful, against nature, hard to bear—but that, because there is in pain no fraud, no wickedness, no malice, no fault, no baseness, it is therefore not an evil. Whoever has heard this will go away no firmer for the bearing of pain than he came, even if he forbears to laugh.

\ciceroSection{53} You, however, deny that anyone can be brave who thinks pain an evil. Why should he be braver if he is going to think it harsh and scarcely bearable—the very thing you yourself concede? For timidity is born of things, not of words. And you say that if a single letter be disturbed, the whole system will totter. Do I seem to you, then, to be disturbing a single letter, or whole pages? For granted that among them there is what you praised—the order of things kept intact, and everything fitted and linked together (for so you used to say)—still, we are not bound to follow it through, if, starting from false first principles, it agrees with itself and does not stray from its own purpose.

\ciceroSection{54} In his very first founding of the system, then, your Zeno departed from nature; and when he had placed the highest good in the excellence of the mind, which we call virtue, and had said that nothing else was good except what was honourable, and that virtue could not hold firm if among the other things there were anything that was better or worse than another—these premises laid down, he held to the consequences throughout. You say rightly; I cannot deny it. But the things that follow are so false that the things from which they were begotten cannot be true.

\ciceroSection{55} For the logicians teach us, as you know, that if the things which follow from some proposition are false, then that very proposition which they follow is false. So this conclusion comes out not only true but so clear that the logicians think no account of it even needs to be given: if the first, then the second; but not the second; therefore not the first either. Thus, when your consequences are removed, your first principles are removed with them. What, then, are the consequences? That all who are not wise are equally wretched; that the wise are all supremely happy; that all right acts are equal, all sins on a par—claims which, though they seemed at first to be magnificently stated, on examination won less approval. For each man's own sense, and the nature of things, and truth herself cried out, in a manner of speaking, that they could not be brought to believe there was no difference among the things that Zeno levelled.

\ciceroSection{56} Afterwards that little Phoenician of yours—for you know that the men of Citium, your clients, set out from Phoenicia—a shrewd fellow, then, finding that his case would not hold because nature fought against it, set about twisting words; and first, with respect to those things which we count goods, he conceded that they might be held to have value and to be suited to nature, and he began to admit that for the wise man—that is, for the supremely happy man—it is nonetheless more convenient if he has those things too which he does not dare to call goods, conceding that they are suited to nature; and he denies that Plato, even if he be not wise, stands in the same case as the tyrant Dionysius; for Dionysius it is best to die, because of his despair of wisdom, while for Plato it is best to live, because of his hope of it. As for sins, some, he says, are tolerable, others in no way so, because some sins overstep more, and others fewer, of the so-called numbers of duty. And again, some of the unwise are so constituted that they can in no way attain wisdom, while others could attain it if they set themselves to the task.

\ciceroSection{57} Here he spoke otherwise than all the rest, yet thought the same as the others. And in fact he reckoned the things he himself denied to be goods to be of no less value than those did who called them goods. What, then, did he mean to gain by changing the names? He should at least have taken away something from their weight, and valued them a little less than the Peripatetics do, so that he might seem to feel differently as well, and not merely to speak differently. What of this? Of the happy life itself, to which everything you say is referred—you deny that it is the life filled out with all those things that nature desires, and you place the whole of it in virtue alone; and since every controversy is wont to be either about a thing or about a name, either kind arises if either the thing is not understood or there is error about the name. If neither is the case, then care must be taken that we use words as customary as possible and as fitting as possible—that is, words that declare the thing.

\ciceroSection{58} Is there any doubt, then, that if no error is made in the thing itself by our predecessors, they use the more convenient words? Let us look, then, at their views, and afterwards come back to the words. They say that the mind's appetite is stirred when something appears to it to be in keeping with nature; that everything in keeping with nature is worthy of some valuation, and is to be valued in proportion to how much weight there is in each; and that, of the things in keeping with nature, some have in themselves nothing of that appetite of which we have already so often spoken—the things that are called neither honourable nor praiseworthy—while others possess pleasure in every living creature, but in a human being reason as well. Of these, the ones that pertain to reason are named honourable, beautiful, praiseworthy, while the former are named natural; and these, joined with the honourable things, complete and perfect the happy life.

\ciceroSection{59} But of all those advantages, to which those who call them goods assign no more weight than Zeno does, who denies that they are goods, by far the most excellent, they hold, is what is honourable and praiseworthy. Yet if two honourable courses be set before us, one with health, the other with sickness, there is no doubt to which of them nature herself will lead us. And nonetheless the power of honour is so great, and it so surpasses and excels all other things, that no torments and no rewards can shift it from what it has decreed to be right; and all things that seem hard, difficult, adverse can be trodden under by those virtues with which we are by nature equipped—not, indeed, that those things are easy or to be despised (for then what would there be of greatness in virtue?), but so that we may judge that in them lies no chief part of living happily or otherwise.

\ciceroSection{60} In sum: the very things that Zeno called worthy of valuation, and to be taken up, and suited to nature, those others call goods; and the happy life they call that which consists of the things I have mentioned, either the most numerous or the weightiest of them. Zeno, on the other hand, calls good only that which has its own peculiar character that makes it worth seeking, and calls the happy life that life alone which is passed in company with virtue. If the dispute must be about the thing, I can have no disagreement with you, Cato; for there is nothing on which you feel otherwise than I, provided that, with the words changed, we compare the things themselves. And Zeno did not fail to see this, but he took delight in the magnificence and the glory of the words. For if he felt about what he says just as the words signify, what difference would there be between him and either Pyrrho or Aristo? But if he did not approve of them, what was the point of differing in words from those with whom he agreed in substance?

\ciceroSection{61} And what if those followers of Plato should come to life again, and after them those who were their hearers, and should speak with you thus? `When we listened to you, Marcus Cato—a most zealous student of philosophy, a most just man, an excellent judge, a most scrupulous witness—we wondered what the reason could be that you set the Stoics before us, men who hold about good and evil the very views that Zeno had learned from this Polemo, but who use names which at first sight stir wonder and, once the matter is unfolded, stir laughter. You, however, if those views won your approval, why did you not hold to them in their own proper words? Or if it was their authority that moved you, did you set this nobody of yours before us all and before Plato himself? Especially when you wished to be a leading man in the commonwealth, and could most of all be equipped and furnished by us, with the highest honour to yourself, for its defence. For it was by us that these matters were investigated, by us that they were set down, marked out, and laid down as precepts; and we have written out in full all the kinds, conditions, and changes in the governing of commonwealths, and the laws too, and the institutions and customs of states. As for eloquence—which is the greatest ornament to leading men, and in which we hear that you have very great power—how much you might have added to yourself out of our records!' When they had said this, what indeed would you answer to such men?

\ciceroSection{62} `I would ask you,' he said, `since you had dictated the speech for them, to speak in turn on my behalf—or rather, to give me a little room to answer them—were it not that I would rather listen to you now, and that I shall in any case answer those men at another time, namely the time when I answer you.' And yet, Cato, if you wished to answer truly, this is what had to be said: not that you disapprove of those men, men of such great genius and such great authority, but that you have noticed that what they, by reason of their antiquity, saw too little of, has been clearly perceived by the Stoics; and that on the very same matters the Stoics have both argued more acutely and felt more weightily and more bravely—seeing that, in the first place, they deny that good health is to be sought and say instead that it is to be chosen, and that being well is to be valued not because it is good but because it is to be reckoned at something (and yet at no more than it seems to those who do not hesitate to call it good); but this you could not bear, that those ancients—bearded, as it were, as we are wont to say of our own forefathers—believed that for a man who lived honourably the life would be more to be wished for, and better, and more to be sought, if he also enjoyed good health, a good name, and abundance, than the life of one who, while equally a good man, was in many ways beset, like Ennius' Alcmaeon, \textit{hemmed in by disease,}

\ciceroSection{63} \textit{by exile and by want.} Those ancients, then, in their less acute fashion, think that life more to be wished for, more excellent, more happy; the Stoics, on the other hand, hold only that it is to be put first in the choosing—not that the life is happier, but that it is better suited to nature. And as for the claim that all who are not wise are equally wretched—this, no doubt, the Stoics saw clearly, while it had escaped their predecessors, who did not think that men polluted by crimes and parricides were no whit more wretched than those who, while they lived chastely and uprightly, had not yet attained that perfect wisdom.

\ciceroSection{64} And at this point you were bringing forward as parallels certain comparisons of theirs which are utterly unlike. For who does not know that, if several men wish to come up from the deep, those who are now drawing near to the surface of the water will indeed be nearer to breathing, but can no more breathe than those who are still in the depths? It does not help, then, to advance and make progress in virtue, so as to be any less wretched, before one has reached it—since in the water it does not help; and since puppies on the point of opening their eyes are as blind as those just born, Plato too, on this showing, since he did not yet see wisdom, must have been as blind in mind as Phalaris?

\ciceroSection{65} Those cases are not parallel, Cato, in which, however far you have advanced, the thing you wish to be free of stands in the same condition until you have made your escape. The man does not breathe again before he has surfaced; and puppies are as blind before their sight opens as if they were going to be so forever. These, rather, are the parallel cases: one man's eyesight is dim, another grows feeble in body; both, once treatment is applied, are relieved day by day — the one grows stronger each day, the other comes to see. Like these are all who pursue virtue: they are relieved of their vices, relieved of their errors — unless, perhaps, you suppose that Tiberius Gracchus the father was no happier than the son, when the one strove to steady the commonwealth and the other to overthrow it. And yet that father was not a sage — for who has ever been one, and when, and where, and from what source? — but because he strove for praise and dignity, he had advanced far in virtue.

\ciceroSection{66} Shall I set your grandfather Drusus beside Gaius Gracchus, who was nearly his contemporary? The wounds the one was inflicting on the commonwealth, the other was healing. If there is nothing that makes men so wretched as impiety and crime, so that all the foolish are now wretched — as they surely are — still the man who looks after his country is not equally wretched with the man who longs to see her destroyed. There comes about, then, a great relief from vices in those who have made some progress toward virtue.

\ciceroSection{67} Your people, however, say that progress toward virtue does take place, but deny that any relief from vices does. It is worth considering the argument these acute men use to prove it. Of those arts, he says, whose highest points can grow, the highest point of their contraries too will be able to be increased; but to the highest point of virtue nothing can be added; therefore vices, which are the contraries of the virtues, will not be able to grow either. Which is it, then, in the end — are doubtful matters laid open by clear ones, or are clear ones done away with by doubtful? Now this is clear, that some vices are greater than others; that is doubtful, whether any addition can be made to what you call the highest good. Yet you, when you ought to illuminate the doubtful by the clear, are trying to do away with the clear by the doubtful.

\ciceroSection{68} And so, by the same line of reasoning I used a little before, you will be stuck once more. For if some vices are not greater than others on the ground that nothing can be added even to the end of goods as you frame it, then, since it is clear that vices are not all equal, you must change your end of goods. For we must hold fast to this necessity: that, when some consequence is false, the thing whose consequence it is cannot be true. What, then, is the cause of these tight corners? A boastful display in fixing the highest good. For when it is maintained that the honourable alone is good, care for one's health is abolished, attention to the household estate, the administration of the commonwealth, the ordering of business to be carried on, the duties of life — in the end, that very honourable itself, in which alone you would have everything consist, must be abandoned. These points are made most carefully against Aristo by Chrysippus. It was out of that difficulty that those, as Accius says, \textit{deceit-speaking} cunning devices were born.

\ciceroSection{69} For since wisdom had no place to set its foot once all duties were taken away — and duties were taken away once all choice and distinction were removed, things that could not exist when all things were thus levelled so that nothing differed among them — out of these tight corners that doctrine emerged, worse than Aristo's. His, at least, was straightforward; yours is sly. For ask Aristo whether these things seem good to him: freedom from pain, riches, health; he will say no. What of the things contrary to these — are they evil? No more so. Ask Zeno; he will answer in just as many words. We, in amazement, may ask of each how on earth we can live our lives if we think it makes no difference to us whether we are well or ill, whether we are free of pain or racked by it, whether or not we can fend off cold and hunger. You will live, says Aristo, splendidly and magnificently, doing whatever has seemed good, never distressed, never in want of anything, never afraid.

\ciceroSection{70} What says Zeno? That these are monstrosities, and that on such a principle there is no way at all to live; that he, for his part, asserts there is an immense — I know not how vast — difference between the honourable and the base, but among other things no difference whatever. So far the same;

\ciceroSection{71} now hear the rest, and hold back your laughter, if you can. Those intermediate things, he says, among which there is no difference, are nonetheless of such a kind that some of them are to be chosen, others rejected, others wholly disregarded — that is, some you should want, others not want, others not care about. — But just now you had said there was nothing in those matters that made a difference. — And I say the same now, he will reply; but I mean that they make no difference with respect to the virtues and the vices.

\ciceroSection{72} Who, I ask you, did not know that? But let us hear him out. — Those things you mentioned, he says — to be healthy, to be wealthy, to be free of pain — I do not call goods, but I shall say in Greek that they are \textit{[Greek: proēgmena]}, and in Latin "advanced" — though I prefer "preferred" or "preeminent," as more bearable and gentler; while those others — sickness, want, pain — I do not call evils, but, if you like, "rejected." And so I do not say I seek the former out, but that I select them; nor that I wish for them, but that I take them up; while the contraries I do not say I flee, but that I set them aside, as it were. What says Aristotle and the rest of Plato's nurslings? That they call all the things in accordance with nature good, and those against it evil. Do you see, then, that your Zeno chimes with Aristo in words but parts from him in substance, while with Aristotle and his school he agrees in substance but differs in words? Why then, when there is agreement about the substance, do we not prefer to speak in the customary way? Or let him show that I shall be readier to despise money if I have counted it among preferred things rather than among goods, and braver in enduring pain if I have called it harsh and hard to bear and contrary to nature rather than if I have called it an evil.

\ciceroSection{73} Wittily our friend Marcus Piso used to mock the Stoics, on many points and on this one too. For what, he would say. You deny that riches are a good, you say they are a preferred thing? What do you gain by it? Do you lessen greed? In what way? If we are to follow the word, "preferred" is in the first place a longer word than "good." — That is nothing to the purpose! — Let it be so; but it is certainly weightier. For from what "good" got its name I do not know; "preferred," I believe, from this, that it is preferred to other things. That seems to me a great matter. And so he used to say that more was granted to riches by Zeno, who placed them among preferred things, than by Aristotle, who admitted that riches were a good — but not a great good, and one to be despised and looked down upon in comparison with right and honourable things, and not greatly to be sought after. And in general, about all those words changed by Zeno, he would argue in this way: that the things which Zeno denied to be goods and which he denied to be evils, the former Zeno called by more cheerful names than we do, the latter by gloomier ones. So Piso, then, in this manner — an excellent man and, as you know, most devoted to you. Let us, with a few additions to all this, at last make an end; for it would take long to answer everything that has been said by you.

\ciceroSection{74} For out of these same verbal sleights kingdoms have been born for you, and empires, and riches — and such riches indeed that you say all things, wherever they are, belong to the sage. He alone, besides, is handsome, he alone free, he alone a citizen, while the foolish are everything contrary, men whom you would even have to be insane. These things they call paradoxes \textit{[Greek: paradoxa]}; let us call them marvels. But what do they contain to marvel at, when you have come up close? I will go through with you what reality you assign to each word; there will be no dispute. You say all faults are equal. I will not jest with you now, as I did about these very matters when I was defending Lucius Murena while you were the prosecutor. Those things were said then before the unlearned, and something even granted to the gallery; now we must proceed more subtly. Faults are equal.

\ciceroSection{75} — In what way? — Because nothing is more honourable than the honourable, nor baser than the base. — Go on further; for about that there is great disagreement. Let us look at the proper arguments for why all faults are equal. — Just as, he says, with a set of strings, unless not one of them is so strung as to be able to keep the harmony, all are equally out of tune, so faults, because they are discordant, are equally discordant; therefore they are equal. — Here we are tricked by an ambiguity. For it happens equally to all the strings that they are out of tune; it does not follow at once that they are equally out of tune. So that comparison helps you not at all. For neither, if we have said that all instances of greed are equally greed, will it follow that we should also say they are equal. Here is another comparison, like yet unlike.

\ciceroSection{76} For just as, he says, the helmsman is equally at fault if he capsizes a ship laden with straw and one laden with gold, so likewise he is equally at fault who unjustly beats a parent and who beats a slave. — Not to see that of what kind the cargo is which the ship carries has nothing to do with the helmsman's art! And so whether he carries gold or straw makes no difference to steering well or badly! But what difference there is between a parent and a slave can be, and ought to be, understood. Therefore in steering it makes no difference, in duty it makes a great deal, what sort of thing is at fault. And if in the very act of steering the ship has been capsized through negligence, the fault is greater with the gold than with the straw. For we hold that to all the arts there belongs that which is called common prudence, which all who preside over any craft ought to have. So not even in this way are faults equal.

\ciceroSection{77} Still they press on and yield nothing. Since, they say, every fault springs from weakness and inconstancy, and since these vices are equally great in all the foolish, faults must necessarily be equal. As though it were granted either that vices are equally great in all the foolish, and that Lucius Tubulus was of the same weakness and inconstancy as Publius Scaevola, on whose motion that man was condemned — and as though there were no difference between the very matters in which the faults are committed, so that, the greater or smaller these are, the greater or smaller too are

\ciceroSection{78} the faults committed in them! And so — for now let the discourse be brought to a close — by this one fault above all your Stoics seem to me to be hard pressed: that they think they can hold two contrary positions. For what is so self-contradictory as that the same man should say that only what is honourable is good, and yet say that the desire for the things suited to living comes from nature? So when they want to keep the things that agree with the former position, they fall in with Aristo; when they flee that, they defend in substance the very things the Peripatetics do, while clinging to their words with their teeth. And again, while they are unwilling to have these words torn from their ranks, they come out the rougher, harsher, harder, both in speech and in manners.

\ciceroSection{79} Fleeing that gloominess and harshness of theirs, Panaetius approved neither the bitterness of their doctrines nor the thorns of their arguing, and was milder in the one respect, clearer in the other, and always had on his lips Plato, Aristotle, Xenocrates, Theophrastus, Dicaearchus, as his own writings show. These men I greatly advise you to study and handle with zeal and care. But since evening is coming on, and I must return to my villa, for now let this be enough;

\ciceroSection{80} but let us do this same thing often. Indeed let us, said he; for what better could we do? And this first service I shall exact of you, that you hear me refuting the things that have been said by you. But remember that you approve all that we hold, except that we use different words, whereas to me nothing of yours is approved. A pebble in my shoe, said I, as you leave; but we shall see. When these things had been said, we parted.

\ciceroBook{5}

\ciceroSection{1} After I had heard Antiochus, Brutus, as was my habit, in the company of M. Piso, in the gymnasium that is called the Ptolemaeum, and with us my brother Quintus and T. Pomponius and Lucius Cicero — my cousin by blood, but a brother in affection — we agreed among ourselves to spend the afternoon walking in the Academy, chiefly because at that hour the place would be free of every crowd. And so at the appointed time we all went to Piso. From there, with talk of this and that, we covered those six stades from the Dipylon. When we had come into the spaces of the Academy, which are not without reason famous, we found the solitude we had wanted.

\ciceroSection{2} Then Piso said: 'Should I call it something given us by nature, or some kind of illusion, that when we see the very places in which we have learned that men worthy of remembrance spent much of their time, we are moved more deeply than when we hear of the deeds of those same men, or read some piece of their writing? Just so I am moved now. For Plato comes into my mind, who we are told was the first to make a practice of holding discussions here; and those little gardens of his close by do not merely bring him to my memory but seem to set the man himself before my eyes. Here was Speusippus, here Xenocrates, here his pupil Polemo — whose very seat that was, the one we see. And for my own part, when I would look upon our own Senate-house — the Hostilian, I mean, not this new one, which seems to me smaller now that it has been made larger — I used to call to mind Scipio, Cato, Laelius, and above all my own grandfather; so great is the power of suggestion that lies in places; with good reason, then, has the art of memory been derived from them.'

\ciceroSection{3} Then Quintus said: 'It is exactly as you say, Piso. For just now, as I was coming here, that spot of Colonus turned my own thoughts toward itself, where its inhabitant Sophocles came before my eyes — Sophocles, whom you know how I admire and how I delight in him. Indeed, a kind of vision stirred me toward an older memory still, of Oedipus coming to this place and asking, in that most tender song, what these very spots might be — to no purpose, of course, yet it stirred me all the same.' Then Pomponius said: 'But as for me, whom you are all in the habit of attacking as a devotee of Epicurus, I spend much time, as you know, with Phaedrus — whom I love beyond all others — in the gardens of Epicurus, which we passed by just now; yet, warned by the old proverb, I bear the living in mind, and still I cannot forget Epicurus, even were I to wish it, since my friends keep his likeness not only on their tablets but even on their cups and on their rings.'

\ciceroSection{4} Here I said: 'Our friend Pomponius is plainly making a joke, and perhaps within his rights. For he has so settled himself at Athens that he is all but one of the Attics, and looks likely to earn that very surname. But I agree with you, Piso, that this is what happens in practice: that at the prompting of places we think rather more keenly and attentively about famous men. For you know that once I came with you to Metapontum, and would not turn aside to my host before I had seen that very place where Pythagoras ended his life, and his seat. But at this present time, though there are throughout every quarter of Athens, in the very places themselves, many tokens of men of the highest greatness, still it is that recess that moves me. For only a little while ago it belonged to Carneades, whom I seem to see — for his likeness is well known — and I think that the very seat, robbed of so great a greatness of genius, longs for that voice of his.'

\ciceroSection{5} Then Piso said: 'Since, then, each of us has something to say, what of our friend Lucius? Does he gladly visit the spot where Demosthenes and Aeschines used to fight it out against each other? For each man is drawn most by his own pursuit.' And he, blushing, said: 'Do not ask it of me — I, who have even gone down to the Phaleric shore, to the place where they say Demosthenes used to declaim against the surf, to train himself to master the roar with his voice. And only just now I turned a little to the right off the road, to approach the tomb of Pericles. Though indeed there is no end of that in this city; for wherever we set foot, we set it down on some piece of history.'

\ciceroSection{6} Then Piso said: 'And yet, Cicero, those pursuits of yours, if they look to the imitation of great men, are the mark of men of talent; but if only to the learning of the tokens of ancient memory, of the merely curious. We all urge you — already at a run, I hope — not only to wish to know the men you wish to know, but to wish to imitate them too.' Here I said: 'Although this young man, as you see, Piso, is doing what you prescribe, still your encouragement is welcome to me.' Then he, in the most friendly way, as was his habit, said: 'Let us all, indeed, bring everything we can to bear upon this young man's coming of age, and above all that he impart something of his own studies to philosophy as well — whether to imitate you, whom he loves, or to be able to do that very thing he pursues with more grace. But are you to be urged by us, Lucius, or are you already inclined of your own accord? To me, at least, you seem to attend to Antiochus, whom you hear, well enough.' Then he, shyly, or rather modestly, said: 'I do indeed; but did you hear, just now, about Carneades? I am carried off that way — yet Antiochus calls me back, and there is no one else for us to hear besides.'

\ciceroSection{7} Then Piso said: 'Although this, perhaps, will not be allowed to pass off so easily, since he is here' — he meant me — 'still I shall make bold to call you back from this New Academy to that old one, in which, as you used to hear Antiochus say, there are reckoned not only those who are called Academics — Speusippus, Xenocrates, Polemo, Crantor, and the rest — but also the old Peripatetics, whose chief was Aristotle, whom, excepting Plato, I should perhaps not be wrong to call the chief of philosophers. Turn yourself, then, to them, I beg you. For from their writings and their teachings one may draw not only all liberal learning, all history, all elegant discourse, but there is besides such a variety of arts in them that no one, without that equipment, can come properly furnished to any of the more distinguished tasks. From these men have arisen orators, from these generals and leaders of commonwealths. To come to lesser matters: mathematicians, poets, musicians, and physicians too have set out, in the end, from this workshop, as it were, of every craftsman.' And I said:

\ciceroSection{8} 'You know I feel the very same, Piso; but you have made mention of the matter at a good moment. For my friend Cicero here is eager to hear what the view is, regarding the ends of goods, of that old Academy you speak of, and of the Peripatetics. And we think you can explain it most easily, since you had Staseas of Naples with you for many years, and since for several months now we have seen you inquiring into these very things at Athens from Antiochus.' And he, laughing, said: 'Come, come' — for you wished, cleverly enough, to make me the starting-point of our discussion — 'let us set it out for the young man, if by chance we are able. For this solitude grants us what, were some god to promise it, I should never have thought possible: that I should hold a discussion in the Academy, like a philosopher. But let me not, while I oblige him, be a burden to you.' 'A burden to me,' I said, 'who asked you for this very thing?' Then, when Quintus and Pomponius had said that they wished the same, Piso began. And do attend, I beg you, Brutus, to his discourse, and judge whether it seems to have embraced sufficiently the view of Antiochus — a view I take to be the one you most approve, since you have often heard his brother Aristus.

\ciceroSection{9} He spoke, then, thus: 'How much ornament there is in the discipline of the Peripatetics has been said by me, as briefly as I could, a little while ago. But the form of that discipline, as of nearly all the rest, is threefold: one part is of nature, a second of reasoning, a third of living. Nature was investigated by them in such a way that no part of sky, sea, or land — to speak in the poet's manner — was passed over; indeed, when they had spoken of the first beginnings of things and of the whole world, so as to draw their conclusions not only by probable argument but even by the necessary reasoning of the mathematicians, they brought the greatest store of material, from things investigated for their own sake, to the knowledge of hidden things.

\ciceroSection{10} Aristotle pursued the origins, the ways of life, and the forms of all living creatures; Theophrastus, for his part, the natures of plants and the causes and accounts of nearly all the things brought forth from the earth; and from this knowledge the investigation of the most hidden things was made easier. By the same men were handed down precepts of reasoning, not by dialectic alone but by oratory too; and by Aristotle, as its chief author, the exercise of speaking on both sides of single questions was established — not so as always to argue against everything, in the manner of Arcesilas, and yet so as, in all matters, to bring out whatever could be said on either side.

\ciceroSection{11} And when the third part sought the precepts of living well, that too was applied by them not only to the conduct of private life but also to the governing of commonwealths. From Aristotle we have learned the customs, institutions, and disciplines of nearly all the states not of Greece only but of the barbarian world as well; from Theophrastus, even their laws. And whereas each of them had taught what kind of leader it befits a man to be in a commonwealth, and had written besides, at greater length, what the best condition of a commonwealth might be, Theophrastus did this further: he showed what tendencies there are in a commonwealth, and what shifts of circumstance, by which one must steer as the situation may require. But as for the manner of passing one's life, the one that most pleased them was indeed the quiet one, set in the contemplation and the knowledge of things; and because this was most like to the life of the gods, it seemed most worthy of the wise man. And on these matters their discourse is splendid and illustrious.

\ciceroSection{12} As to the highest good, however — since there are two kinds of their books, one written for the public, which they called \textit{[Greek: exoterikon]}, the other more finely worked, which they left in their notebooks — they do not always seem to say the same thing; and yet on the substance itself there is no inconsistency, at least among those I have named, nor any disagreement among them. But when the happy life is the question, and this is the one thing that philosophy ought to look to and to follow, whether it lies wholly within the power of the wise man or can be either shaken or torn away by adverse circumstances — on this they do at times seem to differ among themselves and to be in doubt. This is brought out most of all by Theophrastus' book on the happy life, in which a great deal is granted to fortune. And if this is so, wisdom could not guarantee a happy life. This account seems to me too delicate, so to speak, and too soft, for what the force and gravity of virtue demands. Let us hold, then, to Aristotle and to his son Nicomachus, whose carefully written books on ethics are said indeed to be Aristotle's own — though I do not see why the son could not have been like the father. As for Theophrastus, let us call upon him for most things, provided only that we hold to more firmness and strength in virtue than he held to. Let us be content, then, with these men.

\ciceroSection{13} For their successors are, in my view, better than the philosophers of the other schools, but they so degenerate that they seem to have been born of their own selves. First, Theophrastus' pupil Strato wished to be a natural philosopher; and in that, though he is great, still most of his work is novel, and very little of it concerns ethics. His pupil Lyco is rich in language, but rather thin in matters themselves. Lyco's pupil Aristo is then neat and elegant, but that gravity which is looked for in a great philosopher was not in him; his writings are indeed many and polished, but somehow his discourse has no authority.

\ciceroSection{14} I pass over many, among them a learned and agreeable man, Hieronymus — whom I no longer know why I should call a Peripatetic. For he set out freedom from pain as the highest good; and whoever disagrees about the highest good disagrees about the whole scheme of philosophy. Critolaus wished to imitate the ancients, and is indeed nearest to them in gravity, and his discourse runs full; and yet not even he stays within the institutions of his fathers. Diodorus, his pupil, adds freedom from pain to honour. He too is his own man, and, disagreeing about the highest good, cannot truly be called a Peripatetic. But the view of the ancients our friend Antiochus seems to me to follow most carefully, and he teaches that this same view was Aristotle's and Polemo's.

\ciceroSection{15} Our friend Lucius does wisely, then, in wishing above all to hear about the highest good; for once this is settled, everything in philosophy is settled. For in other matters, if anything has been passed over or left unknown, there is no more inconvenience than the worth of whichever of those things it is in which something has been neglected. But if the highest good is unknown, the plan of living must of necessity be unknown too; and from this there follows so great an error that men cannot know to what harbour they should make their way. But once the ends of things are known, when it is understood what is the limit both of goods and of evils, the road of life has been found, and the shaping of all our duties, since it is then asked to what each thing is to be referred;

\ciceroSection{16} from which the plan of living happily — the thing all men long for — can be found and procured. And since there is great disagreement as to wherein this lies, we must call in that division of Carneades which our friend Antiochus is glad to use. For he saw not only how many views of the highest good there had been among philosophers up to now, but how many altogether there could possibly be. He used to deny, then, that there is any art that proceeds from itself; for the thing grasped by an art is always external to it. There is no need to make this longer with examples. For it is plain that no art turns upon itself, but that the art is one thing and the object set before the art another. Since, therefore, just as medicine is the art of health and steering the art of navigation, so prudence is the art of living, it follows of necessity that this too has been established and set going from some thing outside itself.

\ciceroSection{17} On one point, however, nearly all are agreed: that the thing with which wisdom is concerned, and which it aims to attain, must be suited and adapted to nature, and of such a kind that it invites and draws on the mind's appetite of its own accord — what the Greeks call \textit{[Greek: hormē]}. But what that thing is which moves us in this way and is sought by nature at the very first dawning of life — on this there is no agreement, and it is here, when the highest good is investigated, that the whole disagreement among the philosophers lies. For the whole of that inquiry which concerns the ends of good and evil — when we ask what among these is the furthest and the ultimate — has its source to be found in the question of where the first promptings of nature lie; once that is discovered, the entire discussion of the highest good and evil is drawn from it as from a head. Some hold that the first appetite is for pleasure and the first recoil from pain. Others judge that freedom from pain is what is first taken up, and pain what is first turned away from.

\ciceroSection{18} From these others again set out — those who name what they call the primary things in accordance with nature, among which they count the soundness and preservation of all the parts, health, unimpaired senses, freedom from pain, strength, beauty, and the rest of the same kind, of which there are likenesses in our minds, as it were the first sparks and seeds of the virtues. Now since it is some one of these three by which nature is first moved, whether toward seeking or toward repelling, and since there can be nothing at all besides these three, it necessarily follows that the whole duty of avoiding or of pursuing is referred to one of them, so that the wisdom we have called the art of living is concerned with one of those three things, and from it draws the starting-point of the whole of life.

\ciceroSection{19} And from whatever it has settled upon as the thing by which nature is first moved, there will also arise an account of the right and the honourable, which can square with some one of those three, so that it is honourable either to do everything for the sake of pleasure, even if you do not attain it, or for the sake of not feeling pain, even if you cannot reach that, or for the sake of obtaining the things that are in accordance with nature, even if you achieve nothing. So it comes about that, as great as the difference is in the natural beginnings, so great is the unlikeness in the ends of good and evil. Others again, from the same beginnings, will refer every duty either to pleasure, or to not feeling pain, or to securing those primary things in accordance with nature.

\ciceroSection{20} Now that six views about the highest good have been set out, of the three nearest these are the champions: of pleasure, Aristippus; of not feeling pain, Hieronymus; of the enjoyment of those things we have called the primary ones in accordance with nature, Carneades — not, to be sure, as their author, but as their defender for the sake of argument. The earlier three were possible positions, of which one alone has been defended, and that vigorously. For that one should do everything for the sake of pleasure — granted that, even if we attain nothing, still that very purpose of so acting is in itself worth seeking, and honourable, and the sole good — no one has maintained. Nor has anyone thought the mere avoidance of pain to be in itself among the things worth seeking, unless one could actually escape it. But that one should do everything in order to obtain what is in accordance with nature, even if one does not attain it, and that this is both honourable and the sole thing worth seeking for its own sake and the sole good — this the Stoics do say.

\ciceroSection{21} These, then, are the six simple views about the ends of good and evil: two without a patron, four defended. But of joined and twofold expositions of the highest good there have been in all three, and indeed there could be no more, if you look deep into the nature of things. For either pleasure can be joined to honour, as Calliphon and Dinomachus held, or freedom from pain, as Diodorus held, or the primary things of nature, as the ancients held — the same men we have named Academics and Peripatetics. But since all things cannot be said at once, for the present this much will need to be noted: that pleasure is to be set aside, since we are born for certain greater things, as will soon appear. About freedom from pain much the same is generally said as about pleasure. Since, therefore, both about pleasure with Torquatus and about honour, in which alone all good was placed, with Cato, the argument has been carried on, in the first place what was said against pleasure falls in much the same way against freedom from pain.

\ciceroSection{22} Nor in fact need anything else be sought against that view of Carneades. For in whatever way the highest good is so set out that it is empty of honour, neither duties nor virtues nor friendships can stand on that reckoning. And the joining to honour of either pleasure or the absence of pain makes that very honour, which it wishes to embrace, base. For to refer the things you do to ends, one of which — if a man be free of evil — one would call his being in the highest good, while the other is concerned with the lowest part of nature, is to darken all the splendour of honour, not to say to defile it. There remain the Stoics, who, when they had taken over everything from the Peripatetics and the Academics, pursued the same matters under other names. These it is better to refute one by one. But for now, to what we are about;

\ciceroSection{23} of those men, when we please. But the freedom from care of Democritus, which is, as it were, a tranquillity of mind — what they call \textit{[Greek: euthumia]} — had to be set apart from this discussion, because that tranquillity of mind is itself the happy life; and we are inquiring not what it is, but whence it comes. Now the views of Pyrrho, Aristo, and Erillus, already hissed off and cast out, since they cannot fall within the circle we have drawn, were not to be brought in at all. For since this whole inquiry about the ends, and as it were the furthest bounds, of good and evil sets out from what we have said is suited and adapted to nature and is itself sought first for its own sake, this whole foundation is overturned both by those who, in matters in which there is nothing that is not either honourable or base, deny that there is any reason why one thing should be preferred to another, and think there is no difference at all among such things; and by Erillus, who, if he really held that nothing is good except knowledge, did away with every ground for taking counsel and with the very discovery of duty. With the views of the rest thus shut out, and since besides them there can be none, this view of the ancients must necessarily prevail. By the practice of the ancients, then, which the Stoics too employ, let us take our start from here.

\ciceroSection{24} Every living creature loves itself and, the moment it has come into being, sets about preserving itself, because this first appetite for guarding its whole life is given to it by nature — to preserve itself and to be so constituted as to be in the best condition in accordance with nature. At the outset it holds this constitution in a confused and uncertain way, so far only as to protect itself, whatever it may be, but it understands neither what it is nor what it can do nor what its own nature is. But when it has advanced a little and has begun to perceive how far each thing touches it and pertains to it, then it gradually begins to go forward and to recognize itself and to understand the reason why it has the appetite of mind we spoke of, and it begins both to seek the things it feels to be suited to its nature and to repel their opposites. Therefore for every living creature the object of its appetite lies in what is adapted to its nature. And so the end of goods turns out to be to live in accordance with nature, so constituted as to be in the best condition and most perfectly adapted to nature.

\ciceroSection{25} But since each living creature has its own nature, it necessarily follows that the end too is, for all of them, this: that nature be fulfilled — for nothing prevents some things from being common both among the rest of the animals and between man and the beasts, since the nature of all is shared — yet those furthest and highest things which we are seeking are marked off and distributed among the kinds of living creatures, each proper to its own and suited to what the nature of each requires.

\ciceroSection{26} And so when we say that for all living creatures the furthest end is to live in accordance with nature, this is not to be taken as if we were saying that there is one single end for all; rather, just as it can rightly be said to be common to all the arts that they are concerned with some body of knowledge, while the knowledge proper to each art is its own, so to live in accordance with nature is common to all living creatures, but their natures are diverse, so that one thing is natural for a horse, another for an ox, another for a man. And yet in all there is a common highest end — and that not only among living creatures, but also among all those things which nature nourishes, increases, and protects, in which we see that the things produced from the earth do many things in a way themselves of their own accord that avail toward living and growing, so that within their own kind they may arrive at their furthest end; so that it is now permitted to embrace all things in a single comprehension and to say without hesitation that every nature is the preserver of itself and has this set before it as its end and furthest goal: to keep itself in the best state of its own kind; so that the end of all things which thrive by nature must necessarily be similar, not the same. From which it ought to be understood that for man the ultimate good is to live in accordance with nature, which we interpret thus: to live by the nature of man, complete on every side and wanting nothing.

\ciceroSection{27} These things, then, must be unfolded by us — but if rather minutely, you will forgive it; for we owe a service to this young man's age, and to one perhaps now hearing these matters for the first time. 'Quite so,' I said, 'though even what you have said up to now was rightly said in that manner, however much it was suited to a young man.' 'The end of things worth seeking having now been defined,' he said, 'why this matter stands as I have said must next be demonstrated. For that reason let us begin from what I set down first, which is also in fact first, so that we may understand that every living creature loves itself. And although this admits of no doubt — for it is fixed in nature itself and grasped by each one's own senses, so that if anyone wished to say the contrary he would not be heard — still, that we may pass over nothing, I think the reasons too should be brought forward why this is so.

\ciceroSection{28} And yet how can it be understood or even conceived that there should be any living creature that hates itself? For contrary facts would clash. For when that appetite of mind has begun deliberately to draw to itself something that harms it, because it is hostile to itself, then, since it will do this for its own sake, it will both hate itself and at the same time love itself, which cannot be. And it necessarily follows that, if a man is an enemy to himself, he will think the things that are good to be evil, and on the contrary the evil to be good, and will flee what should be sought and seek what should be fled — which is, beyond doubt, the overturning of life. Nor, if some are found who seek out either nooses or other means of destruction, or like that man in Terence who 'resolved that for so long he would do his own son the less wrong' — as he himself says — 'until he himself should be made wretched' — is such a man on that account to be thought an enemy to himself.

\ciceroSection{29} But some are moved by grief, some by desire, and many too are carried away by anger, and when they knowingly rush into evils, they then think they are taking the best counsel for themselves. And so they say, and do not hesitate: 'Such is my way; do you do as you need to do.' And those who had declared war upon themselves would wish to be racked by day and tortured by night, and yet they would not accuse themselves on the ground that they had counselled badly in their own affairs. For that complaint belongs to those who are dear to themselves and love themselves. Therefore, whenever it is said that a man deserves ill of himself and is to himself an enemy and a foe, and in the end flees from life, let it be understood that there underlies some cause of this kind, so that from that very fact it can be understood that each man is dear to himself.

\ciceroSection{30} Nor indeed is it enough that there is no one who hates himself; this too must be understood, that there is no one who thinks it makes no difference to him how he himself fares. For the appetite of mind will be done away with if, just as in those matters between which there is no difference we are inclined to neither side, so in our own selves we shall judge that it makes no difference to us in what state we are. And further, if anyone should wish to say this too, it would be utterly absurd — that each man loves himself in such a way that this force of loving is referred to some other thing, and not to the very man who loves himself. When this is said of friendships, of duties, of virtues, however it is said, what is meant can still be understood; but in our own selves it cannot even be understood that we should love ourselves on account of some other thing — for example, on account of pleasure; for it is on our own account that we love pleasure, not on its account that we love ourselves.

\ciceroSection{31} And yet what is there more plainly evident than that each man is not only dear to himself, but vehemently dear? For who is there, or how few are there, in whom, when death draws near, the 'blood does not flee back in timid fright and the man go pale with fear'? It is, to be sure, a fault to shudder so violently at the dissolution of nature — which is likewise to be censured in the case of pain — but because nearly all are affected in this way, it is argument enough that nature shrinks from destruction; and the more certain men do this in a way that they are even justly censured for it, the more must it be understood that these very excesses would not have arisen in some, had there not been a certain measure set by nature. Nor do I mean the fear of death in those who flee death because they think they are being deprived of the good things of life, or because they dread certain terrors after death, or because they fear they may die in pain; for in little children, who think of none of these things, if ever in play we threaten them that we will throw them down from somewhere, they are seized with terror. Indeed even 'the wild beasts,' says Pacuvius, 'which lack the cunning of understanding for taking precaution,' when the terror of death is thrown upon them, 'shudder.' And who thinks otherwise of the wise man himself, but that, even when he has resolved that he must die, he is still moved by the parting from his own and by the very leaving of the light?

\ciceroSection{32} But most of all in this kind of case the force of nature is plainly visible, since many endure even beggary in order to live, and men worn out by old age are tormented by the approach of death and bear those sufferings which we see Philoctetes bearing in the plays. He, though racked by pains not to be borne, nevertheless prolonged his life by fowling, and 'slow, he transfixed the swift with the stroke of his arrows; standing, he transfixed those in flight,' as it stands in Accius, and from the weaving of their feathers made coverings for his body.

\ciceroSection{33} Am I speaking of the human kind, or of animals at large, when the nature of trees and plants is nearly the same? For whether, as the most learned men have held, some greater and more divine cause has engendered this force, or whether it comes about thus by chance, we see that the things the earth brings forth are kept sound by their bark and roots — which falls to animals through the distribution of the senses and a certain framing of the limbs. On this matter, indeed, though I side with those who hold that all these things are governed by nature, and that nature, were she to neglect them, could not herself exist, still I grant that those who differ on this point may think what they please; and, if they like, let them understand that whenever I say ``the nature of man'' I mean man — for it makes no difference here. For a person could sooner part from himself than lose the appetite for the things that serve his interest. With good reason, then, the weightiest philosophers sought the starting-point of the highest good in nature, and held that this appetite for the things suited to nature is bred into all, because they are held fast by that commendation of nature whereby they love themselves.

\ciceroSection{34} Next we must consider — since it is plain enough that each is dear to himself by nature — what the nature of man is. For that is what we are inquiring into. Now it is evident that man is composed of body and mind, the parts of the mind being first, those of the body second. Then we see this too: that the body is so shaped as to surpass others, and the mind so constituted as to be both furnished with senses and possessed of a commanding intelligence, to which the whole nature of man is subject, and in which there is a certain marvellous power of reason and of knowing and of science and of all the virtues. As for the things that belong to the body, they have no authority to be compared with the parts of the mind, and they are easier to come to know. And so let us begin with these.

\ciceroSection{35} How well suited to nature, then, are the parts of our body, and its whole figure and form and stature, is plain; nor is there any doubt that one understands which the forehead, the eyes, the ears, and the remaining parts are that are proper to man. But surely these need to be sound and vigorous and to have their natural motions and uses, so that none of them is wanting, nor sick nor crippled; for this is what nature requires. There is, moreover, also a certain bearing of the body that keeps to the motions and postures agreeing with nature; and if in these there is some fault — a twisting and distortion, or a deformed motion or posture — as if a man should walk on his hands, or backwards rather than forwards, he would plainly seem to be fleeing from himself, and, stripping the man out of the man, to hate nature. For this reason, too, certain ways of sitting, and bent and broken movements such as belong to wanton or effeminate men, are contrary to nature; so that, even if it come about through a fault of the mind, still the nature of man seems to be altered in the body.

\ciceroSection{36} And so, on the other hand, the moderate and even conditions, affections, and uses of the body seem to be suited to nature. Now the mind, too, ought not merely to be, but to be of a certain quality, so that it both has all its parts unimpaired and lacks none of the virtues. And in the senses each has its own virtue, namely that nothing should hinder each sense from discharging its own office in perceiving swiftly and readily those things that lie subject to the senses. But of the mind, and of that part of the mind which is sovereign and is named the intelligence, the virtues are many; yet there are two primary kinds: the one of those that are bred in by their own nature and are called involuntary, the other of those that, resting in the will, are wont to be called by a more proper name — and the surpassing excellence in the praise of minds belongs to these. To the former kind belong aptness in learning and memory, which are nearly all called by the one name of natural endowment; and those who have these virtues are called gifted. The other kind is that of the great and true virtues, which we call voluntary, such as prudence, temperance, fortitude, justice, and the rest of the same kind. And in sum these were the things to be said about body and mind, by which there is sketched out, as it were, what the nature of man requires.

\ciceroSection{37} From which it is plain that, since we are dear to ourselves and wish all things in mind and in body to be perfect, these very things are dear to us for their own sakes and hold within them the greatest weight toward living well. For one who has set before himself the preservation of himself must hold his own parts dear as well, and the dearer the more perfect they are and the more praiseworthy in their kind. For the life that is sought is one filled out with the virtues of mind and body, and in this the highest good must necessarily be placed, since it must be such as to be the limit of the things to be desired. Once this is grasped, it cannot be doubted that, since men themselves are dear to themselves in their own right and of their own accord, the parts too of both body and mind, and of those things that lie in the motion and posture of each, are cherished with their own dearness and are sought for their own sakes.

\ciceroSection{38} These things being set out, it is an easy inference that those things among our own are most to be sought which have the most worth — so that, of each best part which is sought for its own sake, the virtue is to be sought most of all. Thus it will come about that the virtue of the mind is set before the virtue of the body, and that the voluntary virtues of the mind overcome the involuntary ones — the voluntary, which are properly called virtues and far excel, because they are begotten of reason, than which there is nothing in man more divine. For of all the things that nature both creates and watches over, those that are either without mind or not far otherwise have their highest good in the body; so that the saying about the pig does not seem inept — that a mind was given to that beast in place of salt, to keep it from rotting. There are, moreover, certain beasts in which there is something resembling virtue, as in lions, as in dogs, in horses, in which we see certain movements not of bodies alone, as in pigs, but in some measure of minds as well. In man, however, the whole sum lies in the mind, and in the mind in reason, out of which springs virtue, which is defined as the perfecting of reason — a thing they hold must be expounded again and again.

\ciceroSection{39} Of the things, too, that the earth brings forth, there is a certain rearing and perfecting not unlike that of living creatures. And so we say that a vine both lives and dies, and a tree both young and old, and that it flourishes and ``grows old.'' From which it is not amiss to suppose that, as with living creatures, so with these there are certain things suited to their nature and others alien to it, and that there is a kind of tender of their growth and nourishment — namely the science and art of the farmer, which prunes, lops, raises up, lifts, and props, so that they can go where nature carries them; so that the vines themselves, could they speak, would confess that they ought to be so handled and tended. And now indeed — to speak of the vine above all — that which tends it is something external; for in the vine itself there is too little force present for it to thrive as well as it might, if no cultivation be applied.

\ciceroSection{40} But suppose sense were now to come to the vine, so that it had a certain appetite and moved of its own accord — what do you think it would do? Will it not, of itself, see to the things it formerly obtained through the vinedresser? But do you not see that a care will be added to it, of tending its senses too, and their whole appetite, and any limbs that have been joined to it? Thus to the things it always had it will join those that came to it afterwards, and it will not have the same end its cultivator had, but will wish to live according to that nature which has afterwards been added to it. So its end of good will be like what it was before, yet not the same; for it will no longer seek the good of a plant, but of an animal. What, then, if not sense only had been given to it, but a man's mind as well? Must not those former things both abide, so as to be tended, and these later additions be far dearer — and the better each part of the mind, the dearer — and the end of the highest good rest in that filling-out of nature, since mind and reason far and away excel? Thus that which is the last of all things to be sought, drawn from the first commendation of nature, climbs by many degrees until it arrives at the summit, which is built up out of soundness of body and out of the perfected reason of the intelligence.

\ciceroSection{41} Since, then, such is the form of nature as I have set it out, if, as I said at the beginning, each person knew himself the moment he came to be, and could judge what the force of his whole nature and of its single parts might be, he would at once see what this is that we are inquiring into — the highest and uttermost of all the things we seek — and could go wrong in nothing. But as it is, from the first nature is marvellously hidden, and can be neither seen through nor known. As our years advance, however, little by little, or rather slowly, we come to know ourselves, as it were. And so that first commendation, which nature has made of us to ourselves, is uncertain and obscure, and that first appetite of the mind does only this much: that we may be able to be safe and sound. But when we begin to discern, and to feel what we are and how we differ from the rest of living creatures, then we begin to pursue the things for which we were born.

\ciceroSection{42} The like of this we see in beasts, which at first do not stir from the place where they were born, then each moves by its own appetite. We see snakelings creep, ducklings swim, blackbirds fly, oxen use their horns, scorpions their stings — in short, that for each creature its own nature is its guide to living. And this likeness appears in the human kind as well. For little ones, at the first birth, lie there as if they had no mind at all. But when a little firmness has come to them, they use both mind and senses, and strive to raise themselves up, and use their hands, and recognize those by whom they are reared. Then they take delight in playmates and gladly gather with them, give themselves to play, are drawn by the hearing of stories, wish to be generous to others with what they have to spare, take notice of what goes on at home rather curiously, and begin to ponder something, and to learn, and wish not to be ignorant of the names of those they see; and in the things in which they vie with their playmates, if they have won, they lift themselves up with gladness, and, beaten, they are cast down and let their spirits fall. None of this, we must hold, comes about without a cause.

\ciceroSection{43} For the force of man is so begotten by nature that he seems made for grasping every virtue; and for this reason little ones are moved, without teaching, by likenesses of the virtues whose seeds they hold within them — for these are the first elements of nature, which, once enlarged, bring about virtue's sprouting, so to speak. For since we have been so born and made as to hold within us the beginnings both of doing something, and of loving certain people, and of generosity, and of returning a kindness, and as to have minds apt for science, prudence, fortitude, and estranged from their contraries, it is not without cause that we see in children those sparks of the virtues, as I called them, from which the philosopher's reason must be kindled, so that, following it as a kind of god for its guide, it may arrive at the uttermost of nature. For, as I have often said by now, in the feeble years and the weak mind the force of nature is discerned as if through a fog; but when, as it advances, the mind is strengthened, it does indeed recognize the force of nature, yet in such a way that it can advance further, while in itself it has only been begun.

\ciceroSection{44} We must enter, then, into the nature of things, and see through to the very depths what she requires; for otherwise we cannot know ourselves. And because this precept was too great to seem to come from a man, it was for that reason assigned to a god. The Pythian Apollo therefore bids us know ourselves. But this knowledge of ourselves is one thing: to know the force of body and of mind, and to follow that life which has the full enjoyment of these very things. Now since this appetite of the mind was from the beginning such that we should have those things I spoke of as perfect in nature as may be, it must be confessed that, once we have attained what was sought, nature comes to rest, as it were, in that as in her uttermost, and that this is the highest good — which surely, taken whole, must of its own accord and for its own sake be sought, since it was shown before that even its single parts are to be sought for their own sakes.

\ciceroSection{45} But in reckoning up the advantages of the body, if anyone shall think that pleasure has been passed over by us, let that question be deferred to another time. For whether pleasure is or is not among those things we have called first according to nature makes no difference to what we are about. For if, as it seems to me at least, pleasure does not fill out the goods of nature, it has rightly been passed over; but if it is in them, as some hold, nothing hinders this grasp of ours of the highest good. For to those things which have been established as first according to nature, if pleasure be added, some one advantage of the body will have been added, and it will not have changed that constitution of the highest good which has been set forth.

\ciceroSection{46} And so far, indeed, our reasoning has proceeded in such a way that the whole of it was drawn from the first commendation of nature. But now let us follow another kind of argument: that not only because we love ourselves, but because each part of nature, in both body and mind, has its own force, for that reason we are moved in these matters to the highest degree of our own accord. And, to begin with the body, do you not see how, if anything in the limbs is crooked or crippled or stunted, men hide it? How they even struggle and labour, if they can bring it about, that the body's fault either not appear or appear as little as possible? And how they endure many pains as well, for the sake of a cure, so that, even if the very use of the limbs is to be not greater but actually less, still their appearance may return to nature? For indeed, since all by nature think themselves to be sought whole, and that for no other reason but on their own account, the parts too must necessarily be sought for their own sakes, when the whole is sought for its own sake.

\ciceroSection{47} What? In the motion and posture of the body is there nothing that nature herself judges must be heeded? In what manner a man walks, sits; what the carriage of the face, what the look, may be in each? Is there nothing in these things that we account worthy or unworthy of a free man? Do we not think many men deserving of hatred who, by a certain motion or posture, seem to have scorned the law and measure of nature? And since these things are drawn from the body, what reason is there why beauty itself should not rightly be reckoned worth seeking for its own sake? For if we think the body's crookedness and stunting are to be fled for their own sakes, why should we not also — and perhaps the more — pursue the dignity of form for its own sake? And if we flee ugliness in the posture and motion of the body, what reason is there why we should not pursue beauty? Health, too, and strength, and freedom from pain we shall seek not only for their usefulness but also for their own sakes. For since nature wishes to be filled out in all her parts, she seeks for its own sake that condition of the body which is most in keeping with nature — a nature wholly thrown into disorder if the body is sick, or in pain, or wanting in strength.

\ciceroSection{48} Let us look at the parts of the mind, whose view is more brilliant; and the loftier they are, the clearer the signs of nature they give. So great is the love of knowing and of knowledge bred in us that no one can doubt that the nature of man is swept toward these things lured by no profit. Do we not see how children are not deterred even by beatings from gazing at things and searching them out? How, when driven off, they run back? How they rejoice to know something? How they are eager to tell it to others? How they are held by a procession, by the games and spectacles of that kind, and for the sake of it endure even hunger and thirst? And what of those who take delight in liberal studies and arts — do we not see that they take no account of health or of household estate, but bear all things, captured by knowing and by knowledge itself, and balance against the greatest cares and labours the pleasure they take from learning?

\ciceroSection{49} To me, at any rate, Homer seems to have glimpsed something of this kind in what he invented about the song of the Sirens. For it does not appear that they were in the habit of calling back those who sailed past by the sweetness of their voices, or by some novelty and variety of singing, but because they professed to know many things — so that men clung fast upon their rocks out of a desire to learn. This is how they call to Ulysses (for I have translated this very passage, as I have certain others of Homer's): \textit{O glory of the Argives, why not turn your prow, Ulysses, that with your ears you may take in our song? For no one ever yet has sailed past these blue waters on his course without first standing still, caught by the sweetness of our voices, and then, sated in his eager breast with our manifold music, gliding onward to the shores of his fathers the wiser. We hold fast the heavy struggle of the war, and the ruin that Greece, by the will of heaven, brought upon Troy, and all the traces of things that are spread across the wide earth.} Homer saw that the story could not carry conviction if so great a man were held entangled by mere little ditties; what they promise is knowledge, and it is no wonder that to a man hungry for wisdom this should be dearer than his homeland. And indeed, to wish to know everything, of whatever kind, is the mark of the merely curious; but to be drawn to a desire for knowledge by the contemplation of greater things is what we must reckon the mark of the greatest of men.

\ciceroSection{50} For what burning zeal do you suppose was in Archimedes, who, while he was drawing certain figures in the dust with too close attention, did not even feel that his country had been taken? How much of Aristoxenus' genius do we see consumed in music? With what zeal do we think Aristophanes spent his life upon letters? What of Pythagoras? What shall I say of Plato or of Democritus? Men whom we see, for the desire of learning, to have traversed the farthest lands. Those who do not see this have never loved anything great and worthy of being known. And on this point, those who say that the studies I have named are pursued for the sake of the pleasures of the mind do not understand that they are for this very reason to be sought for their own sakes — namely, because our minds delight in them with no advantage held out, and rejoice in the knowledge itself even where it is going to bring inconvenience.

\ciceroSection{51} But what is the use of demanding more on matters so plain? Let us rather put the question to ourselves: how the movements of the stars and the contemplation of the heavenly bodies, and the knowledge of all those things that are hidden away in the obscurity of nature, move us; and what delight there is in history, which we are accustomed to pursue to the very end, going back over what we have passed by and following up what we have begun. Nor indeed am I unaware that there is utility in history, and not pleasure only.

\ciceroSection{52} What, when we read invented stories, from which no utility can be drawn out, and read them with pleasure? What, when we wish the names of those who have done some deed to be known to us, and their parents, their homeland, and many things besides that are not in the least necessary? What of the fact that men of the lowest fortune, with no hope of conducting affairs — craftsmen, in short — take delight in history? And we can see that those most of all wish to hear and to read of deeds done who are far from any hope of doing them, worn out by old age. From which it must necessarily be understood that in the very things which are learned and known there are inducements present, by which we are moved to learning and to knowing.

\ciceroSection{53} The old philosophers, indeed, imagine in the islands of the blessed what the life of the wise will be like: men freed from all care, requiring no necessary furnishing or provision for life, who they suppose will do nothing else than spend the whole of their time in inquiry and in learning, in the knowledge of nature. We, however, see that this is not only the delight of the happy life but also a relief from miseries. And so many men, when they were in the power of enemies or of tyrants, many in captivity, many in exile, have lightened their pain by the pursuits of learning.

\ciceroSection{54} The chief man of this city, Demetrius of Phalerum, when he had been driven from his country by injustice, betook himself to king Ptolemy at Alexandria. And since he excelled in this very philosophy to which we are urging you, and had been a pupil of Theophrastus, he wrote many splendid things in that calamitous leisure — not for any use of his own, of which he had been stripped, but that cultivation of his mind was for him a kind of food of humanity. For my part, I often used to hear from Cn. Aufidius, a man of praetorian rank, learned, and blind, that he was moved by a longing for the light more than for any use of it. In a word, unless sleep brought rest to our bodies and a kind of medicine for our labour, we should reckon it given against nature, for it takes away our senses and removes all activity. And so, if either nature did not seek rest, or could obtain it in some other way, we should readily put up with it — we who even now are accustomed, for the sake of doing or learning something, to take on our vigils almost against nature.

\ciceroSection{55} There are, moreover, still clearer, even plainly evident and not in the least doubtful signs of nature, above all in man, of course, but in every animal: that the mind craves always to be doing something, and on no terms can endure everlasting rest. This is easy to discern in the earliest little years of children. For although I fear I may seem excessive on this kind of point, still all the old philosophers, and ours above all, go back to the cradle, because they think that in childhood the will of nature can most easily be known. We see, then, how even infants cannot keep still. And when they have advanced a little, they take delight in games even if these are toilsome, so that they cannot be deterred even by beatings, and that desire of doing something grows up together with their years. And so, not even if we supposed we should enjoy the most delightful dreams would we wish the sleep of Endymion to be given us; and if it did befall us, we should reckon it the equal of death.

\ciceroSection{56} Indeed, even the most idle of men, endowed with some singular and indescribable sluggishness, we still see to be forever moved in body and mind, and, when they are hampered by no necessary business, either to call for the dice-board, or to seek out some game, or to look for some conversation; and since they do not have the noble delights that come from learning, to chase after little gatherings and gossiping circles. Indeed, not even the beasts that we shut up for the sake of our delight, although they are fed more abundantly than if they were free, readily endure to be confined, and they long for the free and wandering movements that nature has assigned them.

\ciceroSection{57} And so the better any man is born and trained, the less he would wish to be alive at all if, stripped of affairs to conduct, he could feed on pleasures however well prepared. For men either prefer to conduct some business privately, or, if they are of a loftier mind, take up the commonwealth by winning honours and commands, or devote themselves wholly to the pursuits of learning. In which life they are so far from chasing after pleasures that they even endure cares, anxieties, sleepless nights, and enjoy the best part of a man — which in us must be reckoned divine — the keenness of genius and of mind, neither seeking pleasure nor fleeing toil. Nor indeed do they let up either their wonder at the things that have been discovered by the ancients or their search for new ones. And since they cannot be sated by this zeal, forgetful of all other things, they think nothing low, nothing base; and so great is the power in such pursuits that we see even those who have set themselves other ends of goods, ones they steer by utility or by pleasure, nevertheless wear out their lives in inquiring into things and unfolding the workings of nature.

\ciceroSection{58} Therefore this much, at least, is clear: that we are born for action. Now there are several kinds of action, such that the lesser are even put in the shade by the greater; and the greatest are, first — as it seems to me, at any rate, and to those whose system we are now treating — the consideration and knowledge of the heavenly bodies and of those things which, hidden and lying concealed by nature, reason can track down; next, the administration of public affairs, or the science of administering them; then the prudent, temperate, brave, just exercise of reason and the remaining virtues, and the actions that accord with the virtues — all of which, embracing them in a single word, we call the honourable. To both the knowledge and the practice of these we are led, once we have grown strong, with nature herself going before us. For the beginnings of all things are small, but they grow by employing their own advances — and not without reason: for in the first emergence there is a certain tenderness and softness, such that one can neither see the best things nor do them. The light of virtue and of the happy life, the two things most of all to be sought, appears later, and much later still before it is understood plainly what these are. For Plato says it splendidly: Happy is the man to whom it has fallen, even in old age, to be able to attain wisdom and true opinions! Therefore, since enough has been said about the first advantages of nature, let us now look at the greater things that follow.

\ciceroSection{59} Nature, then, both begot and shaped the body of man in such a way that some things she perfected at the first emergence, others she fashioned as the years advanced, and she made hardly any use at all of external and adventitious aids. The mind, however, she perfected, as she did the body, by the remaining things; for she furnished it with senses fit for perceiving things, so that they needed no aid, or not much, for their own confirmation. But what is most excellent and best in man, that she left undone. And yet she gave such a mind as could receive every virtue, and implanted, without teaching, small notions of the greatest things, and, as it were, set out to teach and led us into the things that were already within — the elements, so to speak, of virtue. But virtue itself she only began; nothing more.

\ciceroSection{60} And so it is our task — when I say ours, I mean it is the task of an art — to seek out, upon those beginnings that we have received, the consequences, until that which we want is brought to completion. And this is worth a good deal more, and far more to be sought for its own sake, than either the senses or those things of the body which we mentioned; for the surpassing perfection of the mind so far excels them that it can scarcely be conceived what the difference is. And so all honour, all admiration, all zeal is referred to virtue and to those actions that are in accord with virtue, and all the things that are either so present in our minds or so done are called by one name the honourable. What the notions of all these are — those, that is, that are signified by the names of things — and what the force and nature of each is, we shall see presently.

\ciceroSection{61} But for now let us only explain that this honourable of which I speak — apart from the fact that we love our very selves — is besides, by its own nature, to be sought for its own sake. Children give the proof, in whom, as in mirrors, nature is discerned. What great zeal there is in those who contend! What great contests in themselves! How they are carried away with joy when they have won! How ashamed the beaten are! How they refuse to be found at fault! How they long to be praised! What labours do they not endure, to be foremost among their fellows! What memory there is in them of those who have done well by them, what desire to return the favour! And these things appear most of all in every best natural disposition, in which this honourable that we have in mind is, as it were, sketched out by nature.

\ciceroSection{62} But these things are in children; in those ages that are now full-formed they are clearly expressed. Who is so unlike a human being as not to be moved both by offence at baseness and by approval of honour? Who is there who does not hate a lustful, shameless youth? Who, on the other hand, does not love modesty and steadiness in that age, even where his own interest is not concerned at all? Who does not hate Pullus Numitorius of Fregellae, the traitor, although he was of service to our commonwealth? Who does not praise above all Codrus, the saviour of his city, who does not praise the daughters of Erechtheus? To whom is the name of Tubulus not hateful? Who does not love Aristides, though he is dead? Or do we forget how greatly we are moved in the hearing and in the reading, when we learn of something done with piety, with friendship, with greatness of soul?

\ciceroSection{63} Why do I speak of ourselves, who were born, taken up, and trained for praise and for honour? What shouts of the crowd and of the unlearned are raised in the theatres when those lines are spoken: 'I am Orestes', and against him from the other: 'No, in truth, I tell you, I am Orestes!' And when, moreover, the way out is given by each to the king, confused and at a loss — when both, then, pray to be put to death together — how often is this acted out, and ever with anything but the greatest admiration? There is no one, then, who does not approve and praise this disposition of mind, by which not only is no advantage sought, but, against advantage, faith is even kept.

\ciceroSection{64} With examples of this kind not only invented stories but histories too are crammed, and ours especially. For we chose the best of men to receive the sacred objects of Ida; we sent guardians to kings; our generals vowed their own lives for the safety of their country; our consuls warned a king, a bitter enemy already drawing near our walls, to beware of poison; in our commonwealth was found a Lucretia, who atoned by a voluntary death for the violation forced upon her, and a man who killed his daughter that she might not be violated. All these things, and countless others besides — who is there who does not understand both that those who did them were drawn by the splendour of worth and were heedless of their own advantage, and that we, when we praise them, are led by nothing else than the honourable? Now that these matters have been set out briefly — for I have not pursued the abundance I might have, since there was no doubt in the case — by all this it is surely concluded that both all the virtues and that honourable which springs from them and clings within them are to be sought for their own sake.

\ciceroSection{65} But in all that honourable good of which we are speaking, nothing is so radiant, nothing of wider reach, than the bond between man and man, a kind of fellowship and sharing of advantages, and that very love of the human race. This love, born from the first begetting — in that children are cherished by those who begot them, and the whole household is bound together by marriage and by stock — creeps gradually outward: first to kinsfolk, then to relations by marriage, then to friends, after that to neighbours, then to fellow citizens and to those who are publicly our allies and friends, and at last to the embrace of the whole human kind. This disposition of the mind, rendering to each his own and guarding generously and fairly that fellowship of human union of which I speak, is called justice; and joined to it are dutifulness, kindness, generosity, good will, courtesy, and whatever else is of the same kind. And these things, while they are proper to justice, are at the same time shared with the rest of the virtues.

\ciceroSection{66} For since man's nature has been begotten in such a way that he has within him something inborn and, as it were, civic and communal, which the Greeks call \textit{[Greek: politikon]}, whatever any virtue does will not be at odds with that community and that love and human fellowship which I have set out; and, in turn, just as justice itself will pour itself into the other virtues, so it will seek them out. For justice cannot be kept except by a brave man, except by a wise one. Such, then, as is this whole concord and consent of the virtues of which I speak, such is that honourable good itself, since the honourable is either virtue itself or a deed done by virtue; and a life that accords with these and answers to the virtues can be judged upright and honourable and steadfast and in keeping with nature.

\ciceroSection{67} And yet this joining and intermingling of the virtues is nonetheless distinguished by the philosophers through a certain method of reasoning. For although they are so coupled and connected that all share in all, and none can be severed from another, still each has its own proper office: courage is discerned in toils and dangers, temperance in the foregoing of pleasures, prudence in the choice of goods and evils, justice in rendering to each his own. Since, then, there is in every virtue a certain care that looks, as it were, outward and reaches after others and embraces them, it comes about that friends, that brothers, that kinsmen, that relations by marriage, that fellow citizens, that everyone in the end — since we hold that there is one fellowship of mankind — are to be sought for their own sake. And yet none of these is of such a kind as to lie within the end and uttermost limit of goods.

\ciceroSection{68} So it comes about that two kinds of things to be sought for their own sake are found. One lies in those things in which that uttermost good is brought to completion — things which belong either to the mind or to the body. But those that are external — that is, those that are present neither in the mind nor in the body, such as friends, parents, children, kinsmen, the fatherland itself — these are indeed dear of their own accord, but they are not of the same class as the others. Nor indeed could anyone ever attain the highest good if all those things that lie outside us, though to be sought, were contained within the highest good.

\ciceroSection{69} "How, then," you will say, "will it be possible for it to be true that all things are referred to the highest good, if friendships, if ties of kinship, if the remaining external things are not contained within the highest good?" By this account, surely: that we guard those things which are external by the duties that arise from each particular kind of virtue. For the cultivation of a friend and of a parent, to one who discharges his duty, profits him in this very thing, that to discharge duty thus belongs among right actions, which are sprung from the virtues. These the wise indeed pursue under nature's guidance, as though they could see; but men who are not yet perfect, and yet endowed with outstanding natural gifts, are often roused by glory, which has the look and likeness of the honourable. But if they could see, through and through, that honourable good itself, perfect and complete on every side … the one thing most splendid of all and most worthy of praise, with what joy would they be filled, when they take such delight in the shadowy notion of it?

\ciceroSection{70} For what man, given over to pleasures, what man set ablaze by the fires of his desires in laying hold of the things he had most fiercely craved, do we suppose to be flooded with so great a gladness as either the elder Africanus, when Hannibal was beaten, or the younger, when Carthage was overthrown? Whom did the descent of the Tiber on that festal day affect with so great a joy as it affected L. Paulus, when he was carried up that same river bringing King Perses captive? Come now, Lucius mine, raise up in your mind the height and excellence of the virtues:

\ciceroSection{71} you will no longer doubt that men who possess them, living with a great and lofty spirit, are always happy — men who understand that all the motions of fortune and the changes of circumstance and the times will be light and feeble, if once they have entered the contest of virtue. For those things which we counted among the goods of the body do indeed fill out the happiest life, but in such a way that the happy life can exist without them. For so small and slight are those additions of goods that, just as the stars are not even seen in the rays of the sun, so these are not even discerned in the splendour of the virtues. And though it is truly said that those advantages of the body count for little toward living happily, still it is too violent to say that they count for nothing;

\ciceroSection{72} for those who argue so seem to me to have forgotten the very first principles of nature which they themselves laid down. Something, then, must be granted to these things — provided you understand how much is to be granted. For it belongs to a philosopher who seeks not so much the showy as the true neither to count as nothing those things which even those showy men confessed to be in accordance with nature, and to see that there is so great a force in virtue and, so to speak, so great an authority in the honourable, that the rest, while not nothing, are yet so small as to seem to be nothing. This is the discourse of one who does not scorn all things save virtue, yet magnifies virtue itself with the praises that are its own; and, in fine, this is the explanation of the highest good, complete and perfect on every side. From here the rest, attempting to snatch up little fragments, each wished his own opinion to seem the one he had contributed.

\ciceroSection{73} Often, by Aristotle, by Theophrastus, the knowledge of things has been wondrously praised for its own sake; caught by this one thing, Erillus maintained that knowledge is the highest good, and that no other thing is to be sought for its own sake. Much was said by the ancients about despising and looking down upon human affairs; this one thing Aristo held to: he denied that, apart from vices and virtues, there is anything to be shunned or to be sought. It was laid down by our school that, among the things which are in accordance with nature, there is the being free from pain; this Hieronymus called the highest good. But Callipho, indeed, and after him Diodorus — the one having grown enamoured of pleasure, the other of freedom from pain — neither could do without the honourable, which has been praised most of all by our school.

\ciceroSection{74} Indeed, even the very devotees of pleasure seek their side-roads, and have the virtues on their lips all day long, and say that pleasure is at first sought only for its own sake, but that afterward, by habit, a kind of second nature is formed, driven by which men do many things while seeking no pleasure at all. The Stoics remain. They, to be sure, have carried over from us not some one thing or another, but our whole philosophy to their own side; and just as the rest of the thieves change the marks on the things they have stolen, so they, in order to use our opinions as their own, changed the names, like the marks upon things. Thus this teaching alone is left worthy of those zealous for the liberal arts, worthy of the learned, worthy of illustrious men, worthy of leading men, worthy of kings. When he had said this, and had paused a little, "What is it?" he said;

\ciceroSection{75} "do I seem to you to have rehearsed enough, in the use of my own right, within your hearing?" And I said, "You, Piso — as often at other times, so today — seemed to me to know these things so well that, if we could have you more often, I should not think we need do much suppliant courting of the Greeks. And this I approved the more, because I remember that Staseas of Naples, that teacher of yours, a Peripatetic of high repute indeed, used to speak of these matters somewhat otherwise, agreeing with those who put much store in fortune favourable or adverse, much in the goods or evils of the body." "It is as you say," he said; "but these things are stated by Antiochus, our friend, much better and more strongly than they used to be stated by Staseas. And yet I am not asking what has been approved by you in my account, but by this Cicero of ours, whom I am eager to carry off from you as my pupil."

\ciceroSection{76} Then Lucius: "For my part, those things have been thoroughly approved by me, as I think they are by my brother too." Then Piso said to me: "What of it, then? Do you grant the young man your indulgence? Or would you rather have him learn those doctrines which, when he has thoroughly mastered them, leave him knowing nothing?" "I, indeed, allow it to him," I said. "But do you not remember that I am at liberty to approve those things you have said? For who can fail to approve the things that seem to him probable?" "But can anyone," he said, "approve what he does not hold to be perceived, grasped, known?" "There is no great dissension here, Piso," I said. "For there is no other reason why nothing seems to me capable of being perceived, except that the power of perceiving is so defined by the Stoics that they deny anything can be perceived unless it is true in such a way that it could not be false. And so this is my dissension with them, but with the Peripatetics none at all. But let us pass these matters by; for they make a discussion both quite long enough and quite contentious enough."

\ciceroSection{77} "This seems to me to have been said by you too hastily — that all the wise are always happy; somehow the discourse flew past it. And unless this is made good, I fear that what Theophrastus said about fortune, about pain, about the torment of the body — things with which he judged the happy life could in no way be joined — may be true. For this is vehemently at odds: that the same man should be happy and yet overwhelmed by many evils. How these things are consistent I do not at all understand." "Which is it, then," he said, "that does not please you — that there is so great a force in virtue that of itself it is sufficient for living happily? Or, if you approve that, do you deny that it can come about that those who possess virtue should be happy even when afflicted with certain evils?" "For my part, I want there to be in virtue the greatest force possible; but how great it is, another time; for now only this: whether it can be so great, if anything outside virtue be counted among goods."

\ciceroSection{78} "And yet," he said, "if you grant the Stoics that virtue alone, if it be present, makes life happy, you grant the same to the Peripatetics as well. For what they do not dare to call evils, but grant to be harsh and inconvenient and to be rejected and foreign to nature — these we call evils, but slight and well-nigh the least of all. And so, if he can be happy who is among harsh and rejectable things, so too can he who is among small evils." And I said, "Piso, if there is anyone who is wont to see sharply, in pleading cases, what the matter at issue is, that man is surely you. So attend, I beg you. For so far — by my own fault, perhaps — you do not perceive what it is I am asking." "Here I am," he said, "and I await what you will answer to the very thing I was asking."

\ciceroSection{79} "I will answer that I am not asking, at this time, what virtue can accomplish," I said, "but what is said consistently, what is at variance with itself." "In what way, then?" he said. "Because," I said, "when by Zeno this is delivered grandly, as though from an oracle: 'Virtue is of itself sufficient for living happily' — and 'Why so?' he says, the answer comes: 'Because, save what is honourable, there is no other good.' I am no longer asking whether this be true; what I say is that the things he says hang together splendidly among themselves.

\ciceroSection{80} Let Epicurus have said this same thing — that the wise man is always happy (which indeed he is given to bubbling out now and then) — Epicurus who says that the wise man, even when he is being worn down by the highest pains, will say, 'How sweet this is! how little I care!'; I would not fight with the man as to why he should hold so much good to lie in nature; this I would press: that he does not understand what he ought to say, when he has called pain the highest evil. The same is now my discourse against you. You say all the same things, good and evil, that those men say who never, as the saying goes, set eyes on a painted philosopher: health, strength, stature, beauty, the soundness of all the little nails, goods; deformity, disease, debility, evils.

\ciceroSection{81} Now those external things — go gently on those, by all means; but since these are goods of the body, you will at least count among goods the things that produce them: friends, children, kinsmen, riches, honours, resources. Notice that against this I say nothing; what I do say is this: if those things are evils into which a wise man can fall, then to be wise is not enough for living happily. ‘On the contrary,’ he said, ‘for living most happily it is too little, but for living happily it is enough.’ ‘I noticed,’ I said, ‘that you put it in that way a little earlier, and I know that our friend Antiochus is in the habit of speaking so; but what could be less to be approved than that someone should be happy and yet not happy enough? For whatever is added to what is already enough is too much; and no one is too happy; and so no one is happier than the happy man.’

\ciceroSection{82} ‘Then, in your view,’ he said, ‘take Q. Metellus, who saw three sons made consuls — one of them even both censor and a celebrant of a triumph, and a fourth a praetor — and who left them all living, and three daughters married, while he himself had been consul, censor, and augur too, and had triumphed: granting that he was a wise man, was he not happier than Regulus — granting that he too was a wise man — who in the power of the enemy was done to death by sleeplessness and starvation?’

\ciceroSection{83} ‘Why ask me that?’ I said. ‘Ask the Stoics.’ ‘What, then,’ he said, ‘do you suppose they will answer?’ ‘That Metellus was no whit happier than Regulus.’ ‘From there, then,’ he said, ‘is where we must begin.’ ‘Even so,’ I said, ‘we are straying from the point. For I am not asking what is true, but what each man is bound to say. Would that they did say one man was happier than another! Then you would see ruins indeed. For since the good is placed in virtue alone and in the honourable itself, and since neither virtue, on their view, nor the honourable can grow, and since this is the only good — a good which whoever possesses must of necessity be happy — then, when the one thing in which being happy is placed cannot be increased, how can anyone be happier than another? Do you see how all this hangs together? And, by Hercules — for I must confess what I think — wonderful is the interweaving of their doctrine. The last answers to the first, the middle to both, all to all. They see what follows, what conflicts. It is as in geometry: grant the first premises, and you must grant the whole. Grant that nothing is good but the honourable: it must be granted that the happy life is placed in virtue. Now look back the other way:

\ciceroSection{84} grant this, and that other must be granted. Your school does not do the same. “Three kinds of goods”: the discourse runs forward down the slope. It comes to the end; it sticks fast in the rut. For it wants to say that nothing can be lacking to the wise man for the happy life — an honourable claim, a Socratic one, Plato’s too. ‘I dare to say it,’ he said. ‘You cannot, unless you unweave all that. If poverty is an evil, no beggar can be happy, however wise he is. Yet Zeno dared to call such a man not merely happy but rich as well. To feel pain is an evil: the man driven onto the cross cannot be happy. Children are a good: bereavement is wretched. Country is a good: exile is wretched. Health is a good: sickness is wretched. Soundness of body is a good: maiming is wretched. Unimpaired sight is a good: blindness is wretched. If wisdom can ease each of these singly by its consolation, how will it bear them all together? Suppose one and the same man blind, maimed, gripped by the gravest sickness, exiled, bereaved, destitute, and racked on the rack: what do you call him, Zeno?’ ‘Happy,’ he says. ‘Most happy too?’ ‘Of course,’ he will say, ‘since I have taught that happiness has no degrees, no more than virtue has, in which happiness itself also resides.’

\ciceroSection{85} ‘To you this is incredible — that he is most happy. Well? Is your own claim credible? For if you summon me to the people’s verdict, you will never convince them that a man so afflicted is happy; if to the verdict of the prudent, on the one point they will perhaps be in doubt — whether there is so much in virtue that those endowed with it are happy even in the bull of Phalaris — but on the other they will not doubt that the Stoics speak consistently with themselves and you with contradiction.’ ‘So you approve,’ he said, ‘of that book of Theophrastus on the happy life?’ ‘Still, we are straying from the point — and, not to drag it out, in a word, Piso,’ I said, ‘if those things are evils, I do approve of it.’ ‘Do they not, then,’ he said, ‘seem to you evils?’

\ciceroSection{86} ‘You are asking,’ I said, ‘a question in which, whichever way I answer, you are bound to twist me this way and that.’ ‘How so, pray?’ he said. ‘Because, if they are evils, the man who is in them will not be happy; if they are not evils, the whole system of the Peripatetics collapses.’ And he, laughing: ‘I see,’ he said, ‘what you are at; you are afraid I shall carry off your pupil.’ ‘You may carry him off, by all means,’ I said, ‘if he will follow; for he will be with me, if he is with you.’ ‘Listen, then, Lucius,’ he said; ‘for it is with you that I must frame my discourse. The whole authority of philosophy, as Theophrastus says, lies in the securing of a happy life; for we are all set ablaze with the desire of living happily. On this your brother and I agree.

\ciceroSection{87} And so what must be looked into is whether the reasoning of the philosophers can give us this. It certainly holds out the promise. For if it did not, why did Plato travel through Egypt to learn numbers and the lore of the heavens from the priests of the barbarians? Why, afterward, did he go to Tarentum to Archytas? Why to the other Pythagoreans — Echecrates, Timaeus, Arion — to the Locrians, so that, having reproduced Socrates, he might add the discipline of the Pythagoreans and learn besides the very things Socrates rejected? Why did Pythagoras himself both range over Egypt and go to the Magi of the Persians? Why did he traverse so many regions of the barbarians on foot, cross so many seas? Why did Democritus do the very same? Who — whether truly or falsely we forbear to ask — is said to have deprived himself of his eyes; and certainly, that his mind might be drawn off as little as possible from its meditations, he neglected his inheritance, left his fields untilled — seeking what else but the happy life? Even if he placed it in the knowledge of things, still it was from that investigation of nature that he wished to gain his end: to be of good cheer. For that highest good he calls cheerfulness \textit{[Greek: euthymia]}, and often fearlessness \textit{[Greek: athambia]}, that is, a mind free of dread.

\ciceroSection{88} But these things, splendid though they are, are not yet brought to a fine finish. For of virtue, indeed, only a little was said by this man, and that not even with full precision. After his time these matters first began to be inquired into in this city by Socrates, then were carried down to this present place, and there was no longer any doubt that in virtue lay all the hope both of living well and of living happily. When Zeno had learned this from our people, he did, as the formula runs in legal pleadings, “the same thing in another manner.” This now you approve in him. So by changing the names of things he escapes the charge of inconsistency — and we cannot escape it! He says that the life of Metellus was no happier than that of Regulus, yet to be preferred — not more to be sought after, but more to be taken up, and, if there were a choice, the life of Metellus to be chosen and that of Regulus rejected. I call the life he says is to be preferred and more to be chosen the happier one, and I assign to that life not the least particle more than the Stoics do.

\ciceroSection{89} What is the difference, except that I call known things by known words, while they look for new names by which to say the same thing? So, just as in the Senate there is always someone to call for an interpreter, in the same way these men must be heard by us with an interpreter at hand. I call good whatever is in keeping with nature, and bad whatever is against it — and not I alone, but you too, Chrysippus, in the Forum, at home; only in the school you leave off. What, then? Do you think men ought to speak one way, philosophers another? What value each thing has, the learned and the unlearned put differently; but once it has been settled among the learned what value each thing has — if they were men, they would speak in the customary way — let them, while the things stay the same, coin words at their own discretion.

\ciceroSection{90} But I come to the charge of inconsistency, lest you say too often that I am straying — the inconsistency you place in words, while I supposed it placed in the thing. If only this much be firmly grasped — and here we have the Stoics as our best helpers — that the force of virtue is so great that, if all else be set on the other side of the scale, it does not so much as appear, then take all the things they certainly call advantages, and to be taken up, and to be chosen, and preferred — which they so define that they are to be valued at a fairly high price — these things, then, which I call by goods’ name where the Stoics call them by so many names, partly new and made up, like those “promoted” and “demoted” of theirs, partly names that mean the same (for what is the difference between “seek after” and “choose”? to me, indeed, what is chosen, and toward which selection is exercised, seems even the finer term) — but, when I have called all those things goods, it matters only how great I say they are, when I have called them things to be sought after, how strongly. But if I no more declare them to be sought after than you to be chosen, and I who call them goods set no higher value on them than you who call them “promoted,” then all those things must of necessity be darkened and disappear and run into the rays of virtue as into the rays of the sun.

\ciceroSection{91} ‘But,’ comes the objection, ‘a life in which there is something of evil cannot be happy.’ Then neither is a cornfield happy, for all its rich and thronging ears, if you spy a wild oat anywhere; nor is a trade profitable, if amid the greatest gains it has taken on some little loss. Or is this true everywhere, but otherwise in life? And will you not judge the whole by its greatest part? Or is it in doubt that virtue holds so great a part in human affairs that it buries all the rest? I shall dare, then, to call the other things that are in keeping with nature goods, and not to cheat them of their old name nor go searching out some newer one instead; but the magnitude of virtue I shall set, as it were, in the other pan of the scale.

\ciceroSection{92} That pan, believe me, will press down both earth and sea. For a whole thing is always named from that which holds the largest parts and is spread most widely. We say someone lives in good cheer; if, then, he has once been made rather gloomy, is his cheerful life lost? Yet this did not happen even to that M. Crassus, who, Lucilius says, laughed but once in his life, so as to be on that account any the less called \textit{never-laughing} \textit{[Greek: agelastos]}, as the same poet has it. They used to call Polycrates of Samos fortunate. Nothing had befallen him that he did not wish, except that he had cast into the sea a ring he delighted in. Was he, then, unfortunate from that one vexation, and fortunate again when that very ring was found in the belly of a fish? But if he was a fool — which surely he was, since he was a tyrant — he was never happy; and if he was wise, he was not wretched even then, when he was driven onto the cross by Oroetes, the satrap of Darius. ‘But he was afflicted with many evils.’ Who denies it? But those evils were buried under the greatness of his virtue.

\ciceroSection{93} Or will you not grant even this to the Peripatetics, that they may say the life of all good men — that is, of the wise, men adorned with every virtue — has in all its parts always more of good than of evil? Who says this? The Stoics, of course. By no means; but those very men who measure all things by pleasure and pain — do they not cry out that the wise man always has more present that he would wish than that he would not? Since, then, men who confess that they would not so much as lift a hand for virtue’s sake, unless it produced pleasure, set so much in virtue, what ought we to do, who say that even the very smallest excellence of mind outstrips all the goods of the body, so that these are not so much as left in sight? For who is there who would dare to say that this could befall the wise man — that, if it were possible, he would cast virtue away forever to be freed from all pain? Which of us would say — we who are not ashamed to call evils the things the Stoics call “hardships” — that it is better to do something shamefully with pleasure than honourably with pain?

\ciceroSection{94} To us that Dionysius of Heraclea seems to have deserted the Stoics disgracefully on account of a pain in his eyes — as if he had learned from Zeno not to feel pain when he was in pain! What he had heard, but had not really learned, was that that pain is not an evil, because it is not shameful, and that it is to be borne by a man. Had he been a Peripatetic, he would have stayed, I believe, in his conviction — the Peripatetics, who, while they say that pain is an evil, give as their precepts about bearing its harshness bravely the very same as the Stoics give. And indeed your own Arcesilas, though he was rather stubborn in disputation, was nonetheless one of ours; for he was Polemo’s pupil. When he was burning with the pains of gout, and Charmides the Epicurean, a close friend, had visited him and was leaving downcast, ‘Stay, I beg you, our Charmides,’ he said; ‘nothing has come from there to here’ — and he pointed to his feet and his breast. And yet he would have preferred not to feel pain.

\ciceroSection{95} This, then, is our system, which seems to you inconsistent: that, on account of a certain heavenly and divine and so surpassing excellence of virtue — such that, where virtue is and great deeds done by virtue, deeds supremely worthy of praise, there misery and tribulation cannot be, though toil can be there, vexation can be — I should not hesitate to say that all wise men are always happy, yet that it can come about that one is happier than another. ‘And yet that very point, Piso, you must establish again and again,’ I said; ‘and if you hold it firm, you will be free to carry off not only my Cicero here but me myself.’

\ciceroSection{96} Then Quintus said: ‘To me, indeed, this seems established well enough, and I rejoice that this philosophy — whose furnishings I once valued more highly than the estates of all the rest, so rich did it seem to me that I could ask of it whatever I might covet in our studies — this philosophy, I say, I rejoice has been found keener too than the others, a thing some used to say it lacked.’ ‘Keener, at any rate, than ours,’ said Pomponius, with a jest; ‘but, by Hercules, your discourse was most welcome to me. For things I did not think could be said in Latin have been said by you in fitting words, and no less plainly than they are said by the Greeks. But it is time, if you please — and straight to my house, in fact.’ When he had said this, and it seemed enough had been argued, we all went on into the town, to Pomponius’s.

